% Traduction de quelques vers d'un poème ; POR y PARA ; debate guiado : « Ciudadanos del mundo » — Espagnol Tle A — Cours signé OnBuch+
% Programme officiel MINESEC. Compiler avec : tectonic E12-por-y-para-poema-ciudadanos-del-mundo.tex
\newcommand{\DOCMATIERE}{Espagnol}
\newcommand{\DOCNIVEAU}{Tle A}
\newcommand{\DOCMODULE}{Módulo 3 : Vie citoyenne et ouverture au monde}
\newcommand{\DOCLECON}{E12}
\newcommand{\DOCTITRE}{Traduction de quelques vers d'un poème ; POR y PARA ; debate guiado : « Ciudadanos del mundo »}
% OnBuch+ — préambule « cours signés » ENRICHI (compilé avec tectonic / XeLaTeX)
% Système visuel « Pop » d'OnBuch+ : fond crème (jamais blanc), bordures encre
% épaisses, ombres « dures » décalées, titres Archivo Black en capitales.
% Version MATHÉMATIQUES : graphes (pgfplots/tikz), boîtes pédagogiques
% (définition, propriété, méthode, attention, le savais-tu, exercice résolu,
% exercice, TP, bilan), page de garde et signature OnBuch+ ancrée en bas.
%
% Macros attendues dans main.tex AVANT \input{preamble} :
%   \DOCMATIERE  \DOCNIVEAU  \DOCMODULE  \DOCLECON (numéro)  \DOCTITRE
\documentclass[11pt]{article}

\usepackage{fontspec}
\usepackage[spanish,french]{babel} % français = langue principale ; l'espagnol via \esp{..} / espagnol
\usepackage[a4paper,margin=2cm,top=3.4cm,bottom=2.8cm,headheight=1.4cm,footskip=1.3cm]{geometry}
\usepackage{amsmath,amssymb,amsfonts}
\usepackage{mathtools} % \xmapsto, \coloneqq... souvent utilisés par les modèles
\usepackage{cancel} % simplifications barrées (bilans de forces)
% \abs{z} (module, valeur absolue) et \norm{u} (norme), utilisés par les modèles
\providecommand{\abs}{}\let\abs\relax\DeclarePairedDelimiter{\abs}{\lvert}{\rvert}
\providecommand{\norm}{}\let\norm\relax\DeclarePairedDelimiter{\norm}{\lVert}{\rVert}
\usepackage[table]{xcolor} % [table] : fournit aussi \rowcolors (alternance de lignes)
\usepackage{pagecolor}
\usepackage{tikz}
\usetikzlibrary{arrows.meta,positioning,calc,shapes.geometric,shapes.misc,shapes.arrows,decorations.pathmorphing,decorations.pathreplacing,decorations.markings,patterns,shadows,mindmap,trees,fit,backgrounds,matrix,intersections,angles,quotes}
\usepackage{pgfplots}
\usepackage[version=4]{mhchem} % équations nucléaires \ce{^{235}_{92}U}
\usepackage[european,siunitx]{circuitikz} % schémas électriques
\pgfplotsset{compat=1.18}
\usepgfplotslibrary{fillbetween,groupplots} % aires entre courbes (calcul intégral)
\usepackage{enumitem}
\usepackage{fancyhdr}
% [most] charge toutes les bibliothèques tcolorbox (listings, minted, raster,
% documentation...) et gonfle inutilement la mémoire TeX sur un document long
% (~300 boîtes) : on ne charge que ce qui est réellement utilisé.
\usepackage[skins,breakable]{tcolorbox}
\usepackage{graphicx}
\usepackage{colortbl}
\usepackage{booktabs}
\usepackage{tabularx}
\usepackage{diagbox} % case d'en-tête diagonale des tableaux à double entrée
% Colonnes extensibles centrée / gauche / droite (C, L, R), souvent utilisées par les modèles
\newcolumntype{C}{>{\centering\arraybackslash}X}
\newcolumntype{L}{>{\raggedright\arraybackslash}X}
\newcolumntype{R}{>{\raggedleft\arraybackslash}X}
\usepackage{array}
\usepackage{multicol}
\usepackage{multirow}
\usepackage{siunitx}
% parse-numbers=false : accepte aussi \SI{2,5 \times 10^{-3}}{...}
\sisetup{locale=FR,output-decimal-marker={,},group-minimum-digits=4,parse-numbers=false}
\DeclareSIUnit\ohm{\ensuremath{\symup{\Omega}}} % U+2126 absent de la police
\DeclareSIUnit\division{div} % réglages d'oscilloscope (ms/div, V/div)
\DeclareSIUnit\tr{tr} % tours (vitesse de rotation, ex. \SI{1500}{\tr\per\minute})
\DeclareSIUnit\molar{M} % concentration molaire (ex. \SI{3}{\molar}, \SI{1}{\micro\molar})
\DeclareSIUnit\an{an} % année (le modèle confond parfois avec \year, un compteur TeX) — ex. \SI{5}{\kilo\meter\per\an}
\DeclareSIUnit\FCFA{FCFA} % franc CFA (ex. \SI{500}{\FCFA\per\kilo\gram})
\DeclareSIUnit\degreeBrix{\textdegree Brix} % teneur en sucre d'un sirop/jus (réfractométrie)
\DeclareSIUnit\month{mois} % le modèle utilise parfois \month, un compteur TeX, comme unité de durée
\usepackage{qrcode}
% \degree : commande fréquemment utilisée par les modèles pour un angle ou une
% température en dehors de \SI{}{} (non fournie par les paquets chargés).
\providecommand{\degree}{^{\circ}}
% \qed (fin de démonstration), utilisé par les modèles sans amsthm
\providecommand{\qed}{\hfill$\square$}
% \barre : double barre des valeurs interdites dans un tableau de variations (tkz-tab)
\providecommand{\barre}{\ensuremath{\|}}
% environnement proof (amsthm non chargé) utilisé par les modèles
\ifdefined\proof\else\newenvironment{proof}[1][Démonstration]{\par\noindent\textit{#1.}\ }{\qed\par}\fi
\usepackage{lastpage}
\usepackage{hyperref}
\hypersetup{hidelinks}

% ============================================================
% Palette « Pop » OnBuch+ — mêmes hex que la classe P (plus_ui.dart)
% ============================================================
\definecolor{poporange}{HTML}{F59321}
\definecolor{popdark}{HTML}{E07A0C}
\definecolor{popcream}{HTML}{FFF6E8}
\definecolor{popink}{HTML}{1C1714}
\definecolor{poppurple}{HTML}{7A5AE0}
\definecolor{popgreen}{HTML}{2E9E6A}
\definecolor{poppink}{HTML}{E0457B}
\definecolor{popblue}{HTML}{2F7DE1}
\definecolor{popgold}{HTML}{C98A12}
\definecolor{popmuted}{HTML}{8A7F76}
\definecolor{popline}{HTML}{EADFD2}
\definecolor{popgray}{HTML}{8A7F76}
\colorlet{popgrey}{popgray}
% Alias tolérants (le modèle nomme parfois les couleurs différemment)
\colorlet{popwhite}{white}
\colorlet{popblack}{popink}
\colorlet{popred}{poppink}
\colorlet{popyellow}{popgold}
% Teintes claires (fonds de tableaux/encadrés) — le modèle nomme parfois
% les couleurs avec un suffixe L (ex. popblueL) sans les avoir définies.
\colorlet{poporangeL}{poporange!20}
\colorlet{popdarkL}{popdark!20}
\colorlet{poppurpleL}{poppurple!20}
\colorlet{popgreenL}{popgreen!20}
\colorlet{poppinkL}{poppink!20}
\colorlet{popblueL}{popblue!20}
\colorlet{popgoldL}{popgold!20}
\colorlet{popredL}{popred!20}
\colorlet{popgrayL}{popgray!20}
% Teintes claires pour fonds de tableaux / zones
\colorlet{popblueL}{popblue!10!white}
\colorlet{poporangeL}{poporange!14!white}
\colorlet{popgreenL}{popgreen!12!white}
\colorlet{poppurpleL}{poppurple!12!white}
\colorlet{poppinkL}{poppink!12!white}
\colorlet{popgoldL}{popgold!14!white}

\pagecolor{popcream}

% ============================================================
% Typographie
% ============================================================
\defaultfontfeatures{Ligatures=TeX}
\setmainfont{PlusJakartaSans-Medium.ttf}[Path=../../../fonts/,BoldFont=PlusJakartaSans-Bold.ttf]
% Symboles, API et emojis absents de la police principale : repli sur DejaVu Sans (caractères actifs)
\newfontface\symfont{DejaVuSans.ttf}[Path=../../../fonts/]
\makeatletter
\newcommand\symactive[1]{\catcode#1=13 \begingroup\lccode`\~=#1 \lowercase{\endgroup\def~}{{\symfont\char#1}}}
\newcommand\symblank[1]{\catcode#1=13 \begingroup\lccode`\~=#1 \lowercase{\endgroup\def~}{}}
\newcommand\symhyph[1]{\catcode#1=13 \begingroup\lccode`\~=#1 \lowercase{\endgroup\def~}{\mbox{-}}}
\newcount\symc
\newcommand\symrange[2]{\symc=#1 \@whilenum\symc<#2 \do{\expandafter\symactive\expandafter{\the\symc}\advance\symc by 1 }}
\newcommand\symblankrange[2]{\symc=#1 \@whilenum\symc<#2 \do{\expandafter\symblank\expandafter{\the\symc}\advance\symc by 1 }}
\makeatother
\symrange{"0250}{"02FF}\symactive{"2713}\symactive{"2714}\symactive{"2717}\symactive{"2718}\symactive{"2610}\symactive{"2611}\symactive{"2612}
\symhyph{"2011}\symhyph{"2E1A}\symblankrange{"1F300}{"1FAFF}\symblank{"FE0F}
% Maths en sans-serif (Fira Math) avec chiffres et lettres droites de Plus
% Jakarta, pour que les formules se fondent dans le texte.
\usepackage[math-style=ISO,bold-style=ISO]{unicode-math}
\setmathfont{FiraMath-Regular.otf}[Path=../../../fonts/]
\setmathfont{PlusJakartaSans-Medium.ttf}[Path=../../../fonts/,range={up/num,up/{Latin,latin},\mathup}]
\usepackage{icomma} % virgule décimale sans espace en mode math
% Caractères mathématiques Unicode absents des polices de texte (Plus Jakarta,
% Archivo Black) : rendus par leur équivalent mathématique, en texte comme en
% formule — sinon ils s'affichent en carré vide (ex. « Arithmétique dans ℤ »).
\usepackage{newunicodechar}
\newunicodechar{ℝ}{\ensuremath{\mathbb{R}}}
\newunicodechar{ℤ}{\ensuremath{\mathbb{Z}}}
\newunicodechar{ℂ}{\ensuremath{\mathbb{C}}}
\newunicodechar{ℕ}{\ensuremath{\mathbb{N}}}
\newunicodechar{ℚ}{\ensuremath{\mathbb{Q}}}
\newunicodechar{𝔻}{\ensuremath{\mathbb{D}}}
\newunicodechar{α}{\ensuremath{\alpha}}
\newunicodechar{β}{\ensuremath{\beta}}
\newunicodechar{γ}{\ensuremath{\gamma}}
\newunicodechar{δ}{\ensuremath{\delta}}
\newunicodechar{ε}{\ensuremath{\varepsilon}}
\newunicodechar{θ}{\ensuremath{\theta}}
\newunicodechar{λ}{\ensuremath{\lambda}}
\newunicodechar{μ}{\ensuremath{\mu}}
\newunicodechar{σ}{\ensuremath{\sigma}}
\newunicodechar{τ}{\ensuremath{\tau}}
\newunicodechar{φ}{\ensuremath{\varphi}}
\newunicodechar{ψ}{\ensuremath{\psi}}
\newunicodechar{ω}{\ensuremath{\omega}}
\newunicodechar{Σ}{\ensuremath{\Sigma}}
% majuscules produites par \MakeUppercase dans les titres (α-aminés -> Α)
\newunicodechar{Α}{\ensuremath{\alpha}}
\newunicodechar{Β}{\ensuremath{\beta}}
\newunicodechar{Γ}{\ensuremath{\Gamma}}
\newunicodechar{Δ}{\ensuremath{\Delta}}
\newunicodechar{Ε}{\ensuremath{\varepsilon}}
\newunicodechar{Θ}{\ensuremath{\Theta}}
\newunicodechar{Λ}{\ensuremath{\Lambda}}
\newunicodechar{Μ}{\ensuremath{\mu}}
\newunicodechar{Τ}{\ensuremath{\tau}}
\newunicodechar{Φ}{\ensuremath{\Phi}}
\newunicodechar{Ψ}{\ensuremath{\Psi}}
\newunicodechar{Ω}{\ensuremath{\Omega}}
\newunicodechar{↦}{\ensuremath{\mapsto}}
\newunicodechar{∈}{\ensuremath{\in}}
\newunicodechar{∉}{\ensuremath{\notin}}
\newunicodechar{∧}{\ensuremath{\wedge}}
\newunicodechar{⇔}{\ensuremath{\Leftrightarrow}}
\newunicodechar{⟺}{\ensuremath{\Longleftrightarrow}}
\newunicodechar{⇒}{\ensuremath{\Rightarrow}}
\newunicodechar{⟹}{\ensuremath{\Longrightarrow}}
\newunicodechar{∘}{\ensuremath{\circ}}
\newunicodechar{∪}{\ensuremath{\cup}}
\newunicodechar{∩}{\ensuremath{\cap}}
\newunicodechar{≡}{\ensuremath{\equiv}}
\newunicodechar{⊂}{\ensuremath{\subset}}
\newunicodechar{⊥}{\ensuremath{\perp}}
\newunicodechar{∃}{\ensuremath{\exists}}
\newunicodechar{∀}{\ensuremath{\forall}}
\newunicodechar{∛}{\ensuremath{\sqrt[3]{\,}}}
\newunicodechar{∜}{\ensuremath{\sqrt[4]{\,}}}
\newunicodechar{⃗}{\ensuremath{{}^{\to}}}
\newunicodechar{ⁿ}{\ifmmode^{n}\else\textsuperscript{\ensuremath{n}}\fi}
\newunicodechar{ˣ}{\ifmmode^{x}\else\textsuperscript{\ensuremath{x}}\fi}
\newunicodechar{ᵇ}{\ifmmode^{b}\else\textsuperscript{\ensuremath{b}}\fi}
\newunicodechar{ʳ}{\ifmmode^{r}\else\textsuperscript{\ensuremath{r}}\fi}
\newunicodechar{ᵘ}{\ifmmode^{u}\else\textsuperscript{\ensuremath{u}}\fi}
\newunicodechar{ʸ}{\ifmmode^{y}\else\textsuperscript{\ensuremath{y}}\fi}
\newunicodechar{ᵃ}{\ifmmode^{a}\else\textsuperscript{\ensuremath{a}}\fi}
\newunicodechar{ᵉ}{\ifmmode^{e}\else\textsuperscript{\ensuremath{e}}\fi}
\newunicodechar{ᵏ}{\ifmmode^{k}\else\textsuperscript{\ensuremath{k}}\fi}
\newunicodechar{ᵖ}{\ifmmode^{p}\else\textsuperscript{\ensuremath{p}}\fi}
\newunicodechar{ᴺ}{\ifmmode^{N}\else\textsuperscript{\ensuremath{N}}\fi}
\newunicodechar{ᶜ}{\ifmmode^{c}\else\textsuperscript{\ensuremath{c}}\fi}
\newunicodechar{ʰ}{\ifmmode^{h}\else\textsuperscript{\ensuremath{h}}\fi}
\newunicodechar{ᵠ}{\ifmmode^{q}\else\textsuperscript{\ensuremath{q}}\fi}
\newunicodechar{ₙ}{\ifmmode_{n}\else\textsubscript{\ensuremath{n}}\fi}
\newunicodechar{ₐ}{\ifmmode_{a}\else\textsubscript{\ensuremath{a}}\fi}
\newunicodechar{ᵢ}{\ifmmode_{i}\else\textsubscript{\ensuremath{i}}\fi}
\newunicodechar{ᵣ}{\ifmmode_{r}\else\textsubscript{\ensuremath{r}}\fi}
\newunicodechar{ₖ}{\ifmmode_{k}\else\textsubscript{\ensuremath{k}}\fi}
\newunicodechar{ᵦ}{\ifmmode_{\beta}\else\textsubscript{\ensuremath{\beta}}\fi}
\newunicodechar{ₓ}{\ifmmode_{x}\else\textsubscript{\ensuremath{x}}\fi}
\newfontfamily\popdisplay{ArchivoBlack-Regular.ttf}[Path=../../../fonts/]
\newfontfamily\poplogo{SpaceGrotesk-Bold.ttf}[Path=../../../fonts/]
\newfontfamily\popsemi{PlusJakartaSans-SemiBold.ttf}[Path=../../../fonts/]
\newfontfamily\popxbold{PlusJakartaSans-ExtraBold.ttf}[Path=../../../fonts/]
\newfontfamily\popxboldls{PlusJakartaSans-ExtraBold.ttf}[Path=../../../fonts/,LetterSpace=3.0]

\setlist[enumerate]{leftmargin=1.7em,itemsep=5pt,topsep=4pt}
\setlist[itemize]{leftmargin=1.7em,itemsep=4pt,topsep=4pt,label={\color{poporange}\textbullet}}
\setlength{\parindent}{0pt}
\setlength{\parskip}{5pt}
\renewcommand{\arraystretch}{1.25}

% pgfplots : style Pop par défaut
\pgfplotsset{
  every axis/.append style={axis line style={popink,line width=1pt},tick style={popink},
    grid style={popline,line width=0.5pt},label style={font=\small},tick label style={font=\footnotesize}},
  popcurve/.style={poporange,line width=1.8pt,no markers,smooth},
}
% Même style disponible dans un \draw TikZ simple (hors pgfplots)
\tikzset{popcurve/.style={poporange,line width=1.8pt}}
\tikzset{popgrid/.style={popline,very thin}}
\colorlet{popcurve}{poporange} % le nom du style est parfois utilisé comme couleur

% ============================================================
% Pill / squiggle
% ============================================================
\newcommand{\poppill}[3]{%
  \tikz[baseline=-0.55ex]\node[draw=popink,line width=1.2pt,rounded corners=2.6pt,%
    fill=#2,inner xsep=6.5pt,inner ysep=3pt]{%
    \popxboldls\fontsize{7}{7}\selectfont\color{#3}\MakeUppercase{#1}};%
}
\newcommand{\popsquiggle}[1]{%
  \tikz[baseline=-0.4ex]{
    \draw[#1,line width=2.6pt,line cap=round]
      plot [smooth] coordinates {(0,0.09) (0.34,0.01) (0.68,0.21) (1.02,0.01) (1.36,0.10)};
  }%
}
% Mot-clé surligné (marqueur orange sous le texte)
\newcommand{\cle}[1]{{\popxbold\color{popdark}#1}}

% ============================================================
% En-tête / pied
% ============================================================
\pagestyle{fancy}
\fancyhf{}
\renewcommand{\headrulewidth}{0pt}
\renewcommand{\footrulewidth}{0pt}
\lhead{\raisebox{-0.15ex}{\poplogo\fontsize{13}{13}\selectfont\color{popink}OnBuch\color{poporange}+}}
\rhead{\poppill{\DOCMATIERE\ \textbullet\ \DOCNIVEAU\ \textbullet\ Leçon \DOCLECON}{white}{popink}}
\lfoot{\popxbold\fontsize{7.6}{7.6}\selectfont\color{popmuted}\MakeUppercase{Cours signé OnBuch+}\ \textbullet\ {\popsemi\DOCTITRE}}
\rfoot{\tikz[baseline=-0.6ex]\node[rounded corners=8.5pt,fill=popink,minimum height=17pt,minimum width=17pt,inner xsep=5pt,inner sep=0]{\popxbold\fontsize{8}{8}\selectfont\color{popcream}\thepage\,/\,\pageref*{LastPage}};}

% ============================================================
% Titres
% ============================================================
\newcommand{\chaptitre}[2]{%
  \begin{tcolorbox}[enhanced,colback=poporange,colframe=popink,boxrule=2.6pt,arc=11pt,
      drop shadow=popink,left=15pt,right=15pt,top=12pt,bottom=11pt,before skip=3pt,after skip=18pt]
    \poppill{#2}{popink}{popcream}\par\vspace{9pt}
    {\raggedright\hyphenpenalty=10000\exhyphenpenalty=10000\popdisplay\fontsize{22}{24}\selectfont\color{popcream}\MakeUppercase{#1}\par}
  \end{tcolorbox}
}
% Section numérotée : pastille orange avec le numéro + titre Archivo
\newcounter{popsec}
\newcommand{\coursec}[1]{%
  \stepcounter{popsec}\setcounter{popsub}{0}%
  \par\vspace{16pt}\noindent
  \begin{minipage}{\linewidth}
  \tikz[baseline=-0.7ex]\node[draw=popink,line width=1.6pt,fill=poporange,rounded corners=4pt,minimum size=22pt,inner sep=0]{\popdisplay\fontsize{12}{12}\selectfont\color{popcream}\thepopsec};%
  \hspace{8pt}{\popdisplay\fontsize{15}{17}\selectfont\color{popink}\MakeUppercase{#1}}\par
  \vspace{4pt}\noindent{\color{poporange}\rule{36pt}{3.6pt}}
  \end{minipage}\par\nopagebreak\vspace{8pt}
}
\newcounter{popsub}
\newcommand{\courssub}[1]{%
  \stepcounter{popsub}%
  \par\vspace{9pt}\noindent{\popxbold\fontsize{11.5}{13}\selectfont\color{popink}{\color{poporange}\thepopsec.\thepopsub}\hspace{6pt}#1}\par\nopagebreak\vspace{4pt}
}

% ============================================================
% Boîtes pédagogiques — même recette : fond blanc, bordure épaisse colorée,
% ombre dure encre, badge plein. Titre optionnel [#1].
% ============================================================
\tcbset{popbase/.style={breakable,enhanced,colback=white,boxrule=2.3pt,arc=9pt,
  shadow={3pt}{-3pt}{0pt}{popink},left=14pt,right=14pt,top=11pt,bottom=11pt,before skip=11pt,after skip=15pt,
  fontupper=\popsemi\fontsize{10.6}{14.5}\selectfont\color{popink},
  fontlower=\popsemi\fontsize{10.6}{14.5}\selectfont\color{popink}}}
% Coche dessinée (le glyphe ✓ n'existe pas dans les polices du document)
\newcommand{\popcheck}[1]{\tikz[baseline=-0.5ex]\draw[#1,line width=1.6pt,line cap=round,line join=round](-0.13,0.0)--(-0.04,-0.09)--(0.13,0.1);}
\providecommand{\coche}{\popcheck{popgreen}}
\providecommand{\resultat}[1]{\par\smallskip\noindent\fcolorbox{popgreen}{popgreenL}{\parbox{\dimexpr\linewidth-2\fboxsep-2\fboxrule\relax}{\textbf{Résultat.} #1}}\par\smallskip}
\newcommand{\popboxhead}[3]{\poppill{#1}{#2}{white}\ifx\relax#3\relax\else\hspace{6pt}{\popxbold\fontsize{10.6}{12}\selectfont\color{popink}#3}\fi\par\vspace{7pt}}

\newtcolorbox{aretenir}[1][]{popbase,colframe=popgreen,before upper={\popboxhead{À retenir}{popgreen}{#1}}}
% Texte en espagnol (document de lecture, dialogue, poème, extrait) : cadre
% d'étude. Un titre optionnel (source, auteur) s'affiche dans l'en-tête.
\newtcolorbox{textebox}[1][]{popbase,colframe=popblue,before upper={\popboxhead{Texto}{popblue}{#1}\begin{otherlanguage}{spanish}},after upper={\end{otherlanguage}}}
% Marques ✓ / ✗ (forme correcte / fautive) souvent utilisées par les modèles
\providecommand{\cross}{\textcolor{poppink}{\ensuremath{\times}}}
\providecommand{\xmark}{\cross}\providecommand{\crossmark}{\cross}\providecommand{\cmark}{\ensuremath{\checkmark}}
% Ligne de réponse / justification à compléter (inventée par les modèles)
\providecommand{\justification}[1]{\par\noindent\textit{Justification~:} \dotfill\par}
% \textipa (package tipa non chargé) : la police ne couvre pas l'API -> simple passage
\providecommand{\textipa}[1]{#1}
% Soulignement (ulem non chargé) : \uline, \ul
\providecommand{\uline}[1]{\underline{#1}}\providecommand{\ul}[1]{\underline{#1}}
% \glossaire{\item ...} inventée par les modèles : liste de glossaire
\providecommand{\glossaire}[1]{\begin{itemize}#1\end{itemize}}
% \alert (beamer) inventée par les modèles : texte mis en évidence
\providecommand{\alert}[1]{\textbf{\textcolor{poppink}{#1}}}
% variantes ASCII d'ordinaux écrites en macro
\providecommand{\iere}{\textsuperscript{re}}\providecommand{\ieme}{\textsuperscript{e}}\providecommand{\ere}{\textsuperscript{re}}\providecommand{\eme}{\textsuperscript{e}}\providecommand{\er}{\textsuperscript{er}}\providecommand{\ie}{\textsuperscript{e}}
% Ordinaux français écrits en macro par les modèles : 1\ière, 3\ième
\newcommand{\ière}{\textsuperscript{re}}\newcommand{\ième}{\textsuperscript{e}}
% \exemple (inventée par les modèles) : étiquette « Exemple : »
\providecommand{\exemple}{\textbf{Exemple~:}\ }
% Mot ou phrase en espagnol à l'intérieur d'un paragraphe français
\newcommand{\esp}[1]{\ifmmode\text{\textit{#1}}\else\begingroup\selectlanguage{spanish}\itshape #1\endgroup\fi}
\newtcolorbox{exemplebox}[1][]{popbase,colframe=poporange,before upper={\popboxhead{Exemple}{poporange}{#1}}}
\newtcolorbox{definition}[1][]{popbase,colframe=popblue,before upper={\popboxhead{Définition}{popblue}{#1}}}
\newtcolorbox{propriete}[1][]{popbase,colframe=poppurple,before upper={\popboxhead{Propriété}{poppurple}{#1}}}
\newtcolorbox{methode}[1][]{popbase,colframe=popgold,before upper={\popboxhead{Méthode}{popgold}{#1}}}
\newtcolorbox{attention}[1][]{popbase,colframe=poppink,before upper={\popboxhead{Attention}{poppink}{#1}}}
\newtcolorbox{experience}[1][]{popbase,colframe=popblue,colback=popblueL,before upper={\popboxhead{Expérience}{popblue}{#1}}}
% « Le savais-tu ? » : carte violette pleine (contexte, culture, Cameroun)
\newtcolorbox{savaistu}[1][]{popbase,colback=poppurple,colframe=popink,
  fontupper=\popsemi\fontsize{10.4}{14.2}\selectfont\color{white},
  before upper={\poppill{Le savais-tu\,?}{white}{poppurple}\ifx\relax#1\relax\else\hspace{6pt}{\popxbold\color{white}#1}\fi\par\vspace{7pt}}}
% Exercice résolu : énoncé en haut, \tcblower puis la solution sur fond crème
\newcounter{popexo}\newcounter{popexr}
\newtcolorbox{exoresolu}[1][]{popbase,colframe=poporange,colbacklower=popcream,
  segmentation style={popink,line width=1.2pt,dashed},
  before upper={\stepcounter{popexr}\popboxhead{Exercice résolu \thepopexr}{poporange}{#1}},
  before lower={\poppill{Solution}{popink}{popcream}\par\vspace{6pt}}}
% Exercice (non résolu en place) — la correction est dans la partie Corrigés
\newtcolorbox{exercice}[1][]{popbase,colframe=popink,
  before upper={\stepcounter{popexo}\popboxhead{Exercice \thepopexo}{popink}{#1}}}
% Corrigé
\newtcolorbox{corrige}[1][]{popbase,colframe=popgreen,colback=popgreenL,
  before upper={\popboxhead{Corrigé}{popgreen}{#1}}}
% Objectifs / prérequis / bilan
\newtcolorbox{objectifs}{popbase,colframe=popink,colback=poporangeL,
  before upper={\popboxhead{Objectifs de la leçon}{popink}{}}}
\newtcolorbox{prerequis}{popbase,colframe=popink,colback=popblueL,
  before upper={\popboxhead{Prérequis}{popblue}{}}}
\newtcolorbox{bilan}{popbase,colframe=popink,colback=white,boxrule=2.8pt,
  before upper={\popboxhead{Fiche bilan}{poporange}{L'essentiel à connaître pour l'examen}}}
% Figure encadrée avec légende
\newtcolorbox{popfigure}[1][]{enhanced,breakable=false,colback=white,colframe=popline,boxrule=1.4pt,arc=8pt,
  left=10pt,right=10pt,top=10pt,bottom=8pt,before skip=10pt,after skip=12pt,halign=center,
  lower separated=false,halign lower=center,
  fontlower=\popsemi\fontsize{9}{11.5}\selectfont\color{popmuted},
  #1}
\newcommand{\legende}[1]{\par\vspace{6pt}{\popsemi\fontsize{9}{11.5}\selectfont\color{popmuted}#1}}

% ============================================================
% Signature OnBuch+
% ============================================================
\newcommand{\poplogoinline}[1]{{\poplogo\fontsize{#1}{#1}\selectfont\color{popink}OnBuch\color{poporange}+}}
\newcommand{\signatureonbuch}{%
  \begin{tcolorbox}[enhanced,colback=popink,colframe=popink,boxrule=2.2pt,arc=10pt,
      drop shadow=poporange,left=14pt,right=14pt,top=10pt,bottom=10pt,before skip=0pt,after skip=0pt]
    \tikz[baseline=-0.7ex]\node[circle,fill=popgreen,draw=popcream,line width=1.3pt,minimum size=22pt,inner sep=0]{\popcheck{white}};%
    \hspace{9pt}\begin{minipage}[c]{0.62\linewidth}
      {\popxbold\fontsize{12}{13}\selectfont\color{popcream}Signé\ \textbullet\ {\poplogo OnBuch\color{poporange}+}}\par\vspace{2pt}
      {\raggedright\popsemi\fontsize{8}{10.5}\selectfont\color{popline}\MakeUppercase{Équipe pédagogique OnBuch+}\\\MakeUppercase{Conforme au programme officiel MINESEC}\par}
    \end{minipage}\hfill
    {\poplogo\fontsize{22}{22}\selectfont\color{popcream}OnBuch\color{poporange}+}
  \end{tcolorbox}%
}

% ============================================================
% Page de garde
% #1 titre · #2 matière · #3 niveau · #4 module · #5 n° leçon · #6 durée ·
% #7 description / objectifs (texte ou itemize)
% ============================================================
\newcommand{\pagegarde}[7]{%
  \thispagestyle{empty}
  \newgeometry{a4paper,margin=1.7cm,top=1.6cm,bottom=1.4cm}%
  \begin{tikzpicture}[remember picture,overlay]
    \fill[poporange!18] ([xshift=-3.2cm,yshift=-3.6cm]current page.north east) circle (4.6cm);
    \fill[poppurple!14] ([xshift=2.4cm,yshift=7.6cm]current page.south west) circle (3.8cm);
  \end{tikzpicture}%
  {\poplogo\fontsize{30}{30}\selectfont\color{popink}OnBuch\color{poporange}+}\hfill
  \raisebox{0.6ex}{\poppill{Cours signé \textbullet\ Premium}{popink}{popcream}}\par
  \vspace{1pt}\popsquiggle{poppurple}\par
  \vspace{8pt}
  \begin{tcolorbox}[enhanced,colback=poporange,colframe=popink,boxrule=2.8pt,arc=13pt,
      drop shadow=popink,left=18pt,right=18pt,top=12pt,bottom=13pt,after skip=10pt]
    \poppill{#2 \textbullet\ #3}{popink}{popcream}\hspace{5pt}\poppill{#4}{white}{popink}\par\vspace{8pt}
    {\popdisplay\fontsize{14}{14}\selectfont\color{popink}LEÇON #5}\par\vspace{4pt}
    {\raggedright\hyphenpenalty=10000\exhyphenpenalty=10000\popdisplay\fontsize{25}{27}\selectfont\color{popcream}\MakeUppercase{#1}\par}
    \vspace{8pt}
    {\popxbold\fontsize{9.5}{9.5}\selectfont\color{popink}\MakeUppercase{Durée conseillée : #6}}
  \end{tcolorbox}
  \begin{tcolorbox}[enhanced,colback=white,colframe=popink,boxrule=2.2pt,arc=10pt,
      drop shadow=popink,left=16pt,right=16pt,top=10pt,bottom=10pt,after skip=8pt]
    \poppill{Dans cette leçon}{poppurple}{white}\par\vspace{5pt}
    \popsemi\fontsize{9.3}{12.6}\selectfont\color{popink}#7
  \end{tcolorbox}
  \noindent\begin{minipage}[t]{0.487\linewidth}
    \begin{tcolorbox}[enhanced,colback=popgreen,colframe=popink,boxrule=2.2pt,arc=10pt,
        drop shadow=popink,left=11pt,right=11pt,top=10pt,bottom=10pt,height=3.1cm,valign=center]
      \tcbox[colback=white,colframe=popink,boxrule=1.3pt,arc=5pt,boxsep=0pt,nobeforeafter,
        left=3pt,right=3pt,top=3pt,bottom=3pt,valign=center]{\qrcode[height=1.9cm]{https://play.google.com/store/apps/details?id=com.luvvix.onbuch&hl=fr}}%
      \hspace{8pt}\begin{minipage}[c]{0.5\linewidth}
        {\popdisplay\fontsize{11}{12.5}\selectfont\color{white}\MakeUppercase{Télécharge OnBuch}}\par\vspace{4pt}
        {\raggedright\popsemi\fontsize{8.4}{11}\selectfont\color{white}Gratuit \textbullet\ Google Play\\Tuteur IA, annales, concours, résultats\par}
      \end{minipage}
    \end{tcolorbox}
  \end{minipage}\hfill
  \begin{minipage}[t]{0.487\linewidth}
    \begin{tcolorbox}[enhanced,colback=poppurple,colframe=popink,boxrule=2.2pt,arc=10pt,
        drop shadow=popink,left=11pt,right=11pt,top=10pt,bottom=10pt,height=3.1cm,valign=center]
      \tikz[baseline=-0.7ex]\node[circle,fill=white,draw=popink,line width=1.3pt,minimum size=1.6cm,inner sep=0]{\popdisplay\fontsize{24}{24}\selectfont\color{poppurple}+};%
      \hspace{8pt}\begin{minipage}[c]{0.6\linewidth}
        {\popdisplay\fontsize{11}{12.5}\selectfont\color{white}\MakeUppercase{Passe à OnBuch+}}\par\vspace{4pt}
        {\raggedright\popsemi\fontsize{8.4}{11}\selectfont\color{white}Tous les cours signés, Tuteur IA illimité, fiches PDF — onglet OnBuch+ de l'app\par}
      \end{minipage}
    \end{tcolorbox}
  \end{minipage}
  \vspace{8pt}
  \popcontenu
  \vfill
  \signatureonbuch
  \restoregeometry
  \newpage
}

% Rangée « Ce cours contient » (page de garde)
\newcommand{\poptile}[3]{% #1 couleur #2 gros texte #3 légende
  \begin{tcolorbox}[enhanced,colback=#1,colframe=popink,boxrule=2pt,arc=9pt,shadow={2.5pt}{-2.5pt}{0pt}{popink},
      left=6pt,right=6pt,top=9pt,bottom=9pt,halign=center,valign=center,height=1.75cm,width=\linewidth,nobeforeafter]
    {\popdisplay\fontsize{12.5}{13}\selectfont\color{white}\MakeUppercase{#2}}\par\vspace{3pt}
    {\centering\popsemi\fontsize{7.8}{9.5}\selectfont\color{white}#3\par}
  \end{tcolorbox}}
\newcommand{\popcontenu}{%
  \par\noindent{\popxboldls\fontsize{7.5}{7.5}\selectfont\color{popmuted}CE COURS SIGNÉ CONTIENT}\par\vspace{7pt}
  \noindent\begin{minipage}[t]{0.235\linewidth}\poptile{popblue}{Cours}{complet, illustré, pas à pas}\end{minipage}\hfill
  \begin{minipage}[t]{0.235\linewidth}\poptile{poporange}{Méthodes}{et exercices résolus}\end{minipage}\hfill
  \begin{minipage}[t]{0.235\linewidth}\poptile{poppink}{Exercices}{type examen + corrigés}\end{minipage}\hfill
  \begin{minipage}[t]{0.235\linewidth}\poptile{popgreen}{Fiche bilan}{l'essentiel à retenir}\end{minipage}\par}

% Sommaire
\newtcolorbox{sommaire}{enhanced,colback=white,colframe=popink,boxrule=2.2pt,arc=10pt,
  drop shadow=popink,left=18pt,right=18pt,top=13pt,bottom=13pt,before skip=6pt,after skip=20pt,
  before upper={\poppill{Sommaire}{poppurple}{white}\par\vspace{9pt}\popsemi\fontsize{10.6}{14.5}\selectfont\color{popink}}}

% Signature de fin de cours — poussée tout en bas de la dernière page
\newcommand{\findecours}{%
  \par\vfill\nopagebreak
  \begin{center}\popsquiggle{poporange}\end{center}\vspace{4pt}
  \signatureonbuch
}
\AtBeginDocument{\def\llbracket{\mathopen{[\mkern-3.5mu[}}\def\rrbracket{\mathclose{]\mkern-3.5mu]}}} % intervalles d'entiers (glyphe ⟦ absent de la police)
% Glyphes absents de Fira Math : redéfinis à partir de glyphes disponibles
\AtBeginDocument{%
  \def\setminus{\mathbin{\backslash}}\let\smallsetminus\setminus
  \def\vdots{\mathord{\vbox{\baselineskip3.2pt\lineskiplimit0pt\kern4pt\hbox{.}\hbox{.}\hbox{.}}}}%
  \def\ddots{\mathinner{\mkern1mu\raise6.5pt\hbox{.}\mkern2mu\raise3.6pt\hbox{.}\mkern2mu\raise0.7pt\hbox{.}\mkern1mu}}%
  \def\triangle{\mathord{\scalebox{0.9}{$\symup{\Delta}$}}}%
  \def\nabla{\mathord{\raisebox{\depth}{\scalebox{1}[-1]{$\symup{\Delta}$}}}}%
  \def\checkmark{\popcheck{popdark}}%
  \def\star{\mathbin{\ast}}\def\bigstar{\mathord{\ast}}%
  \def\diamond{\mathbin{\scalebox{0.6}{$\Diamond$}}}%
  \def\bigcup{\mathop{\mathchoice{\raisebox{-0.3em}{\scalebox{1.7}{$\cup$}}}{\scalebox{1.25}{$\cup$}}{\cup}{\cup}}\displaylimits}%
  \def\bigcap{\mathop{\mathchoice{\raisebox{-0.3em}{\scalebox{1.7}{$\cap$}}}{\scalebox{1.25}{$\cap$}}{\cap}{\cap}}\displaylimits}%
  \def\lesseqqgtr{\mathrel{\lessgtr}}\def\bigoplus{\mathop{\oplus}\displaylimits}%
}
\providecommand{\ssi}{si et seulement si} % abréviation fréquente
\AtBeginDocument{\providecommand{\wideparen}{\overparen}} % arc de cercle

%% FIGFIX-BEGIN v3 — correctifs globaux des figures (docs/onbuch-generation/tools/figfix)
\usepackage{adjustbox}\usepackage{etoolbox}
% toute figure TikZ/pgfplots plus large que la ligne est réduite à la largeur disponible
\BeforeBeginEnvironment{tikzpicture}{\begin{adjustbox}{max width=\linewidth}}
\AfterEndEnvironment{tikzpicture}{\end{adjustbox}}
% légendes pgfplots lisibles par-dessus les courbes
\pgfplotsset{every axis/.append style={legend style={fill=white,fill opacity=0.92,text opacity=1,draw=black!15,inner sep=3pt}}}
% étiquettes d'arêtes (syntaxe edge["..."]) avec fond blanc
\tikzset{every edge quotes/.append style={fill=white,inner sep=1.2pt}}
% bulles de cartes mentales : texte à la ligne dans la bulle, bulle un peu plus grande
\tikzset{every concept/.append style={text width=2.0cm,align=center,minimum size=2.3cm,inner sep=2pt}}
%% FIGFIX-END
\begin{document}

\pagegarde{Traduction de quelques vers d'un poème ; POR y PARA ; debate guiado : « Ciudadanos del mundo »}{Espagnol}{Tle A}{Módulo 3 : Vie citoyenne et ouverture au monde}{E12}{2 h}{Leçon mixte de Terminale A : traduction guidée de quelques vers d'un poème engagé, étude complète des prépositions POR et PARA (cause, finalité, temps, lieu, agent, échange, destinataire), puis débat guidé sur les droits et la citoyenneté mondiale.
\vspace{4pt}
\begin{multicols}{2}\small
\begin{itemize}
\item À la fin de cette leçon, je sais traduire quelques vers d'un poème espagnol en respectant le sens, le rythme et les images
\item je sais distinguer les emplois de POR (cause, agent, échange, durée, trajet) et de PARA (finalité, destinataire, destination, date limite)
\item je sais choisir correctement entre por et para dans une phrase à traduire du français vers l'espagnol
\item je sais éviter les pièges fréquents des francophones (pour + infinitif, pendant, par + agent)
\item je sais utiliser le lexique des droits, des devoirs et de la citoyenneté
\item je sais participer à un débat guidé en espagnol : donner mon avis, argumenter, réagir à un contradicteur
\end{itemize}\end{multicols}}

\begin{sommaire}
\begin{multicols}{2}
\begin{enumerate}
\item Texte support : vers d'un poème engagé
\item Compréhension et lexique thématique
\item Traduction de vers : méthode et pratique
\item Grammaire : POR et PARA
\item Exercices de thème et de version
\item Débat guidé : « Ciudadanos del mundo »
\item Activité d'intégration
\item Méthodes et exercices résolus
\item Exercices
\item Corrigés des exercices
\item Fiche bilan
\end{enumerate}
\end{multicols}
\end{sommaire}

\begin{objectifs}
\begin{itemize}
\item À la fin de cette leçon, je sais traduire quelques vers d'un poème espagnol en respectant le sens, le rythme et les images
\item je sais distinguer les emplois de POR (cause, agent, échange, durée, trajet) et de PARA (finalité, destinataire, destination, date limite)
\item je sais choisir correctement entre por et para dans une phrase à traduire du français vers l'espagnol
\item je sais éviter les pièges fréquents des francophones (pour + infinitif, pendant, par + agent)
\item je sais utiliser le lexique des droits, des devoirs et de la citoyenneté
\item je sais participer à un débat guidé en espagnol : donner mon avis, argumenter, réagir à un contradicteur
\item je sais relier un texte poétique à une problématique de vie citoyenne et d'ouverture au monde
\end{itemize}
\end{objectifs}

\begin{prerequis}
\begin{itemize}
\item Les prépositions simples espagnoles vues en Première (a, de, en, con, sin)
\item Le présent de l'indicatif et de l'impératif (utiles pour le débat)
\item Les connecteurs d'opinion de Première : creo que, pienso que, me parece que
\item La notion de thème et de version (traduction français-espagnol et espagnol-français)
\item Le vocabulaire de base de la société et de la politique (ley, derecho, voto)
\end{itemize}
\end{prerequis}

\newpage

\coursec{Texte support : vers d'un poème engagé}

\courssub{Situation d'accroche et problématique}

Imaginez la salle des congrès de Yaoundé, un soir de juillet. Le sommet international de la jeunesse bat son plein : des délégués venus de Douala, de Madrid, de Mexico, de Malabo et de Bogotá se succèdent à la tribune. Une jeune Camerounaise, Aïcha, monte sur l'estrade et lit, d'une voix ferme, quelques vers de Federico García Lorca. Le silence est total. À la fin de sa lecture, un journaliste mexicain lève la main et lui pose une question en apparence anodine : \esp{¿Luchas por la paz o para la paz ?} « Combats-tu \emph{par} la paix ou \emph{pour} la paix ? »

La nuance est subtile, et toute la salle la sent : \esp{por} exprime la cause, le motif qui nous pousse à agir ; \esp{para} exprime la finalité, le but vers lequel nous marchons. Traduire un poème, c'est donc déjà débattre du monde : chaque mot choisi engage une vision de la citoyenneté. Aujourd'hui, nous apprendrons à lire et à traduire des vers engagés, à maîtriser l'opposition \esp{por}/\esp{para}, puis à débattre de ce que signifie être \esp{ciudadanos del mundo}, « citoyens du monde ».

\textbf{Problématique :} comment la poésie engagée hispanophone fait-elle du poète un citoyen du monde, et comment la traduire sans trahir son message ?

\courssub{Présentation du texte}

Le texte support de cette leçon est un poème original, écrit pour cette leçon, dans la grande tradition de la \cle{poésie engagée} hispanophone : celle de Federico García Lorca, dénonçant l'injustice ; celle de Pablo Neruda, chantant les humbles ; celle de León Felipe, poète de l'exil et de la fraternité universelle. Ce courant, né au XXe siècle, affirme que le poète n'est pas un rêveur isolé : il est un \cle{citoyen} qui prend la parole pour ceux qui ne l'ont pas.

Notre poème s'intitule \esp{No tengo patria, tengo hermanos}. Son thème : la fraternité au-delà des frontières, c'est-à-dire exactement le cœur du module « Vie citoyenne et ouverture au monde ».

\begin{textebox}[No tengo patria, tengo hermanos (poème original)]
No tengo patria, tengo hermanos ;\\
mi frontera es el alba,\\
mi bandera, un pan compartido.\\
Caminé por el mundo\\
para sembrar palabras de paz.\\
\\
No tengo patria, tengo hermanos :\\
los busqué por las plazas de Madrid,\\
por los mercados de Yaundé,\\
por las calles de Malabo al amanecer.\\
Mi pasaporte es mi voz,\\
mi himno, el llanto de los niños\\
que piden escuela y no guerra.\\
\\
Caminé por el mundo\\
para sembrar palabras de paz :\\
una palabra por cada surco,\\
una mano tendida por cada muro,\\
un abrazo para cada exiliado.\\
\\
No tengo patria, tengo hermanos.\\
Mi frontera es el alba :\\
allí donde nace la luz,\\
allí donde el pan se parte en dos,\\
allí empieza mi país.\\
\\
Y si me preguntan de dónde soy,\\
responderé con los ojos abiertos :\\
soy de aquí, soy de allá,\\
soy del lugar donde un hombre\\
tiende la mano a otro hombre.\\
No tengo patria, tengo hermanos,\\
y mi bandera es un pan compartido.
\end{textebox}

\begin{definition}[Glossaire : les mots difficiles du poème]
\begin{itemize}
\item \esp{la patria} : la patrie, le pays natal (mot plus solennel que \esp{país}).
\item \esp{el alba} : l'aube, le lever du jour (féminin malgré le \esp{el} : on dit \esp{el alba} mais \esp{el alba serena}).
\item \esp{el surco} : le sillon, trace laissée par la charrue dans la terre ; image du travail et de ce qu'on sème.
\item \esp{el hermano / la hermana} : le frère / la sœur ; au sens poétique, tout être humain fraternel.
\item \esp{la frontera} : la frontière.
\item \esp{la bandera} : le drapeau.
\item \esp{el pan compartido} : le pain partagé (\esp{compartir} = partager).
\item \esp{sembrar} : semer ; \esp{palabras de paz} : des paroles de paix.
\item \esp{el llanto} : les pleurs, les larmes (registre soutenu ; le mot courant est \esp{las lágrimas}).
\item \esp{el muro} : le mur.
\item \esp{el exiliado} : l'exilé (personne chassée de son pays).
\item \esp{el himno} : l'hymne.
\item \esp{amanecer} : (nom masculin) l'aube, le lever du jour ; (verbe) se lever (du soleil).
\item \esp{tender la mano} : tendre la main (verbe irrégulier \esp{tender}, e $\rightarrow$ ie : \esp{tiendo}).
\end{itemize}
\end{definition}

\courssub{Lecture à voix haute : pauses, accent tonique, musicalité}

Avant d'analyser, il faut \emph{entendre} le poème. La lecture à voix haute (\esp{la lectura en voz alta}) obéit à trois règles simples.

\begin{propriete}[Lire un vers espagnol à voix haute]
\begin{enumerate}
\item \textbf{Respecter les pauses.} On marque une courte pause en fin de vers, une pause plus longue à la fin d'une strophe. On ne s'arrête jamais au milieu d'un groupe de sens : on lit \esp{mi frontera / es el alba} et non \esp{mi / frontera es / el alba}.
\item \textbf{Placer l'accent tonique.} En espagnol, chaque mot porte un accent d'intensité sur une syllabe précise : \esp{paTRIa}, \esp{herMAnos}, \esp{fronTEra}, \esp{comparTIdo}. C'est cet accent régulier qui crée la musicalité du vers.
\item \textbf{Ne pas « chanter » artificiellement.} La voix suit le sens : elle monte légèrement sur une énumération (\esp{por las plazas de Madrid, por los mercados de Yaundé...}) et retombe sur la chute du vers (\esp{allí empieza mi país}).
\end{enumerate}
\end{propriete}

\begin{attention}[Piège de prononciation pour les francophones]
Le \esp{h} espagnol est toujours muet : \esp{hermanos} se prononce « er-manos », \esp{himno} « im-no ». De même, le \esp{ll} de \esp{llanto} se prononce comme un « y » fort (ou « li » mouillé en Espagne), jamais comme un « l » simple. Enfin, le \esp{y} de \esp{y mi bandera} se prononce « i ».
\end{attention}

\courssub{Première compréhension globale : qui parle ? à qui ? pour quoi ?}

Après une première lecture silencieuse puis orale, posons les trois questions fondamentales de toute compréhension de texte.

\begin{itemize}
\item \textbf{Qui parle ?} Une voix poétique à la première personne (\esp{no tengo}, \esp{caminé}, \esp{responderé}) : un voyageur sans patrie fixe, qui se définit par ses liens humains et non par ses papiers d'identité. C'est la figure du poète-citoyen du monde.
\item \textbf{À qui s'adresse-t-il ?} À tous les hommes, et singulièrement à ceux qui l'interrogent (\esp{si me preguntan de dónde soy}) : le lecteur, le journaliste, le douanier, c'est-à-dire chacun d'entre nous.
\item \textbf{Pour quoi ?} Pour proposer une autre définition de l'appartenance : non pas la frontière héritée, mais la fraternité choisie ; non pas le drapeau, mais le pain partagé. Le poème est un acte citoyen : il invite à repenser la citoyenneté.
\end{itemize}

\begin{exemplebox}[Questions de compréhension (avec réponses attendues)]
\begin{enumerate}
\item \esp{¿Qué tiene el poeta en lugar de una patria ?} — Il a des frères : \esp{tiene hermanos}.
\item \esp{¿Cuál es su frontera ?} — Sa frontière est l'aube : \esp{su frontera es el alba}.
\item \esp{¿Para qué caminó por el mundo ?} — Pour semer des paroles de paix : \esp{para sembrar palabras de paz}.
\item \esp{¿Qué piden los niños del poema ?} — Ils demandent l'école et non la guerre : \esp{piden escuela y no guerra}.
\end{enumerate}
\end{exemplebox}

\courssub{Repérage des procédés poétiques}

Passons maintenant à la lecture analytique : comment le poème fabrique-t-il son émotion et son message ?

\begin{definition}[L'anaphore]
L'\cle{anaphore} est la répétition d'un même mot ou groupe de mots en tête de vers ou de phrases successifs. Dans notre poème : \esp{mi frontera...}, \esp{mi bandera...}, \esp{mi pasaporte...}, \esp{mi himno...} ; puis \esp{una palabra...}, \esp{una mano...}, \esp{un abrazo...} ; et surtout \esp{No tengo patria, tengo hermanos}, repris comme un refrain. L'anaphore crée le rythme, martèle l'idée et donne au poème une force de manifeste.
\end{definition}

\begin{propriete}[Les principaux procédés à repérer dans un poème]
\begin{itemize}
\item \textbf{La métaphore} : une image sans outil de comparaison. \esp{Mi bandera es un pan compartido} : le drapeau \emph{est} un pain partagé (la fraternité remplace le symbole national). \esp{Mi pasaporte es mi voz} : la voix tient lieu de passeport.
\item \textbf{La répétition} : reprise d'un mot ou d'un vers entier (\esp{Caminé por el mundo / para sembrar palabras de paz} revient deux fois) pour insister et structurer le poème.
\item \textbf{Le rythme} : régularité des accents et des pauses, effet de balancement qui imite la marche (\esp{caminé}) du poète à travers le monde.
\item \textbf{L'antithèse} : opposition de deux termes. \esp{Escuela y no guerra} ; \esp{no tengo patria / tengo hermanos} : le poème avance par contrastes.
\item \textbf{L'énumération} : accumulation qui donne l'ampleur du voyage : \esp{por las plazas de Madrid, por los mercados de Yaundé, por las calles de Malabo}.
\end{itemize}
\end{propriete}

\begin{center}
\begin{tabularx}{\linewidth}{|X|X|X|}
\hline
\rowcolor{popblueL} \textbf{Procédé} & \textbf{Exemple du poème} & \textbf{Effet produit} \\
\hline
Anaphore & \esp{mi frontera... mi bandera...} & rythme, insistance, force de manifeste \\
\hline
Métaphore & \esp{mi bandera, un pan compartido} & la fraternité remplace la nation \\
\hline
Répétition & \esp{No tengo patria, tengo hermanos} & refrain mémorable, thème central \\
\hline
Antithèse & \esp{escuela y no guerra} & contraste moral frappant \\
\hline
Énumération & \esp{Madrid, Yaundé, Malabo} & ouverture au monde, universalité \\
\hline
\end{tabularx}
\end{center}

\begin{popfigure}
\begin{tikzpicture}[mindmap, grow cyclic, every node/.style={concept, text=white, align=center},
root concept/.append style={concept color=popblue, font=\large\bfseries},
level 1/.append style={level distance=3.2cm, sibling angle=90},
level 2/.append style={level distance=2.2cm, sibling angle=50}]
\node[root concept]{el poema}
child[concept color=poporange]{ node {voz poética} child { node {yo lírico} } child { node {viajero sin patria} } }
child[concept color=popgreen]{ node {imágenes} child { node {pan compartido} } child { node {alba, surco} } }
child[concept color=poppurple]{ node {ritmo} child { node {anáfora} } child { node {repetición} } }
child[concept color=poppink]{ node {mensaje ciudadano} child { node {fraternidad} } child { node {paz, escuela} } };
\end{tikzpicture}
\legende{Carte mentale du poème : quatre entrées pour l'analyser — la voix poétique, les images, le rythme, le message citoyen.}
\end{popfigure}

\courssub{Méthode : lire un poème en trois temps}

\begin{methode}[Les trois lectures d'un poème]
\begin{enumerate}
\item \textbf{Lecture globale} (sens général) : lire le poème entier, sans dictionnaire, et répondre à trois questions : qui parle ? à qui ? de quoi ? Formuler le thème en une phrase : « un poète sans patrie se dit frère de tous les hommes ».
\item \textbf{Lecture analytique} (images, sonorités) : souligner les procédés (anaphores, métaphores, répétitions), relever le champ lexical dominant (ici : la famille, le pain, la lumière, le voyage), et demander à chaque fois : quel effet ? pourquoi ce choix ?
\item \textbf{Lecture personnelle} (réaction) : dire ce que le poème nous fait penser et ressentir, le relier à notre expérience et à l'actualité. C'est cette lecture qui prépare le débat : suis-je, moi aussi, \esp{ciudadano del mundo} ?
\end{enumerate}
\end{methode}

\courssub{Premiers pas vers la traduction : observer por et para dans le poème}

Le poème contient déjà la leçon de grammaire du jour. Observons deux vers :

\esp{Caminé por el mundo.} « J'ai marché \emph{à travers} le monde. » — \esp{por} indique ici le lieu parcouru, le déplacement dans l'espace.

\esp{Para sembrar palabras de paz.} « \emph{Pour} semer des paroles de paix. » — \esp{para} + infinitif indique la finalité, le but de la marche.

Et dans la question du journaliste mexicain : \esp{¿Luchas por la paz ?} signifierait « tu luttes à cause de la paix, pour la défendre (elle est ton motif) », tandis que \esp{¿luchas para la paz ?} insisterait sur le but à atteindre. Nous étudierons en détail cette opposition dans la section grammaire ; retenons dès maintenant le réflexe : \esp{por} regarde en arrière (la cause), \esp{para} regarde en avant (le but).

\begin{exoresolu}[Première traduction guidée de deux vers]
Énoncé : traduisez en français les vers \esp{No tengo patria, tengo hermanos ; mi frontera es el alba, mi bandera, un pan compartido.}
\tcblower
Solution détaillée, étape par étape :
\begin{enumerate}
\item \textbf{Comprendre} : le poète refuse la patrie au profit des frères ; sa frontière n'est pas une ligne mais l'aube ; son drapeau n'est pas un tissu mais un pain partagé.
\item \textbf{Traduire mot à mot puis reformuler} : \esp{no tengo patria} = « je n'ai pas de patrie » ; \esp{tengo hermanos} = « j'ai des frères » ; \esp{mi frontera es el alba} = « ma frontière est l'aube » ; \esp{mi bandera, un pan compartido} = « mon drapeau, un pain partagé ».
\item \textbf{Version finale} : « Je n'ai pas de patrie, j'ai des frères ; ma frontière est l'aube, mon drapeau, un pain partagé. »
\end{enumerate}
Remarque : on conserve les métaphores telles quelles (l'aube, le pain) : trahir l'image, c'est trahir le poème.
\end{exoresolu}

\begin{attention}[Erreurs fréquentes des francophones]
\begin{itemize}
\item Ne traduisez pas \esp{patria} par « patrie » dans tous les contextes sans réfléchir : ici c'est exact, mais dans \esp{patria potestad} il s'agit de l'autorité parentale.
\item \esp{el alba} est féminin : « l'aube », et non « l'albe » ; attention à l'article : \esp{el alba} mais \esp{la misma alba}.
\item Ne confondez pas \esp{hermano} (frère) et \esp{hermanito} (petit frère) ; et méfiez-vous du faux ami \esp{éxito} (succès) que l'on rencontre souvent dans les poèmes engagés : ce n'est pas « exit ».
\item Dans \esp{caminé por el mundo}, \esp{por} ne se traduit pas par « pour » mais par « à travers, dans » : « j'ai marché à travers le monde ».
\end{itemize}
\end{attention}

\courssub{Le poète, citoyen du monde : ouverture culturelle}

\begin{savaistu}[Federico García Lorca, poète de la liberté]
Federico García Lorca (1898-1936), poète et dramaturge andalou, auteur de \emph{Romancero gitano} et de \emph{La casa de Bernarda Alba}, fut assassiné au début de la guerre d'Espagne, en août 1936, pour ses idées de liberté et sa défense des humbles. Il est devenu un symbole mondial du poète engagé : son corps disparu et sa voix restée vivante incarnent l'idée que la poésie est un acte citoyen. C'est dans cette tradition — avec Neruda au Chili, León Felipe en exil, ou Gabriela Mistral, prix Nobel 1945 — que s'inscrit notre poème : le poète n'a pas de frontière, parce que la fraternité n'en a pas.
\end{savaistu}

\begin{experience}[Activité orale : la lecture engagée]
Par groupes de quatre : chaque élève choisit une strophe du poème, la prépare (pauses, accents, intonation), puis la lit devant le groupe. Les trois autres élèves désignent ensuite le mot le plus important de la strophe et justifient leur choix en une phrase : \esp{Para mí, la palabra clave es « hermanos » porque...}. L'activité se termine par une lecture chorale du refrain : \esp{No tengo patria, tengo hermanos}.
\end{experience}

Nous avons entendu le poème, compris son sens global et repéré ses procédés. La section suivante approfondira le lexique thématique de la citoyenneté, avant que nous passions à l'art difficile et passionnant de la traduction des vers.

\coursec{Compréhension et lexique thématique}

Avant d'aborder la traduction littéraire et l'étude de \esp{por} et \esp{para}, nous devons nous approprier le texte support de cette leçon : un \textbf{poème original} écrit pour ce manuel, intitulé \esp{Ciudadanos del mundo}. On ne traduit bien que ce que l'on comprend parfaitement. Cette section a donc deux objectifs : vérifier la \cle{compréhension} du texte (sens littéral, \cle{locuteur}, thème, \cle{images}) et construire le \textbf{lexique thématique} de la citoyenneté, des droits humains, de l'ouverture au monde et de l'engagement, qui servira dans tout le module 3, y compris au débat guidé final.

\begin{textebox}[Poème original : Ciudadanos del mundo]
Yo soñé que la frontera \textbf{fuera} un alba,\\
que la bandera \textbf{fuera} un pan compartido,\\
y que todos los hombres, sin patria ni muro,\\
\textbf{se llamaran} hermanos en la misma lengua.
\end{textebox}

\emph{Traduction littérale :} « J'ai rêvé que la frontière \textbf{fût} une aube, que le drapeau \textbf{fût} un pain partagé, et que tous les hommes, sans patrie ni mur, \textbf{s'appellent} frères dans la même langue. »

\begin{attention}[Grammaire : concordance des temps / subjonctif imparfait]
Après un verbe de volonté, d'émotion ou d'irréel au passé (\esp{soñé}, prétérit), la subordonnée introduite par \esp{que} exige le \textbf{subjonctif imparfait} (\esp{fuera, se llamaran}) et non l'indicatif imparfait (\esp{era, llamaban}). C'est la norme grammaticale exigeante pour le Baccalauréat. Le poète exprime un vœu, une hypothèse irréelle, d'où le subjonctif.
\end{attention}

\courssub{Questions de compréhension du poème}

\begin{methode}[Comprendre un poème avant de le traduire]
\begin{enumerate}
\item \textbf{Première lecture globale} : lire le texte en entier sans dictionnaire, pour saisir l'impression générale (ton, émotion, sujet).
\item \textbf{Repérage du locuteur} : qui parle ? Chercher les marques de la première personne (\esp{yo}, \esp{mi}, \esp{soñé}, \esp{quiero}) et les temps verbaux.
\item \textbf{Sens littéral} : reformuler chaque vers en prose espagnole simple, puis en français, mot à mot si nécessaire.
\item \textbf{Sens figuré} : identifier les \cle{images} (métaphores, symboles) et se demander ce qu'elles représentent.
\item \textbf{Thème central} : résumer le poème en une phrase : de quoi parle-t-il vraiment ?
\end{enumerate}
\end{methode}

Appliquons cette méthode aux vers étudiés :

\begin{enumerate}
\item \textbf{Locuteur} : le \esp{yo} poétique, c'est-à-dire le poète lui-même, qui se présente comme un rêveur (\esp{soñé} : j'ai rêvé, prétérit indéfini de \esp{soñar}). Ce n'est pas un récit réel mais la vision d'un monde idéal (\cle{monde idéal}).
\item \textbf{Sens littéral} : « J'ai rêvé que la frontière était une aube, que le drapeau était un pain partagé, et que tous les hommes, sans patrie ni mur, s'appelaient frères dans la même langue. »
\item \textbf{Thème central} : le rêve d'une humanité unie, sans frontières ni divisions, fondée sur la fraternité et le partage. C'est le thème de la \textbf{citoyenneté mondiale} (\cle{ciudadanía mundial}).
\end{enumerate}

\begin{exemplebox}[Interprétation des images du poème]
\textbf{Image 1 : \esp{la frontera fuera un alba}.} La frontière, symbole de séparation et de fermeture, devient une \esp{alba} (une aube, un lever de soleil), symbole de commencement, de lumière et d'espérance. Le poète transforme ce qui divise en ce qui illumine : la limite entre les pays devient la promesse d'un jour nouveau.

\textbf{Image 2 : \esp{la bandera fuera un pan compartido}.} Le drapeau, symbole d'appartenance nationale et parfois de conflit, devient un \esp{pan compartido} (un pain partagé), symbole de nourriture, de solidarité et de vie commune. Au lieu de se battre pour un drapeau, les hommes se réunissent autour d'un même pain.

\textbf{Image 3 : \esp{sin patria ni muro... hermanos en la misma lengua}.} Supprimées les patries et les murs, les hommes se reconnaissent comme frères : la vraie « langue commune » est la fraternité (\esp{hermanos}).
\end{exemplebox}

\begin{exoresolu}[Questions de compréhension]
\textbf{Énoncé.} Réponds en espagnol, avec des phrases complètes.
\begin{enumerate}
\item ¿Quién habla en el poema y qué hace ?
\item ¿Qué significa la palabra \esp{alba}? ¿Por qué el poeta compara la frontera con un alba?
\item ¿Qué representa el \esp{pan compartido}?
\item ¿Cuál es el mensaje principal del poema?
\end{enumerate}
\tcblower
\textbf{Solution détaillée.}
\begin{enumerate}
\item \esp{Habla el poeta (un « yo » lírico), que cuenta un sueño: soñó con un mundo sin fronteras.}
\item \esp{« Alba » significa « amanecer », el momento en que sale el sol. El poeta compara la frontera con un alba porque quiere transformar un símbolo de separación en un símbolo de esperanza y de nuevo comienzo.}
\item \esp{El « pan compartido » representa la solidaridad y la vida en comunidad: en lugar de luchar por una bandera, los hombres comparten el mismo pan.}
\item \esp{El mensaje principal es que todos los seres humanos somos hermanos y ciudadanos del mismo mundo, más allá de las patrias, las fronteras y los muros.}
\end{enumerate}
\end{exoresolu}

\courssub{Un texte pour réutiliser le lexique : la citoyenneté au quotidien}

Pour fixer le vocabulaire, lisons maintenant un court texte original qui reprend les mots-clés du poème dans un contexte réaliste, celui d'un lycée camerounais.

\begin{textebox}[Un debate en el liceo de Yaoundé]
El viernes pasado, en el Liceo Bilingüe de Yaoundé, los alumnos de Terminale organizaron un debate titulado « Ciudadanos del mundo ». La profesora de español abrió la sesión con una pregunta sencilla: ¿qué significa ser ciudadano?

Amina, la delegada de la clase, respondió primero: « Ser ciudadano es tener derechos, pero también deberes. En Camerún, los ciudadanos votan para elegir a sus representantes, respetan las leyes y pagan sus impuestos. » Su compañero Rodrigue añadió que la ciudadanía no se limita a la nacionalidad: « También somos ciudadanos del mundo. Los derechos humanos, como la libertad de expresión, la igualdad y la dignidad, pertenecen a todas las personas, \textbf{vivan donde vivan}. »

Después, la conversación se volvió más viva. Algunos alumnos hablaron de los migrantes que cruzan las fronteras para buscar una vida mejor. « Nadie debería ser tratado como un extraño », dijo Kevin. « El extranjero de hoy puede ser el hermano de mañana. » Otros recordaron que la paz no cae del cielo: hay que sembrarla cada día, con diálogo, con justicia y con compromiso.

Al final del debate, la clase entera adoptó una conclusión común: defender los derechos de los demás es defender los propios. Luchar por la libertad, por la igualdad y por la paz no es solo tarea de los gobiernos: es el deber de cada ciudadano, aquí en Yaoundé, en Douala, en Malabo, en Madrid o en cualquier rincón del planeta. Como escribió el poeta, la verdadera bandera del ser humano es un pan compartido.
\end{textebox}

\begin{propriete}[Subjonctif dans les relatives à antécédent indéfini/négatif]
Remarquez \esp{vivan donde vivan} (présent subjonctif). La relative \esp{donde vivan} porte sur un antécédent indéfini (\esp{todas las personas}) : on ne sait pas *qui* précisément, ni *où* elles vivent. La règle : \textbf{antécédent indéfini, négatif ou superlatif $\rightarrow$ subjonctif dans la relative}. Cf. \esp{busco un libro que \textbf{trate} de la paz} (je cherche un livre qui traite de la paix — je ne l'ai pas encore trouvé).
\end{propriete}

\textbf{Glossaire français des mots difficiles :}

\begin{center}
\begin{tabularx}{\linewidth}{|X|X|}
\hline
\rowcolor{popblueL} \textbf{Español} & \textbf{Français} \\
\hline
\esp{el liceo} & le lycée (terme usité en Espagne et en Guinée équatoriale) \\
\hline
\esp{la delegada} & la déléguée (de classe) \\
\hline
\esp{pagar impuestos} & payer des impôts \\
\hline
\esp{vivan donde vivan} & où qu'ils vivent / où qu'elles vivent (subjonctif présent, antécédent indéfini) \\
\hline
\esp{tratar como} & traiter comme \\
\hline
\esp{el extraño} & l'étranger, l'inconnu (nom) \\
\hline
\esp{los propios (derechos)} & ses propres (droits) \\
\hline
\esp{tarea} & tâche, mission \\
\hline
\esp{cualquier rincón} & n'importe quel coin, le moindre recoin \\
\hline
\end{tabularx}
\end{center}

\textbf{Questions sur le texte :}
\begin{enumerate}
\item ¿Dónde tiene lugar el debate y quién lo organiza?
\item Según Amina, ¿qué significa ser ciudadano?
\item ¿Qué derechos humanos menciona Rodrigue?
\item ¿Cuál es la conclusión final de la clase?
\end{enumerate}

\courssub{Champ lexical 1 : la citoyenneté (\cle{la ciudadanía})}

\begin{definition}[La ciudadanía]
\esp{La ciudadanía} (la citoyenneté) désigne l'appartenance d'une personne à un État, avec les \esp{derechos} (droits) et les \esp{deberes} (devoirs) que cela implique. Le citoyen est \esp{el ciudadano} ; la nationalité est \esp{la nacionalidad} ; le vote est \esp{el voto} ; voter est \esp{votar}.
\end{definition}

\begin{center}
\begin{tabularx}{\linewidth}{|X|X|X|}
\hline
\rowcolor{popblueL} \textbf{Palabra} & \textbf{Traducción} & \textbf{Ejemplo en contexto} \\
\hline
\esp{la ciudadanía} & la citoyenneté & \esp{La ciudadanía se aprende en la escuela.} « La citoyenneté s'apprend à l'école. » \\
\hline
\esp{el / la ciudadano/a} & le / la citoyen(ne) & \esp{Los ciudadanos respetan las leyes.} « Les citoyens respectent les lois. » \\
\hline
\esp{el derecho} & le droit (juridique) & \esp{La educación es un derecho fundamental.} « L'éducation est un droit fondamental. » \\
\hline
\esp{el deber} & le devoir & \esp{Votar es un derecho y también un deber.} « Voter est un droit et aussi un devoir. » \\
\hline
\esp{el voto} & le vote, le suffrage & \esp{El voto es secreto, libre y obligatorio.} « Le vote est secret, libre et obligatoire. » \\
\hline
\esp{votar} & voter & \esp{Los jóvenes votan por primera vez.} « Les jeunes votent pour la première fois. » \\
\hline
\esp{la nacionalidad} & la nationalité & \esp{Tiene la nacionalidad camerunesa.} « Il a la nationalité camerounaise. » \\
\hline
\esp{la ley} & la loi & \esp{Todos somos iguales ante la ley.} « Nous sommes tous égaux devant la loi. » \\
\hline
\esp{elegir (a sus representantes)} & élire (ses représentants) & \esp{El pueblo elige a sus diputados.} « Le peuple élit ses députés. » \\
\hline
\end{tabularx}
\end{center}

\begin{attention}[Derecho : un mot, plusieurs sens — \cle{PIÈGE MAJEUR}]
\begin{itemize}
\item \esp{el derecho} (nm) = le droit (juridique) : \esp{los derechos humanos} = les droits de l'homme. \textbf{Jamais} « les humains de droite » !
\item \esp{derecho, -a} (adj) = droit (opposé à gauche) : \esp{la mano derecha} = la main droite.
\item \esp{la derecha} (nf) = la droite (direction ou tendance politique) : \esp{gira a la derecha} = tourne à droite.
\item \esp{derecho} (adv) = tout droit : \esp{sigue todo derecho} = continue tout droit.
\end{itemize}
\end{attention}

\courssub{Champ lexical 2 : les droits humains (\cle{los derechos humanos})}

\begin{center}
\begin{tabularx}{\linewidth}{|X|X|X|}
\hline
\rowcolor{popgreenL} \textbf{Palabra} & \textbf{Traducción} & \textbf{Ejemplo en contexto} \\
\hline
\esp{los derechos humanos} & les droits de l'homme & \esp{Los derechos humanos son universales, inalienables e indivisibles.} \\
\hline
\esp{la libertad} & la liberté & \esp{Luchar por la libertad es un deber de todos.} \\
\hline
\esp{la libertad de expresión} & la liberté d'expression & \esp{La libertad de expresión es pilar de la democracia.} \\
\hline
\esp{la igualdad} & l'égalité & \esp{Hombres y mujeres: igualdad de derechos y oportunidades.} \\
\hline
\esp{la justicia} & la justice & \esp{Sin justicia no hay paz duradera.} \\
\hline
\esp{la dignidad} & la dignité & \esp{Toda persona merece vivir con dignidad.} \\
\hline
\esp{el respeto} & le respect & \esp{El respeto a la diversidad enriquece a la sociedad.} \\
\hline
\end{tabularx}
\end{center}

\begin{exemplebox}[Exemples traduits à retenir par cœur]
\esp{Luchar \textbf{por} la libertad.} = « Lutter pour la liberté. » (Note le \esp{por} de la cause/motif : on lutte \emph{pour} une cause.)

\esp{Los derechos humanos pertenecen a todos.} = « Les droits de l'homme appartiennent à tous. »

\esp{Todos los seres humanos nacen libres e iguales en dignidad y derechos.} = « Tous les êtres humains naissent libres et égaux en dignité et en droits. » (Article 1 de la Déclaration universelle de 1948 : phrase culte à connaître pour le commentaire et la version.)
\end{exemplebox}

\begin{savaistu}[Le document le plus traduit du monde]
La \esp{Declaración Universal de los Derechos Humanos}, adoptée par l'ONU en 1948 après la Seconde Guerre mondiale, est le document le plus traduit de l'histoire : plus de 500 langues, dont plusieurs langues nationales du Cameroun (bassa, bamoun, douala, fulfulde...). Son article premier (\esp{Todos los seres humanos nacen libres e iguales en dignidad y derechos}) est un excellent entraînement à la version pour le Bac.
\end{savaistu}

\courssub{Champ lexical 3 : l'ouverture au monde (\cle{la apertura al mundo})}

\begin{center}
\begin{tabularx}{\linewidth}{|X|X|X|}
\hline
\rowcolor{poporangeL} \textbf{Palabra} & \textbf{Traducción} & \textbf{Ejemplo en contexto} \\
\hline
\esp{la frontera} & la frontière & \esp{Cruzar la frontera puede ser muy peligroso.} \\
\hline
\esp{el / la extranjero/a} & l'étranger / l'étrangère & \esp{Los extranjeros son bienvenidos en nuestra comunidad.} \\
\hline
\esp{el / la migrante} & le / la migrant(e) & \esp{Los migrantes buscan una vida digna y segura.} \\
\hline
\esp{la solidaridad} & la solidarité & \esp{La solidaridad une a los pueblos más que las fronteras.} \\
\hline
\esp{el diálogo} & le dialogue & \esp{El diálogo es el único camino para resolver los conflictos.} \\
\hline
\esp{el muro} & le mur & \esp{El poeta sueña con un mundo sin muros.} \\
\hline
\esp{la patria} & la patrie & \esp{La patria no debe separar a los seres humanos.} \\
\hline
\esp{el hermano / la hermana} & le frère / la sœur & \esp{Todos los hombres son hermanos en la humanidad.} \\
\hline
\end{tabularx}
\end{center}

\begin{exemplebox}[Point de conjugaison : \esp{cruzar}]
\esp{Cruzar la frontera.} = « Traverser la frontière. » Attention à l'orthographe au prétérit indéfini (z $\rightarrow$ c devant \emph{e}) : \esp{yo crucé, tú cruzaste, él cruzó, nosotros cruzamos, vosotros cruzasteis, ellos cruzaron}. Même règle pour \esp{empezar, comenzar, almorzar}.
\end{exemplebox}

\courssub{Champ lexical 4 : l'engagement et la paix (\cle{el compromiso y la paz})}

\begin{center}
\begin{tabularx}{\linewidth}{|X|X|X|}
\hline
\rowcolor{poppinkL} \textbf{Palabra} & \textbf{Traducción} & \textbf{Ejemplo en contexto} \\
\hline
\esp{luchar (por)} & lutter (pour) & \esp{Luchamos por la justicia social.} \\
\hline
\esp{sembrar} & semer (au sens propre et figuré) & \esp{Hay que sembrar la paz cada día con pequeños gestos.} \\
\hline
\esp{defender (los derechos)} & défendre (les droits) & \esp{Defender los derechos de los demás es defender los propios.} \\
\hline
\esp{la paz} & la paix & \esp{La paz se construye con justicia y diálogo.} \\
\hline
\esp{el compromiso} & l'engagement & \esp{El compromiso de la juventud es indispensable.} \\
\hline
\esp{compartir} & partager & \esp{Compartir el pan es el primer gesto de hermandad.} \\
\hline
\end{tabularx}
\end{center}

\begin{attention}[Faux amis et pièges de conjugaison — \cle{À MAÎTRISER ABSOLUMENT}]
\begin{itemize}
\item \esp{sembrar} = semer. \textbf{Faux ami} : ne pas confondre avec FR « sembler » (\esp{parecer}). \textbf{Diphtongaison} \emph{e} $\rightarrow$ \emph{ie} : \esp{siembro, siembras, siembra, sembramos, sembráis, siembran} (présent indicatif/subjonctif).
\item \esp{defender} = défendre. \textbf{Diphtongaison} \emph{e} $\rightarrow$ \emph{ie} : \esp{defiendo, defiendes, defiende, defendemos, defendéis, defienden}.
\item \esp{compromiso} = engagement (ou situation délicate). \textbf{Faux ami} : FR « compromis » = \esp{acuerdo, transacción}. Verbe « s'engager » = \esp{comprometerse} (diphtongaison \emph{e}$\rightarrow$\emph{ie} : \esp{me comprometo}).
\item \esp{la paz} : nom féminin (sauf \esp{el paz} archaïque). Pluriel rare \esp{las paces} (traité de paix). Presque toujours au singulier.
\end{itemize}
\end{attention}

\courssub{Carte mentale du lexique et stratégie de mémorisation}

La figure suivante organise les quatre champs lexicaux autour du thème central. Reproduis-la dans ton cahier et ajoute chaque semaine deux ou trois mots nouveaux découverts dans l'actualité.

\begin{popfigure}
\begin{center}
\begin{tikzpicture}[
  mindmap,
  grow cyclic,
  every node/.style={concept, align=center, font=\small, execute at begin node=\hskip0pt},
  concept color=popblue,
  level 1 concept/.append style={level distance=3.6cm, sibling angle=90, font=\small\bfseries},
  level 2 concept/.append style={level distance=2.6cm, sibling angle=45, font=\footnotesize}
]
\node[concept color=popblue, text=white, align=center]{Ciudadanos\\ del mundo}
  child[concept color=popgreen] { node {Ciudadanía}
    child { node[concept color=popgreenL, text=popdark, align=center]{derechos\\ deberes} }
    child { node[concept color=popgreenL, text=popdark, align=center]{voto\\ nacionalidad\\ ley} }
  }
  child[concept color=poporange] { node[align=center]{Derechos\\ humanos}
    child { node[concept color=poporangeL, text=popdark, align=center]{libertad\\ expresión} }
    child { node[concept color=poporangeL, text=popdark, align=center]{igualdad\\ justicia\\ dignidad} }
  }
  child[concept color=poppurple] { node[align=center]{Apertura\\ al mundo}
    child { node[concept color=poppurpleL, text=popdark, align=center]{frontera\\ migrante\\ extranjero} }
    child { node[concept color=poppurpleL, text=popdark, align=center]{solidaridad\\ diálogo\\ hermano} }
  }
  child[concept color=popgold] { node[align=center]{Compromiso\\ y paz}
    child { node[concept color=popgoldL, text=popdark, align=center]{luchar\\ defender} }
    child { node[concept color=popgoldL, text=popdark, align=center]{sembrar\\ paz\\ compartir} }
  };
\end{tikzpicture}
\end{center}
\legende{Carte mentale des quatre champs lexicaux du thème « Ciudadanos del mundo ».}
\end{popfigure}

\begin{methode}[Retenir le lexique : la phrase personnelle contextualisée]
Un mot isolé s'oublie vite ; un mot dans \textbf{ta} phrase ancrée dans ta réalité se retient. Pour chaque mot nouveau :
\begin{enumerate}
\item Lis l'exemple du cours et traduis-le mentalement.
\item Construis une phrase personnelle sur le Cameroun (ton lycée, ta ville, ton quartier, l'actualité).
\item Vérifie l'accord (genre/nombre) et la conjugaison, puis apprends la phrase entière à l'oral.
\end{enumerate}
\textbf{Modèles :}
\begin{itemize}
\item \esp{En Camerún, los ciudadanos votan para elegir a sus representantes.} (Au Cameroun, les citoyens votent pour élire leurs représentants.)
\item \esp{En mi liceo de Douala, aprendemos a respetar los derechos de los demás.} (Dans mon lycée de Douala, nous apprenons à respecter les droits des autres.)
\item \esp{Los jóvenes cameruneses quieren sembrar la paz en su país.} (Les jeunes Camerounais veulent semer la paix dans leur pays.)
\item \esp{La solidaridad entre los vecinos de mi barrio en Yaoundé es muy fuerte.} (La solidarité entre les voisins de mon quartier à Yaoundé est très forte.)
\end{itemize}
\end{methode}

\begin{exoresolu}[Lexique en contexte : complète et traduis]
\textbf{Énoncé.} Complète les phrases avec le mot qui convient (\esp{frontera, derechos, voto, paz, extranjero, deber}), puis traduis-les en français.
\begin{enumerate}
\item Los \ldots{} humanos son universales.
\item Para entrar en otro país, hay que cruzar la \ldots{}.
\item El \ldots{} es un derecho y también un deber del ciudadano.
\item Sin justicia no hay \ldots{} verdadera.
\item Respetar las leyes es un \ldots{} de todos.
\end{enumerate}
\tcblower
\textbf{Solution détaillée.}
\begin{enumerate}
\item \esp{Los \textbf{derechos} humanos son universales.} « Les droits de l'homme sont universels. »
\item \esp{Para entrar en otro país, hay que cruzar la \textbf{frontera}.} « Pour entrer dans un autre pays, il faut traverser la frontière. »
\item \esp{El \textbf{voto} es un derecho y también un deber del ciudadano.} « Le vote est un droit et aussi un devoir du citoyen. »
\item \esp{Sin justicia no hay \textbf{paz} verdadera.} « Sans justice, il n'y a pas de paix véritable. »
\item \esp{Respetar las leyes es un \textbf{deber} de todos.} « Respecter les lois est un devoir de tous. »
\end{enumerate}
\end{exoresolu}

\begin{exercice}[À toi de jouer]
\begin{enumerate}
\item Classe les mots suivants dans les quatre champs lexicaux du cours : \esp{igualdad, migrante, deber, luchar, diálogo, nacionalidad, dignidad, sembrar, ley, solidaridad}.
\item Écris trois phrases personnelles sur le Cameroun en utilisant : \esp{votar}, \esp{los derechos humanos}, \esp{la solidaridad}.
\item Traduis en espagnol :
\begin{enumerate}
\item La liberté d'expression est un droit fondamental.
\item Les migrants traversent la frontière pour chercher une vie meilleure.
\end{enumerate}
\end{enumerate}
\end{exercice}

\begin{corrige}[Exercice : À toi de jouer]
\textbf{1. Classification :}
\begin{itemize}
\item Citoyenneté : \esp{deber, nacionalidad, ley}
\item Droits humains : \esp{igualdad, dignidad}
\item Ouverture au monde : \esp{migrante, diálogo, solidaridad}
\item Engagement et paix : \esp{luchar, sembrar}
\end{itemize}

\textbf{2. Exemples de réponses personnelles (à adapter) :}
\begin{itemize}
\item \esp{En Camerún, los ciudadanos \textbf{votan} para elegir a sus representantes en las elecciones municipales.}
\item \esp{Los \textbf{derechos humanos} se enseñan en las escuelas camerunesas desde la primaria.}
\item \esp{La \textbf{solidaridad} entre los vecinos de mi barrio en Yaoundé se ve en los \emph{tontines} y en las fiestas comunes.}
\end{itemize}

\textbf{3. Traduction (Thème) :}
\begin{enumerate}
\item \esp{La libertad de expresión es un derecho fundamental.}
\item \esp{Los migrantes cruzan la frontera \textbf{para} buscar una vida mejor.} (\cle{Note} : \esp{para + infinitif} exprime le \textbf{but} / la finalité ; nous l'étudierons en détail dans la section \esp{Por y Para}.)
\end{enumerate}
\end{corrige}

\begin{aretenir}[L'essentiel de la section]
\begin{itemize}
\item Avant de traduire un poème, on vérifie toujours : \textbf{sens littéral}, \textbf{locuteur}, \textbf{images} (métaphores/symboles), \textbf{thème central}.
\item Le poème oppose deux séries de symboles : \esp{frontera / bandera / muro / patria} (division, fermeture) contre \esp{alba / pan compartido / hermanos / misma lengua} (union, ouverture, vie).
\item \textbf{Quatre champs lexicaux} à maîtriser pour le module 3 :
\begin{itemize}
\item Citoyenneté : \esp{ciudadanía, derechos, deberes, voto, nacionalidad, ley, elegir}.
\item Droits humains : \esp{libertad, expresión, igualdad, justicia, dignidad, respeto}.
\item Ouverture au monde : \esp{frontera, extranjero, migrante, solidaridad, diálogo, muro, patria, hermano}.
\item Engagement/Paix : \esp{luchar, sembrar, defender, paz, compromiso, compartir}.
\end{itemize}
\item \cle{Piège majeur} : \esp{el derecho} (droit juridique) $\neq$ \esp{derecho/a} (adj. droit) $\neq$ \esp{la derecha} (droite, direction/politique).
\item \cle{Grammaire} : Après \esp{soñé que} (passé + irréel) $\rightarrow$ \textbf{subjonctif imparfait} (\esp{fuera, se llamaran}).
\item \cle{Mémorisation} : Une phrase personnelle camerounaise par mot nouveau.
\end{itemize}
\end{aretenir}

Ce socle lexical est désormais solide. Nous pouvons passer à la section suivante : la \textbf{traduction littéraire} des vers du poème, où chacun de ces mots trouvera son équivalent français exact, avant d'aborder la grammaire de \esp{por} et \esp{para} et le débat guidé.

\coursec{Leçon E12 : Ciudadanos del mundo — Traduction poétique, POR / PARA et débat citoyen}

\courssub{Traduction de vers : méthode et pratique}

Dans la section précédente, nous avons lu et compris le poème engagé \esp{Ciudadanos del mundo}, et nous avons constitué notre lexique thématique de la citoyenneté (\esp{la bandera}, \esp{el pan compartido}, \esp{el camino}, \esp{la patria}, \esp{los derechos}, \esp{la frontera}, \esp{el pueblo}...). Nous voici maintenant devant une tâche plus exigeante et plus délicate : \cle{traduire} ces vers en français. Traduire un poème n'est pas traduire une phrase administrative : il faut transporter à la fois le \cle{sens}, les \cle{images} et la \cle{musique} du texte. C'est un exercice qui tombe régulièrement au baccalauréat (version et thème), et qui développe une compétence précieuse : la précision dans les deux langues.

\begin{definition}[Traduire un poème]
Traduire un poème, c'est produire dans la langue d'arrivée un texte qui respecte trois fidélités hiérarchisées :
\begin{itemize}
\item la \textbf{fidélité au sens} : ne rien ajouter, ne rien retirer, ne pas déformer l'idée du poète ;
\item la \textbf{fidélité aux images} : conserver si possible les métaphores et les symboles (\esp{el pan compartido} = le pain partagé, symbole de fraternité) ;
\item la \textbf{fidélité au rythme} : chercher en français une phrase qui « sonne » bien, fluide, naturelle, sans calque de la syntaxe espagnole.
\end{itemize}
Quand les trois fidélités entrent en conflit, le sens prime toujours, puis l'image, puis le rythme.
\end{definition}

Observons d'abord deux vers du texte support et leurs traductions possibles.

\begin{exemplebox}[Deux vers, deux traductions]
Vers espagnol : \esp{Mi bandera, un pan compartido.}

\textbf{Version littérale} : « Mon drapeau, un pain partagé. »

\textbf{Version soignée} : « Pour drapeau, j'ai un pain partagé. »

\textbf{Commentaire} : la version littérale est correcte mais sèche ; elle laisse le lecteur français deviner la relation entre « mon drapeau » et « un pain ». La version soignée explicite la construction elliptique du vers (le poète sous-entend : « j'ai pour drapeau... » ou « mon drapeau est... ») et rend le français plus naturel, sans trahir l'image.
\end{exemplebox}

\begin{exemplebox}[Un deuxième exemple]
Vers espagnol : \esp{Camino por el mundo con las manos abiertas.}

\textbf{Version littérale fautive} : « Je chemine par le monde avec les mains ouvertes. »

\textbf{Version soignée} : « Je marche à travers le monde, les mains ouvertes. »

\textbf{Commentaire} : \esp{caminar} se traduit le plus souvent par « marcher » ; « cheminer » est un faux raffinement. \esp{por} exprime ici le déplacement à travers un espace : « à travers », et non « par ». Le complément \esp{con las manos abiertas} devient plus élégant avec l'apposition « les mains ouvertes ».
\end{exemplebox}

\begin{attention}[Le piège du mot à mot]
Ne traduisez \textbf{jamais} un poème mot à mot. Le mot à mot produit des contresens ou du « fragnol ». Exemples :
\begin{itemize}
\item \esp{camino por el mundo} $\neq$ « je chemine par le monde », mais « je marche à travers le monde » ;
\item \esp{mi patria no tiene fronteras} $\neq$ « ma patrie n'a pas de frontières » (acceptable) $>$ mieux : « ma patrie est sans frontières » ;
\item \esp{somos ciudadanos del mundo} $\neq$ « nous sommes citoyens du monde » est ici correct, car le mot à mot tombe juste : preuve qu'il faut \emph{vérifier}, pas interdire systématiquement.
\end{itemize}
Règle d'or : après votre traduction, relisez la phrase française \textbf{en oubliant l'espagnol}. Si un francophone qui ne connaît pas l'espagnol la trouve bizarre, elle est mauvaise.
\end{attention}

\courssub{La méthode en quatre étapes}

\begin{methode}[Traduire un poème en 4 étapes]
\begin{enumerate}
\item \textbf{Le sens exact} : lisez le vers entier, identifiez le sujet, le verbe, les compléments ; vérifiez chaque mot dans le contexte du poème (un mot peut avoir plusieurs sens : \esp{pan} = pain, mais aussi, dans d'autres contextes, métaphore de la nourriture, du partage).
\item \textbf{Les images} : repérez les métaphores et les symboles ; demandez-vous si l'image existe aussi en français (\esp{un pan compartido} : oui, le pain partagé évoque la fraternité en français aussi).
\item \textbf{Le rythme et les sonorités} : comptez approximativement les syllabes, cherchez une phrase française fluide, évitez les lourdeurs (répétitions de « de », enchaînements de pronoms).
\item \textbf{La relecture à voix haute en français} : lisez votre traduction à voix haute ; si elle « accroche » l'oreille ou la grammaire, corrigez-la. Vérifiez enfin que rien du sens original n'a été perdu.
\end{enumerate}
\end{methode}

\begin{popfigure}
\begin{tikzpicture}[
  node distance=0.5cm,
  etape/.style={draw=popblue, fill=popblueL, rounded corners, align=center, text width=3.2cm, minimum height=1.1cm, font=\small},
  fleche/.style={-{Stealth[length=3mm]}, thick, poporange}
]
\node[etape, align=center] (v){\textbf{Vers espagnol}\\ \esp{Mi bandera, un pan compartido}};
\node[etape, below=of v, align=center] (c){\textbf{1. Compréhension}\\ sujet, verbe, sens des mots-clés};
\node[etape, below=of c, align=center] (l){\textbf{2. Version littérale}\\ « mon drapeau, un pain partagé »};
\node[etape, below=of l, align=center] (s){\textbf{3. Version soignée}\\ « pour drapeau, j'ai un pain partagé »};
\node[etape, below=of s, fill=popgreenL, draw=popgreen, align=center] (ver){\textbf{4. Vérification}\\ sens, grammaire, style};
\draw[fleche] (v) -- (c);
\draw[fleche] (c) -- (l);
\draw[fleche] (l) -- (s);
\draw[fleche] (s) -- (ver);
\draw[fleche, dashed, popmuted] (ver.west) to[bend right=60] node[left, font=\footnotesize, align=center]{si problème,\\ on recommence} (c.west);
\end{tikzpicture}
\legende{Organigramme de la traduction poétique : du vers espagnol à la version française vérifiée.}
\end{popfigure}

\courssub{Les choix du traducteur : garder ou adapter une métaphore ?}

Toute traduction poétique est une série de décisions. Face à une image espagnole, le traducteur a trois options :

\begin{center}
\begin{tabularx}{\linewidth}{|X|X|X|}
\hline
\rowcolor{popblueL} \textbf{Option} & \textbf{Quand l'employer} & \textbf{Exemple} \\
\hline
Garder l'image telle quelle & L'image est claire et naturelle en français & \esp{las manos abiertas} $\rightarrow$ « les mains ouvertes » \\
\hline
Adapter légèrement & L'image existe mais la formulation française diffère & \esp{camino por el mundo} $\rightarrow$ « je marche à travers le monde » \\
\hline
Remplacer par une image équivalente & L'image littérale serait incompréhensible en français & \esp{tener el corazón en la mano} $\rightarrow$ « avoir le cœur sur la main » (et non « dans la main ») \\
\hline
\end{tabularx}
\end{center}

\begin{propriete}[L'ordre des mots en poésie espagnole]
En poésie, l'ordre des mots espagnol est beaucoup plus libre qu'en français : le poète peut placer le verbe avant le sujet, le complément avant le verbe, pour des raisons de rythme ou de mise en relief. Le traducteur français doit souvent \textbf{rétablir l'ordre sujet -- verbe -- complément}, obligatoire en français.
\begin{itemize}
\item \esp{Mi bandera, un pan compartido.} (complément d'abord, sujet implicite) $\rightarrow$ « Pour drapeau, j'ai un pain partagé » ou « J'ai pour drapeau un pain partagé ».
\item \esp{Con las manos abiertas camino.} $\rightarrow$ « Je marche les mains ouvertes » (et non « avec les mains ouvertes je marche », calque lourd).
\end{itemize}
Attention : cette liberté est \emph{poétique}. Dans la prose courante, l'espagnol suit, comme le français, l'ordre sujet -- verbe -- complément.
\end{propriete}

\courssub{Version guidée : traduction collective des vers 3 à 5}

Reprenons les vers 3 à 5 du poème support et appliquons la méthode ensemble.

\begin{textebox}[Ciudadanos del mundo, vers 3 à 5]
Mi patria no tiene fronteras,\\
camino por el mundo con las manos abiertas,\\
y mi lengua es la de todos los pueblos.
\end{textebox}

\textbf{Glossaire} : \esp{la patria} = la patrie ; \esp{la frontera} = la frontière ; \esp{caminar} = marcher ; \esp{las manos abiertas} = les mains ouvertes ; \esp{la lengua} = la langue ; \esp{el pueblo} = le peuple.

\textbf{Étape 1 -- sens} : le poète affirme que sa patrie dépasse les frontières, qu'il se déplace dans le monde en toute ouverture, et que sa langue est celle de l'humanité entière.

\textbf{Étape 2 -- version littérale} : « Ma patrie n'a pas de frontières, / je marche par le monde avec les mains ouvertes, / et ma langue est celle de tous les peuples. »

\textbf{Étape 3 -- version soignée} : « Ma patrie est sans frontières, / je marche à travers le monde, les mains ouvertes, / et ma langue est celle de tous les peuples. »

\textbf{Étape 4 -- vérification} : sens intégral conservé ; « à travers » rend correctement \esp{por} ; l'apposition « les mains ouvertes » allège la phrase ; le rythme français est fluide.

\courssub{Thème : du français vers l'espagnol}

Le thème poétique est un excellent complément d'entraînement pour le Bac : il oblige à \emph{produire} en espagnol les images que l'on a apprises à reconnaître.

\begin{exoresolu}[Thème : deux vers français à traduire en espagnol]
Énoncé : traduisez en espagnol les deux vers suivants : « Nous sommes les enfants de la même terre, / et notre pain est le pain de tous. »
\tcblower
Solution détaillée :
\begin{itemize}
\item « Nous sommes » $\rightarrow$ \esp{Somos} (verbe \esp{ser}, identité : pas \esp{estar}).
\item « les enfants de la même terre » $\rightarrow$ \esp{los hijos de la misma tierra} (\esp{tierra} = terre, monde ; \esp{misma} s'accorde avec \esp{tierra}, féminin singulier).
\item « notre pain est le pain de tous » $\rightarrow$ \esp{nuestro pan es el pan de todos} (\esp{ser} pour une définition ; possessif \esp{nuestro} accordé avec \esp{pan}, masculin singulier).
\end{itemize}
Traduction complète : \esp{Somos los hijos de la misma tierra, / y nuestro pan es el pan de todos.}
Pièges évités : pas de \esp{estamos} (identité, pas état) ; pas de \esp{mismo tierra} (accord féminin) ; pas d'article devant \esp{todos} après \esp{de}.
\end{exoresolu}

\courssub{Correction collective : comparer et justifier}

Comparer plusieurs propositions de traduction est la meilleure façon de progresser. Voici trois propositions d'élèves pour le vers \esp{Camino por el mundo con las manos abiertas} :

\begin{center}
\begin{tabularx}{\linewidth}{|X|X|X|}
\hline
\rowcolor{popblueL} \textbf{Proposition} & \textbf{Qualités} & \textbf{Défauts} \\
\hline
« Je chemine par le monde avec les mains ouvertes. » & Respecte l'ordre espagnol & « chemine » trop littéraire ; « par le monde » est un calque de \esp{por} ; « avec » alourdit \\
\hline
« Je vais dans le monde, mains ouvertes. » & Fluide & « je vais » affaiblit \esp{camino} (marcher, avancer) ; « mains ouvertes » sans article est trop elliptique en français \\
\hline
« Je marche à travers le monde, les mains ouvertes. » & Sens exact, image conservée, rythme naturel & Aucun défaut majeur \\
\hline
\end{tabularx}
\end{center}

La troisième proposition est la meilleure : elle respecte les trois fidélités (sens, image, rythme). En correction collective, exigez toujours une \textbf{justification} : un bon traducteur sait dire \emph{pourquoi} il a choisi telle formulation.

\begin{exercice}[À vous de traduire]
Traduisez en français soigné les vers suivants, en appliquant les quatre étapes de la méthode :
\begin{enumerate}
\item \esp{Mi bandera, un pan compartido.}
\item \esp{En mi mesa hay sitio para el que llega.}
\item \esp{No conozco fronteras, conozco rostros.}
\end{enumerate}
Puis traduisez en espagnol : « Mon cœur n'a pas de frontières, et ma maison est ouverte à tous. »
\end{exercice}

\begin{corrige}[Exercice : traduction de vers]
\begin{enumerate}
\item « Pour drapeau, j'ai un pain partagé. » (ou : « Mon drapeau, c'est un pain partagé. »)
\item « À ma table, il y a une place pour celui qui arrive. » (\esp{hay sitio} = il y a de la place ; \esp{el que llega} = celui qui arrive, le nouveau venu.)
\item « Je ne connais pas de frontières, je connais des visages. » (\esp{rostros} = visages ; le parallélisme \esp{no conozco... conozco...} doit être conservé.)
\end{enumerate}
Thème : \esp{Mi corazón no tiene fronteras, y mi casa está abierta a todos.} (On accepte \esp{es abierta} si l'élève justifie un trait permanent, mais \esp{estar} est préférable pour un état concret : la porte est ouverte.)
\end{corrige}

\begin{aretenir}[L'essentiel de la section traduction]
\begin{itemize}
\item Traduire un poème exige trois fidélités : le \textbf{sens}, les \textbf{images}, le \textbf{rythme} ; en cas de conflit, le sens prime.
\item La méthode : comprendre $\rightarrow$ version littérale $\rightarrow$ version soignée $\rightarrow$ relecture à voix haute en français.
\item Jamais de mot à mot : \esp{camino por el mundo} = « je marche à travers le monde ».
\item L'ordre des mots poétique espagnol est libre ; le français exige sujet -- verbe -- complément.
\item Au thème, surveillez \esp{ser}/\esp{estar}, les accords et les articles : ce sont les fautes les plus sanctionnées.
\end{itemize}
\end{aretenir}

\courssub{Grammaire : POR et PARA — Règles, emplois et pièges}

Le poème support utilise \esp{por} (\esp{camino \textbf{por} el mundo}) et \esp{para} (\esp{hay sitio \textbf{para} el que llega}). Ces deux prépositions se traduisent souvent par « pour » ou « par » en français, mais leurs emplois sont distincts et non interchangeables. Maîtriser cette opposition est indispensable pour le thème, la version et l'expression orale.

\begin{definition}[Por vs Para : distinction fondamentale]
\begin{itemize}
\item \textbf{POR} exprime la \cle{cause}, le \cle{moyen}, le \cle{mouvement à travers}, la \cle{durée}, l'\cle{échange}, l'\cle{agent} (passif), la \cle{distribution}. Idée de \emph{parcours}, de \emph{raison d'être}, de \emph{substitution}.
\item \textbf{PARA} exprime la \cle{finalité}, le \cle{but}, le \cle{destinataire}, la \cle{destination}, l'\cle{échéance}, l'\cle{opinion}, la \cle{comparison}. Idée de \emph{projection vers l'avant}, de \emph{flèche} orientée vers un objectif.
\end{itemize}
\end{definition}

\begin{propriete}[Emplois principaux de POR]
\begin{center}
\begin{tabularx}{\linewidth}{|l|X|X|}
\hline
\rowcolor{popblueL} \textbf{Valeur} & \textbf{Explication} & \textbf{Exemple traduit} \\
\hline
Cause / Motif & Pourquoi on fait l'action (raison antérieure) & \esp{Lo hice \textbf{por} ti.} (Je l'ai fait \textbf{pour} toi / \textbf{à cause de} toi.) \\
\hline
Moyen / Instrument & Par quel moyen l'action se fait & \esp{Te llamo \textbf{por} teléfono.} (Je t'appelle \textbf{par} téléphone.) \\
\hline
Mouvement / Lieu & Traversée, déplacement \emph{à travers} & \esp{Caminamos \textbf{por} el parque.} (Nous marchons \textbf{à travers} le parc.) \\
\hline
Durée / Moment vague & Pendant combien de temps, moment de la journée & \esp{Estudié \textbf{por} tres horas.} (J'ai étudié \textbf{pendant} trois heures.) \\
\hline
Échange / Substitution & Échanger une chose contre une autre & \esp{Te doy mi libro \textbf{por} el tuyo.} (Je te donne mon livre \textbf{contre} le tien.) \\
\hline
Agent (passif) & Par qui l'action est faite & \esp{Fue escrito \textbf{por} Lorca.} (Il a été écrit \textbf{par} Lorca.) \\
\hline
Distribution & Répartition (par tête, par jour) & \esp{Una vez \textbf{por} semana.} (Une fois \textbf{par} semaine.) \\
\hline
\end{tabularx}
\end{center}
\end{propriete}

\begin{propriete}[Emplois principaux de PARA]
\begin{center}
\begin{tabularx}{\linewidth}{|l|X|X|}
\hline
\rowcolor{poporangeL} \textbf{Valeur} & \textbf{Explication} & \textbf{Exemple traduit} \\
\hline
Finalité / But & Pour quoi on fait l'action (objectif futur) + infinitif & \esp{Estudio \textbf{para} aprobar.} (J'étudie \textbf{pour} réussir.) \\
\hline
Destinataire / Bénéficiaire & À qui est destiné l'objet/l'action & \esp{Este regalo es \textbf{para} ti.} (Ce cadeau est \textbf{pour} toi.) \\
\hline
Destination & Vers où l'on se dirige & \esp{Salgo \textbf{para} Madrid.} (Je pars \textbf{pour} Madrid.) \\
\hline
Échéance / Date limite & Moment limite (deadline) & \esp{La tarea es \textbf{para} el lunes.} (Le devoir est \textbf{pour} lundi.) \\
\hline
Opinion / Jugement & Selon quelqu'un & \esp{\textbf{Para} mí, es fácil.} (\textbf{Pour} moi, c'est facile.) \\
\hline
Comparaison / Exception & Par rapport à une norme & \esp{Es alto \textbf{para} su edad.} (Il est grand \textbf{pour} son âge.) \\
\hline
\end{tabularx}
\end{center}
\end{propriete}

\begin{attention}[Pièges majeurs pour francophones]
\begin{itemize}
\item \textbf{« Pour » = POR ou PARA ?} \esp{Lo hice \textbf{por} ti} (cause/motif : à cause de toi, en ton honneur) vs \esp{Lo hice \textbf{para} ti} (destinataire : pour que tu l'aies, à ton intention).
\item \textbf{« Par » = POR (surtout) :} \esp{por teléfono, por correo, por avión}. Mais « par » agent passif = \esp{por}.
\item \textbf{« Pendant » = POR (durée accomplie) :} \esp{Viví allí \textbf{por} dos años.} Mais \esp{PARA} si c'est une durée prévue/requise : \esp{Necesito el coche \textbf{para} dos días} (il me faut la voiture pour deux jours / deadline).
\item \textbf{Verbes de mouvement :} \esp{IR / SALIR / PARTIR + PARA} = destination (\esp{Salgo para Yaoundé}). \esp{PASAR / CAMINAR / VIAJAR + POR} = traverser (\esp{Viajamos por Douala}).
\item \textbf{Expressions figées à connaître par cœur :}
\begin{itemize}
\item \esp{Por favor, por supuesto, por eso, por fin, por lo general, por escrito, por casualidad, por ciento, por ciento, multiplicar por, dividir por.}
\item \esp{Para siempre, para variar, para empezar, para colmo, estar para + infinitif (être sur le point de), ser para (être bon pour / destiné à).}
\end{itemize}
\end{itemize}
\end{attention}

\begin{exemplebox}[Contraste POR / PARA sur le même verbe]
\begin{itemize}
\item \esp{Compro pan \textbf{para} el desayuno.} (Finalité : le pain \emph{est destiné au} petit-déjeuner.)
\item \esp{Compro pan \textbf{por} mi madre.} (Substitution : j'achète \emph{à la place de} ma mère / \emph{pour le compte de} ma mère.)
\item \esp{Voy \textbf{para} la escuela.} (Destination : je me dirige \emph{vers} l'école.)
\item \esp{Voy \textbf{por} la escuela.} (Trajet : je passe \emph{par} l'école / je traverse l'école.)
\end{itemize}
\end{exemplebox}

\begin{exoresolu}[Choix entre POR et PARA]
Énoncé : Complétez par \esp{por} ou \esp{para}.
\begin{enumerate}
\item Este libro fue escrito \_\_\_\_\_\_ Gabriel García Márquez.
\item Salimos \_\_\_\_\_\_ la estación a las ocho.
\item Este regalo es \_\_\_\_\_\_ mi hermana.
\item He estado esperando \_\_\_\_\_\_ dos horas.
\item Lo hice \_\_\_\_\_\_ ti, no \_\_\_\_\_\_ mí.
\end{enumerate}
\tcblower
Solution détaillée :
\begin{enumerate}
\item \textbf{por} (agent passif : \emph{écrit par}).
\item \textbf{para} (destination : \emph{partir pour/vers} la gare).
\item \textbf{para} (destinataire : \emph{pour} ma sœur).
\item \textbf{por} (durée accomplie : \emph{pendant} deux heures).
\item \textbf{por} (cause/motif : \emph{à cause de / pour} toi), \textbf{para} (opposition / comparaison : \emph{au lieu de} moi / \emph{pour} moi). \textit{Note :} \esp{Lo hice por ti, no para mí} = Je l'ai fait pour toi (à cause de toi), pas pour moi.
\end{enumerate}
\end{exoresolu}

\begin{exercice}[Entraînement : Por ou Para ?]
Complétez les phrases suivantes par \esp{por} ou \esp{para}. Justifiez brièvement votre choix (cause, finalité, destinataire, durée, destination, agent, moyen...).
\begin{enumerate}
\item Mi padre trabaja \_\_\_\_\_\_ una empresa multinacional.
\item El tren pasa \_\_\_\_\_\_ el túnel.
\item Esta carta es \_\_\_\_\_\_ ti.
\item Estudiamos \_\_\_\_\_\_ aprender.
\item \_\_\_\_\_\_ mí, el examen fue difícil.
\item Vamos a quedarnos \_\_\_\_\_\_ una semana.
\item El puente fue construido \_\_\_\_\_\_ los romanos.
\item Pagué 1000 francos \_\_\_\_\_\_ el libro.
\item Voy \_\_\_\_\_\_ Douala mañana.
\item \_\_\_\_\_\_ favor, abre la ventana.
\end{enumerate}
\end{exercice}


\begin{corrige}[Exercice : Por ou Para ?]
\begin{enumerate}
\item \textbf{para} (destination / lieu de travail vu comme but) — \textit{Note :} \esp{trabajar en} aussi possible, mais \esp{para} insiste sur l'employeur/bénéficiaire.
\item \textbf{por} (mouvement à travers : le train traverse le tunnel).
\item \textbf{para} (destinataire : cette lettre est pour toi).
\item \textbf{para} (finalité : étudier pour/apprendre).
\item \textbf{Para} (opinion : pour moi / selon moi).
\item \textbf{por} (durée : rester pendant une semaine).
\item \textbf{por} (agent passif : construit par les Romains).
\item \textbf{por} (échange / prix : payer 1000 francs \emph{contre} le livre).
\item \textbf{para} (destination : je vais vers Douala).
\item \textbf{Por} (expression figée : \esp{por favor}).
\end{enumerate}
\end{corrige}

\begin{aretenir}[L'essentiel : POR vs PARA]
\begin{itemize}
\item \textbf{POR} = flèche arrière / horizontalité : cause, moyen, traversée, durée, échange, agent, distribution. Question clé : \emph{« Pourquoi ? Par quoi ? Combien de temps ? »}
\item \textbf{PARA} = flèche avant / verticalité : but, destinataire, destination, échéance, opinion, comparaison. Question clé : \emph{« Pour quoi faire ? Pour qui ? Jusqu'à quand ? Vers où ? »}
\item Au thème : \esp{« pour »} $\rightarrow$ analysez la logique (cause ? $\rightarrow$ \esp{por} ; but ? $\rightarrow$ \esp{para}).
\item Expressions figées : à mémoriser (\esp{por favor, para siempre, por escrito, para empezar...}).
\end{itemize}
\end{aretenir}

\courssub{Débat guidé : « Ciudadanos del mundo »}

Le programme exige la compétence « débattre des droits et de la citoyenneté ». Le poème \esp{Ciudadanos del mundo} est le déclencheur idéal. Le débat guidé structure la prise de parole, oblige à l'argumentation et réactive le lexique et la grammaire (subjonctif pour l'opinion, \esp{por/para} pour l'argumentation).

\begin{methode}[Conduire un débat guidé en 5 étapes]
\begin{enumerate}
\item \textbf{Lecture du déclencheur} : texte, citation, image, statistique (ici : le poème).
\item \textbf{Formulation de la motion} : une phrase affirmative claire, sujette à discussion. Exemple : \esp{« La ciudadanía mundial es una utopía necesaria. »} ou \esp{« Las fronteras son obsoletas en el siglo XXI. »}
\item \textbf{Préparation par groupes} (5--10 min) : chaque groupe (Pour / Contre / Nuancés) prépare 3 arguments avec connecteurs logiques (\esp{porque, ya que, por un lado... por otro, sin embargo, además, en conclusión}).
\item \textbf{Débat oral} (15--20 min) : tour de table, modérateur (élève ou prof), respect du temps de parole, reformulation obligatoire de l'argument adverse avant de répondre.
\item \textbf{Synthèse écrite} : chaque élève rédige un paragraphe personnel (80--100 mots) reprenant sa position finale, en utilisant le lexique et \esp{por/para}.
\end{enumerate}
\end{methode}

\begin{textebox}[Sujet de débat : Motion proposée]
\textbf{Moción :} \esp{« Ser ciudadano del mundo implica renunciar a la identidad nacional. »}
(Être citoyen du monde implique de renoncer à son identité nationale.)

\textbf{Contexte camerounais :} Au Cameroun, nous sommes fiers de nos identités locales (Bamiléké, Bassa, Douala, Peul, etc.) et nationale (Cameroun), tout en étant ouverts à la CEMAC, à l'UA, à la Francophonie, au Commonwealth. Ce débat nous concerne directement.
\end{textebox}

\begin{experience}[Mise en œuvre : Débat « Ciudadanos del mundo »]
\textbf{Organisation :}
\begin{itemize}
\item \textbf{Groupe A (Pour la motion)} : La citoyenneté mondiale dilue les responsabilités concrètes ; on ne peut pas être solidaire de 8 milliards d'humains concrètement ; l'identité nationale est le socle des droits réels (vote, protection consulaire, service militaire/civique).
\item \textbf{Groupe B (Contre la motion)} : Identités multiples et non exclusives (je suis Bassa, Camerounais, Africain, Citoyen du monde) ; les défis (climat, pandémies, migration) exigent une conscience planétaire ; \esp{por} la dignidad humana, \esp{para} un futuro compartido.
\item \textbf{Groupe C (Juges / Observateurs)} : Notent les arguments, l'usage correct de \esp{por/para}, du subjonctif (\esp{es importante que...}, \esp{no creo que...}), des connecteurs. Ils rendent un verdict argumenté.
\end{itemize}
\textbf{Expressions utiles au tableau :}
\esp{Estoy de acuerdo / No estoy de acuerdo. -- Me parece que... -- Desde mi punto de vista... -- Por un lado... por otro lado... -- Sin embargo / No obstante... -- Además / Asimismo... -- En conclusión... -- ¿Qué opinas tú? -- ¿Me permites una intervención?}
\end{experience}

\begin{savaistu}[Cameroun, Guinée équatoriale et citoyenneté]
La Guinée équatoriale (pays hispanophone voisin) est membre de la CEMAC et de l'UA comme le Cameroun. Ses citoyens ont la double nationalité fréquente (espagnole/équato-guinéenne). Le débat sur la « ciudadanía mundial » y est vif : entre l'héritage espagnol, l'appartenance africaine (Fang, Bubi) et la mondialisation pétrolière. Au Cameroun, la « double nationalité » est un sujet constitutionnel sensible (Code de la nationalité, loi de 1968 modifiée). Être « citoyen du monde » ne signifie pas juridiquement « apatride », mais moralement « responsable au-delà des frontières ».
\end{savaistu}

\begin{exercice}[Production écrite : Synthèse du débat]
Rédigez en espagnol un paragraphe de 80--100 mots donnant votre opinion personnelle sur la motion : \esp{« Ser ciudadano del mundo implica renunciar a la identidad nacional. »}
Utilisez au moins : 3 connecteurs logiques, 2 verbes au subjonctif (après expression d'opinion/doute), 1 emploi correct de \esp{por} et 1 de \esp{para}.
\end{exercice}


\begin{corrige}[Exercice : Production écrite (proposition de corrigé type)]
\esp{En mi opinión, ser ciudadano del mundo no implica renunciar a la identidad nacional, sino enriquecerla. Por un lado, nuestras raíces (nuestra patria, nuestra lengua, nuestras costumbres) nos dan un punto de apoyo firme. Por otro lado, los grandes retos de hoy -- el cambio climático, la migración, las pandemias -- no entienden de fronteras y exigen una solidaridad planetaria. Es fundamental que entendamos que la humanidad es una sola familia. Para mí, ser camerunés y ciudadano del mundo son compatibles: amo mi tierra \textbf{por} su cultura y trabajo \textbf{para} un futuro mejor para todos.}
\end{corrige}

\begin{aretenir}[L'essentiel de la leçon E12]
\begin{itemize}
\item \textbf{Traduction poétique} : 3 fidélités (sens, image, rythme) ; méthode en 4 étapes ; pas de mot à mot ; rétablir l'ordre SVO en français.
\item \textbf{POR} : cause, moyen, traversée, durée, échange, agent, distribution. \textbf{PARA} : but, destinataire, destination, échéance, opinion, comparaison.
\item \textbf{Débat citoyen} : motion claire, préparation arguments (Pour/Contre), connecteurs, subjonctif, \esp{por/para}, synthèse écrite personnelle.
\item \textbf{Compétences Bac} : Version/Thème (vers), Grammaire en contexte (por/para), Expression orale (débat), Expression écrite (essai argumenté).
\end{itemize}
\end{aretenir}

\coursec{Leçon E12 : Traduction de vers, POR / PARA, Débat « Ciudadanos del mundo »}

\coursec{Texte support : un poème engagé}

\begin{textebox}[Poème original : « Ciudadano del mundo »]
\esp{Camino por el mundo sin fronteras,} \\
\esp{Busco la justicia, siembro verdades.} \\
\esp{No tengo patria, ni banderas,} \\
\esp{Solo la tierra y sus libertades.} \\
\esp{Lucho por la paz, por la dignidad,} \\
\esp{Para que el niño tenga un pan,} \\
\esp{Para que el hombre sea verdad,} \\
\esp{Para un mundo sin más fronteras.}
\end{textebox}

\begin{definition}[Glossaire du poème]
\begin{itemize}
\item \esp{camino} (verbe \esp{caminar}, 1ère pers. sg. présent) : je chemine, je marche.
\item \esp{sin fronteras} : sans frontières.
\item \esp{busco} (verbe \esp{buscar}) : je cherche.
\item \esp{siembro} (verbe \esp{sembrar}) : je sème.
\item \esp{verdades} : vérités (pluriel de \esp{verdad}).
\item \esp{patria} : patrie, pays natal.
\item \esp{dignidad} : dignité.
\item \esp{pan} : pain (symbole de nourriture, de survie).
\end{itemize}
\end{definition}

\begin{exemplebox}[Traduction littéraire proposée]
\esp{Je chemine à travers le monde sans frontières,} \\
\esp{Je cherche la justice, je sème des vérités.} \\
\esp{Je n'ai ni patrie, ni drapeaux,} \\
\esp{Seulement la terre et ses libertés.} \\
\esp{Je lutte pour la paix, pour la dignité,} \\
\esp{Afin que l'enfant ait du pain,} \\
\esp{Afin que l'homme soit vérité,} \\
\esp{Pour un monde sans plus de frontières.}
\end{exemplebox}

\courssub{Compréhension et lexique thématique : la citoyenneté mondiale}

\begin{exercice}[Questions de compréhension]
\begin{enumerate}
\item Quel est le statut du « je » poétique ? Justifiez par deux vers.
\item Identifiez les verbes d'action au présent de l'indicatif. Que traduisent-ils sur l'attitude du citoyen du monde ?
\item Expliquez l'opposition entre les vers 3-4 et le vers 8.
\item Relevez les deux occurrences de \esp{por} et les trois de \esp{para} dans le poème. Prédisez leur valeur (cause, but, etc.) avant d'étudier la règle.
\end{enumerate}
\end{exercice}

\begin{corrige}[Exercice de compréhension]
\begin{enumerate}
\item Le « je » est un \textbf{citoyen du monde} (\esp{ciudadano del mundo}), apatride par choix (\esp{No tengo patria, ni banderas}) mais engagé (\esp{Lucho por la paz}).
\item \esp{Camino, busco, siembro, tengo, lucha} (3e pers.). Attitude active, volontaire, en mouvement perpétuel.
\item Vers 3-4 : refus des appartenances étroites (patrie, drapeaux). Vers 8 : but positif, construction d'un monde uni (\esp{Para un mundo sin más fronteras}). Passage du négatif (refus) au positif (projet).
\item \esp{por} (v.5 x2) : cause/motif (\esp{lucho por la paz}). \esp{para} (v.6, 7, 8) : finalité/but (\esp{para que...}, \esp{para un mundo}).
\end{enumerate}
\end{corrige}

\begin{aretenir}[Lexique clé : Ciudadanía y derechos]
\begin{center}
\begin{tabularx}{\linewidth}{|X|X|}
\hline
\rowcolor{popblueL} Español & Français \\
\hline
\esp{la ciudadanía} & la citoyenneté \\
\hline
\esp{los derechos humanos} & les droits de l'homme \\
\hline
\esp{la justicia social} & la justice sociale \\
\hline
\esp{la igualdad} & l'égalité \\
\hline
\esp{la solidaridad} & la solidarité \\
\hline
\esp{las fronteras} & les frontières \\
\hline
\esp{la identidad} & l'identité \\
\hline
\esp{la pertenencia} & l'appartenance \\
\hline
\esp{un ciudadano del mundo} & un citoyen du monde \\
\hline
\esp{comprometido/a} & engagé(e) \\
\hline
\end{tabularx}
\end{center}
\end{aretenir}

\courssub{Traduction de vers : méthode et pratique}

\begin{methode}[Comment traduire un vers poétique ?]
\begin{enumerate}
\item \textbf{Analyser la structure} : mètre, rimes, enjambements, ponctuation.
\item \textbf{Identifier le sens littéral} : grammaire, vocabulaire, syntaxe (sujet, verbe, COD, CCL).
\item \textbf{Résoudre les problèmes de langue} : \esp{por/para}, \esp{ser/estar}, subjonctif (\esp{para que + subj.}), champs lexicaux, connotations.
\item \textbf{Choisir le registre français} : soutenu, poétique, respectant le rythme si possible (vers libres ou rimés).
\item \textbf{Relire à voix haute} : vérifier la musicalité et l'émotion.
\end{enumerate}
\end{methode}

\begin{exoresolu}[Traduction commentée : vers 5 à 8]
Énoncé. Traduisez et justifiez les choix pour : \esp{Lucho por la paz, por la dignidad. / Para que el niño tenga un pan. / Para que el hombre sea verdad. / Para un mundo sin más fronteras.}
\tcblower
Solution détaillée :
\begin{enumerate}
\item \esp{Lucho por la paz, por la dignidad.} $\rightarrow$ « Je lutte pour la paix, pour la dignité. » \esp{Por} = cause, motif profond qui pousse à l'action (¿por qué lucho?).
\item \esp{Para que el niño tenga un pan.} $\rightarrow$ « Afin que l'enfant ait du pain. » \esp{Para que} + \textbf{subjonctif présent} (\esp{tenga} de \esp{tener}) = but, finalité. « Afin que » rend mieux la finalité forte que « pour que ».
\item \esp{Para que el hombre sea verdad.} $\rightarrow$ « Afin que l'homme soit vérité / vrai. » \esp{Sea} = subjonctif de \esp{ser} (identité essentielle). \esp{Verdad} = vérité (nom) ou « vrai » (attribut).
\item \esp{Para un mundo sin más fronteras.} $\rightarrow$ « Pour un monde sans plus de frontières. » \esp{Para} + nom = destination/but final du combat.
\end{enumerate}
\end{exoresolu}

\coursec{Grammaire : POR et PARA}

Dans le poème, deux vers ont retenu notre attention : \esp{camino \textbf{por} el mundo} (« je chemine \textbf{à travers} le monde ») et \esp{\textbf{para} sembrar palabras de paz} / \esp{\textbf{para} que el niño tenga...} (« \textbf{pour} semer... », « \textbf{afin que} l'enfant... »). Pourquoi le poète a-t-il choisi \esp{por} dans le premier cas et \esp{para} dans les seconds, alors que le français utilise souvent « pour » ou des prépositions voisines ? C'est toute la question de cette section : apprendre à choisir correctement entre \cle{por} et \cle{para}, l'une des difficultés majeures pour les hispanisants francophones.

\courssub{Observation : por et para dans le poème}

\begin{exemplebox}[Quatre vers, deux prépositions]
\esp{Camino \textbf{por} el mundo} = Je chemine \textbf{à travers} le monde (lieu parcouru, trajet).\par
\esp{Lucho \textbf{por} la paz} = Je lutte \textbf{pour} la paix (cause, motif qui me pousse).\par
\esp{\textbf{Para} que el niño tenga pan} = \textbf{Afin que} l'enfant ait du pain (but, finalité).\par
\esp{Este poema es \textbf{para} ti} = Ce poème est \textbf{pour} toi (destinataire).
\end{exemplebox}

Observons : \esp{por} regarde \emph{en arrière}, vers l'origine, la cause, le chemin parcouru, le moyen ; \esp{para} regarde \emph{en avant}, vers le but, la destination, le destinataire, l'échéance. Retenons cette image simple : \textbf{POR = l'origine, la cause, le trajet ; PARA = le but, la destination, la finalité}.

\begin{methode}[La question magique : ¿por qué? ou ¿para qué?]
Pour choisir entre \esp{por} et \esp{para}, posez-vous l'une de ces deux questions :
\begin{enumerate}
\item \esp{¿POR QUÉ?} (pourquoi ?) : si la réponse donne la \textbf{cause} ou le \textbf{motif}, utilisez \cle{POR}. \esp{Lucho por la justicia.} (Pourquoi ? Parce que la justice est ma cause.)
\item \esp{¿PARA QUÉ?} (dans quel but ?) : si la réponse donne le \textbf{but} ou la \textbf{destination}, utilisez \cle{PARA}. \esp{Estudio para ser médico.} (Dans quel but ? Pour devenir médecin.)
\end{enumerate}
Comparez : \esp{Trabajo por mi familia} (ma famille est la cause de mon effort, je la remplace ou j'agis pour elle) et \esp{Trabajo para una empresa} (l'entreprise est le destinataire de mon travail, mon employeur).
\end{methode}

\courssub{PARA : le but et la destination}

\begin{propriete}[PARA + infinitif = la finalité]
\esp{Para} suivi d'un \textbf{infinitif} exprime toujours la \cle{finalité}, le but : il correspond au français « pour » (+ infinitif) ou « afin de ».\par
\esp{Estudio \textbf{para} aprobar el examen.} = J'étudie \textbf{pour} réussir l'examen.\par
\esp{Vinimos a Yaoundé \textbf{para} aprender español.} = Nous sommes venus à Yaoundé \textbf{pour} apprendre l'espagnol.\par
\emph{Remarque :} en espagnol, l'infinitif est \textbf{obligatoire} après \esp{para} ; on ne peut jamais conjuguer le verbe. On dit \esp{para hablar}, jamais \esp{para hablo}.
\end{propriete}

\begin{propriete}[Les quatre autres emplois de PARA]
\begin{enumerate}
\item \textbf{Le destinataire} (à qui est destinée la chose/action) : \esp{Este regalo es \textbf{para} ti.} = Ce cadeau est \textbf{pour} toi. \esp{Esta carta es \textbf{para} el director.} = Cette lettre est \textbf{pour} le proviseur.
\item \textbf{La destination} (vers quel lieu on se dirige) : \esp{Salgo \textbf{para} Douala.} = Je pars \textbf{pour} Douala. \esp{El autobús \textbf{para} Bata sale a las ocho.} = Le bus \textbf{pour} Bata part à huit heures.
\item \textbf{La date limite} (pour quand, au plus tard) : \esp{El deber es \textbf{para} mañana.} = Le devoir est \textbf{pour} demain. \esp{Lo necesito \textbf{para} el lunes.} = J'en ai besoin \textbf{pour} lundi.
\item \textbf{L'opinion} (selon, du point de vue de) : \esp{\textbf{Para} mí, la paz es lo más importante.} = \textbf{Pour} moi, la paix est ce qu'il y a de plus important. \esp{\textbf{Para} un niño, este libro es difícil.} = \textbf{Pour} un enfant, ce livre est difficile.
\end{enumerate}
\end{propriete}

\courssub{POR : la cause, le chemin et le moyen}

\begin{propriete}[POR exprime la cause ou le motif]
\esp{Por} introduit la \cle{cause}, la raison qui pousse à agir : « à cause de, pour, en raison de, par ».\par
\esp{Luchamos \textbf{por} la paz.} = Nous luttons \textbf{pour} la paix (la paix est la cause qui nous mobilise).\par
\esp{Lo hago \textbf{por} ti.} = Je le fais \textbf{pour} toi (à cause de toi, par amour pour toi).\par
\esp{Gracias \textbf{por} tu ayuda.} = Merci \textbf{pour} ton aide.\par
\esp{No pude venir \textbf{por} la lluvia.} = Je n'ai pas pu venir \textbf{à cause de} la pluie.
\end{propriete}

\begin{propriete}[Les autres emplois de POR]
\begin{enumerate}
\item \textbf{L'agent du passif} (par) : \esp{Este poema fue escrito \textbf{por} un ciudadano del mundo.} = Ce poème a été écrit \textbf{par} un citoyen du monde. \esp{Fue construido \textbf{por} los obreros.} = Il a été construit \textbf{par} les ouvriers.
\item \textbf{L'échange et le prix} (pour, contre) : \esp{Te doy dos \textbf{por} uno.} = Je t'en donne deux \textbf{pour} un. \esp{Compré mangos \textbf{por} mil francos.} = J'ai acheté des mangues \textbf{pour} mille francs. \esp{Cambié mi bicicleta \textbf{por} una moto.} = J'ai échangé mon vélo \textbf{contre} une moto.
\item \textbf{La durée} (pendant, durant) : \esp{Estudié \textbf{por} dos horas.} = J'ai étudié \textbf{pendant} deux heures. \esp{Vivió en España \textbf{por} tres años.} = Il a vécu en Espagne \textbf{pendant} trois ans. \emph{Note :} \esp{dos horas} (sans préposition) est aussi fréquent.
\item \textbf{Le trajet, le lieu approximatif} (par, à travers, dans, par ici) : \esp{Pasamos \textbf{por} la ciudad.} = Nous passons \textbf{par} la ville. \esp{Camino \textbf{por} el mundo.} = Je chemine \textbf{à travers} le monde. \esp{Hay mucha gente \textbf{por} aquí.} = Il y a beaucoup de monde \textbf{par} ici / dans le coin.
\item \textbf{Le moyen, le canal de communication} (par) : \esp{Hablamos \textbf{por} teléfono.} = Nous parlons \textbf{par} téléphone. \esp{Envía el documento \textbf{por} correo.} = Envoie le document \textbf{par} courrier / mail. \esp{Lo supe \textbf{por} la radio.} = Je l'ai su \textbf{par} la radio.
\item \textbf{Le moment de la journée} (le, dans) : \esp{\textbf{Por} la mañana} = le matin ; \esp{\textbf{por} la tarde} = l'après-midi ; \esp{\textbf{por} la noche} = le soir, la nuit.
\end{enumerate}
\end{propriete}

\begin{popfigure}
\begin{tikzpicture}[node distance=0.5cm, every node/.style={font=\small}]
\node[draw=popblue, fill=popblueL, rounded corners, font=\bfseries\large, align=center] (por) {POR\\ \small origine, cause, trajet};
\node[draw=poporange, fill=poporangeL, rounded corners, font=\bfseries\large, align=center, right=8cm of por] (para) {PARA\\ \small but, destination, limite};
% POR branches
\node[draw=popblue, fill=white, rounded corners, align=center, below left=0.7cm and -1cm of por] (p1) {Cause / Motif\\ \esp{por la paz}\\ \esp{gracias por...}};
\node[draw=popblue, fill=white, rounded corners, align=center, below=0.7cm of por] (p2) {Agent passif / Moyen\\ \esp{escrito por Lorca}\\ \esp{por teléfono}};
\node[draw=popblue, fill=white, rounded corners, align=center, below right=0.7cm and -1cm of por] (p3) {Échange / Durée / Trajet\\ \esp{dos por uno}\\ \esp{por dos horas}\\ \esp{por la ciudad}};
\node[draw=popblue, fill=white, rounded corners, align=center, right=0.5cm of p3] (p4) {Moment de la journée\\ \esp{por la mañana}};
% PARA branches
\node[draw=poporange, fill=white, rounded corners, align=center, below left=0.7cm and -1cm of para] (q1) {Finalité (infinitif)\\ \esp{para sembrar}\\ \esp{para que + subj.}};
\node[draw=poporange, fill=white, rounded corners, align=center, below=0.7cm of para] (q2) {Destinataire\\ \esp{para ti}\\ \esp{para el director}};
\node[draw=poporange, fill=white, rounded corners, align=center, below right=0.7cm and -1cm of para] (q3) {Destination / Date limite\\ \esp{salgo para Douala}\\ \esp{para mañana}};
\node[draw=poporange, fill=white, rounded corners, align=center, right=0.5cm of q3] (q4) {Opinion\\ \esp{para mí}};
% Arrows
\draw[-{Stealth}, thick, popblue] (por.south west) -- (p1.north);
\draw[-{Stealth}, thick, popblue] (por.south) -- (p2.north);
\draw[-{Stealth}, thick, popblue] (por.south east) -- (p3.north west);
\draw[-{Stealth}, thick, popblue] (p3.east) -- (p4.west);
\draw[-{Stealth}, thick, poporange] (para.south west) -- (q1.north);
\draw[-{Stealth}, thick, poporange] (para.south) -- (q2.north);
\draw[-{Stealth}, thick, poporange] (para.south east) -- (q3.north west);
\draw[-{Stealth}, thick, poporange] (q3.east) -- (q4.west);
\end{tikzpicture}
\legende{Carte mentale : POR regarde vers l'origine (cause, agent, trajet), PARA vers le but (finalité, destinataire, limite).}
\end{popfigure}

\courssub{Comparaison avec le français et tableau récapitulatif}

Le français « pour » est un piège : il correspond tantôt à \esp{para}, tantôt à \esp{por}. Seule l'analyse du sens (cause ou but ?) permet de choisir.

\begin{center}
\begin{tabularx}{\linewidth}{|X|X|X|}
\hline
\rowcolor{popblueL} Français & Español & Remarque \\
\hline
Je pars \textbf{pour} Yaoundé & Salgo \textbf{para} Yaoundé & Destination \\
\hline
J'étudie \textbf{pour} réussir & Estudio \textbf{para} aprobar & Finalité (para + inf.) \\
\hline
Ce cadeau est \textbf{pour} toi & Este regalo es \textbf{para} ti & Destinataire \\
\hline
Je le fais \textbf{pour} toi (à cause de) & Lo hago \textbf{por} ti & Cause, motif \\
\hline
Merci \textbf{pour} ton aide & Gracias \textbf{por} tu ayuda & Cause de la gratitude \\
\hline
Deux \textbf{pour} un & Dos \textbf{por} uno & Échange \\
\hline
Écrit \textbf{par} Lorca & Escrito \textbf{por} Lorca & Agent du passif \\
\hline
\textbf{Pendant} deux heures & \textbf{Por} dos horas & Durée \\
\hline
Nous passons \textbf{par} la ville & Pasamos \textbf{por} la ciudad & Trajet \\
\hline
\textbf{Le} matin & \textbf{Por} la mañana & Moment de la journée \\
\hline
\end{tabularx}
\end{center}

\begin{popfigure}
\begin{tikzpicture}
\draw[-{Stealth}, very thick, popmuted] (0,0) -- (13,0);
\node[circle, draw=popblue, fill=popblueL, inner sep=2.5pt] at (1.5,0) {};
\node[above=0.2cm, align=center, font=\small] at (1.5,0) {\esp{por la mañana}\\ le matin};
\node[circle, draw=popblue, fill=popblueL, inner sep=2.5pt] at (5.5,0) {};
\node[above=0.2cm, align=center, font=\small] at (5.5,0) {\esp{por dos horas}\\ pendant deux heures};
\node[circle, draw=poporange, fill=poporangeL, inner sep=2.5pt] at (11,0) {};
\node[above=0.2cm, align=center, font=\small] at (11,0) {\esp{para mañana}\\ pour demain (limite)};
\node[below=0.4cm, font=\small\itshape, popmuted] at (6.5,0) {Le temps : POR situe ou mesure (moment, durée), PARA fixe une échéance.};
\end{tikzpicture}
\legende{Frise du temps : POR situe ou mesure, PARA fixe une échéance.}
\end{popfigure}

\courssub{Expressions figées et cas difficiles}

\begin{aretenir}[Expressions figées à mémoriser]
\begin{itemize}
\item \esp{por eso} = c'est pourquoi ; \esp{por fin} = enfin ; \esp{por favor} = s'il vous plaît ; \esp{por lo tanto} = par conséquent ; \esp{por lo menos} = au moins ; \esp{por supuesto} = bien sûr.
\item \esp{para siempre} = pour toujours ; \esp{para qué} = à quoi bon ; \esp{para variar} = pour changer.
\end{itemize}
\esp{Por favor, termina el trabajo para mañana; por lo tanto, no salgas esta noche.} = S'il te plaît, termine le travail pour demain ; par conséquent, ne sors pas ce soir.
\end{aretenir}

\begin{propriete}[Estar por / estar para ; trabajar por / trabajar para]
\begin{itemize}
\item \esp{Estar por + infinitif} = être sur le point de, rester à faire : \esp{El tren está por salir.} = Le train est sur le point de partir. \esp{La tarea está por hacer.} = Le devoir reste à faire.
\item \esp{Estar para + infinitif} = être prêt à, être d'humeur à : \esp{Estoy para salir.} = Je suis prêt à sortir. \esp{No estoy para bromas.} = Je ne suis pas d'humeur à plaisanter.
\item \esp{Trabajar por} = travailler pour (une cause) ou à la place de quelqu'un : \esp{Trabaja por la paz.} = Il travaille pour la paix. \esp{Hoy trabajo por mi colega.} = Aujourd'hui je travaille \textbf{à la place de} mon collègue.
\item \esp{Trabajar para} = travailler pour (un employeur, un destinataire) : \esp{Trabaja para una ONG.} = Il travaille pour une ONG.
\end{itemize}
\end{propriete}

\begin{attention}[Erreurs fréquentes des francophones -- \cle{PIÈGES À ÉVITER}]
\begin{itemize}
\item \textbf{Gracias \textbf{por}...} : on dit \esp{gracias \textbf{por} tu ayuda}, \textbf{jamais} \esp{gracias para}. On remercie \emph{à cause de} quelque chose.
\item \textbf{Durée} : « Pendant deux heures » = \esp{\textbf{por} dos horas} (ou \esp{dos horas} seul), \textbf{jamais} \esp{para dos horas}.
\item \textbf{Date limite vs Moment} : « Pour demain » (échéance) = \esp{\textbf{para} mañana} ; « Demain matin » = \esp{mañana \textbf{por} la mañana}.
\item \textbf{Agent du passif} : exige \esp{por} : \esp{escrito \textbf{por} Lorca}, \textbf{jamais} \esp{escrito de Lorca} ni \esp{para Lorca}.
\item \textbf{Faux ami} : \esp{para} ne signifie \textbf{jamais} « pendant ».
\item \textbf{Subjonctif} : \esp{para que} appelle \textbf{obligatoirement} le subjonctif (\esp{para que el niño \textbf{tenga} pan}), pas l'indicatif.
\end{itemize}
\end{attention}

\begin{savaistu}[Por qué, porque, porqué, por que : quatre écritures, quatre fonctions]
L'espagnol distingue soigneusement :
\begin{enumerate}
\item \esp{¿\textbf{por qué}?} (deux mots, accentué) : « pourquoi ? » (questions directes/indirectes). \esp{¿Por qué luchas?}
\item \esp{\textbf{porque}} (un mot, pas d'accent) : « parce que » (réponse, cause). \esp{Lucho \textbf{porque} creo en la paz.}
\item \esp{el \textbf{porqué}} (un mot, accentué, nom masc.) : « le pourquoi », « la raison ». \esp{No entiendo el \textbf{porqué} de la guerra.}
\item \esp{\textbf{por que}} (deux mots, pas d'accent) : « pour que » (préposition \esp{por} + conjonction \esp{que} ou relatif), suivi du \textbf{subjonctif}. \esp{Lucho \textbf{por que} haya justicia.} (Je lutte pour qu'il y ait la justice). \esp{La razón \textbf{por que} lucho.} (La raison pour laquelle je lutte).
\end{enumerate}
\end{savaistu}

\begin{exoresolu}[Choisissez POR ou PARA et justifiez]
Énoncé. Complétez avec \esp{por} ou \esp{para}, puis traduisez :
\begin{enumerate}
\item \esp{Compré este libro \_\_\_ mi hermana.}
\item \esp{Luchamos \_\_\_ los derechos humanos.}
\item \esp{El tren sale \_\_\_ la mañana.}
\item \esp{Estudio español \_\_\_ viajar a Guinea Ecuatorial.}
\item \esp{Este poema fue escrito \_\_\_ una alumna de Terminale.}
\item \esp{Lo hago \_\_\_ ti, \_\_\_ que sonrías.}
\item \esp{Caminamos \_\_\_ el parque \_\_\_ una hora.}
\end{enumerate}
\tcblower
Solution détaillée :
\begin{enumerate}
\item \esp{Compré este libro \textbf{para} mi hermana.} = J'ai acheté ce livre \textbf{pour} ma sœur. (Destinataire : ¿para quién?)
\item \esp{Luchamos \textbf{por} los derechos humanos.} = Nous luttons \textbf{pour} les droits de l'homme. (Cause, motif : ¿por qué?)
\item \esp{El tren sale \textbf{por} la mañana.} = Le train part \textbf{le} matin. (Moment de la journée : expression figée.)
\item \esp{Estudio español \textbf{para} viajar a Guinea Ecuatorial.} = J'étudie l'espagnol \textbf{pour} voyager en Guinée équatoriale. (Finalité : para + infinitif.)
\item \esp{Este poema fue escrito \textbf{por} una alumna de Terminale.} = Ce poème a été écrit \textbf{par} une élève de Terminale. (Agent du passif.)
\item \esp{Lo hago \textbf{por} ti, \textbf{para} que sonrías.} = Je le fais \textbf{pour} toi (cause), \textbf{afin que} tu souries (but).
\item \esp{Caminamos \textbf{por} el parque \textbf{por} una hora.} = Nous marchons \textbf{par/à travers} le parc \textbf{pendant} une heure. (Trajet + Durée.)
\end{enumerate}
\end{exoresolu}

\courssub{Exercices de thème et de version (Consolidation)}

\begin{exercice}[Thème : Traduisez en espagnol]
\begin{enumerate}
\item Je pars pour Douala demain matin.
\item Merci pour ta lettre, elle est arrivée par courrier.
\item Nous marchons à travers la ville pendant deux heures.
\item Ce cadeau est pour toi, parce que tu es un citoyen du monde.
\item Pour moi, la paix est essentielle ; c'est pourquoi je lutte pour elle.
\item Il a été élu pour quatre ans. (Durée du mandat -> \esp{por}).
\item Ce gâteau est fait par ma mère. (Agent passif).
\item J'ai acheté ce téléphone pour cinq cents francs. (Prix).
\item Nous luttons pour que la justice triomphe. (\esp{para que} + subjonctif).
\item Le devoir est pour lundi.
\end{enumerate}
\end{exercice}

\begin{corrige}[Exercice de thème]
\begin{enumerate}
\item \esp{Salgo \textbf{para} Douala mañana \textbf{por} la mañana.} (Destination + moment.)
\item \esp{Gracias \textbf{por} tu carta, llegó \textbf{por} correo.} (Cause gratitude + moyen.)
\item \esp{Caminamos \textbf{por} la ciudad \textbf{por} dos horas.} (Trajet + durée.)
\item \esp{Este regalo es \textbf{para} ti, \textbf{porque} eres un ciudadano del mundo.} (Destinataire + cause introduite par \esp{porque}.)
\item \esp{\textbf{Para} mí, la paz es esencial; \textbf{por eso} lucho \textbf{por} ella.} (Opinion + expression figée + cause.)
\item \esp{Fue elegido \textbf{por} cuatro años.} (Durée du mandat.)
\item \esp{Este pastel está hecho \textbf{por} mi madre.} (Agent passif + \esp{estar} pour résultat.)
\item \esp{Compré este teléfono \textbf{por} quinientos francos.} (Prix / échange.)
\item \esp{Luchamos \textbf{para que} la justicia \textbf{triunfe}.} (Finalité + subjonctif présent \esp{triunfe}.)
\item \esp{El deber es \textbf{para} el lunes.} (Date limite.)
\end{enumerate}
\end{corrige}

\begin{exercice}[Version : Traduisez en français]
\begin{enumerate}
\item \esp{No estoy para bromas, tengo mucho trabajo por hacer.}
\item \esp{Por lo general, voy a la escuela por bus, pero hoy voy para casa a pie.}
\item \esp{El libro fue escrito por un autor anónimo para niños.}
\item \esp{Por qué no vienes? Es para que me ayudes.}
\item \esp{Trabajo por una ONG que lucha por los derechos de la infancia.}
\end{enumerate}
\end{exercice}

\begin{corrige}[Exercice de version]
\begin{enumerate}
\item Je ne suis pas d'humeur à plaisanter, j'ai beaucoup de travail \textbf{à faire} (\esp{por hacer} = restant à faire).
\item D'habitude (\esp{por lo general}), je vais à l'école \textbf{en} bus, mais aujourd'hui je rentre \textbf{chez moi} \textbf{à pied} (\esp{para casa} = destination, \esp{a pie} = moyen).
\item Le livre a été écrit \textbf{par} un auteur anonyme \textbf{pour} des enfants. (Agent + destinataire).
\item Pourquoi ne viens-tu pas ? C'est \textbf{pour que} tu m'aides. (\esp{Por qué} question, \esp{para que} but + subjonctif).
\item Je travaille \textbf{pour} une ONG qui lutte \textbf{pour} les droits de l'enfance. (Employeur \esp{para} + Cause \esp{por}).
\end{enumerate}
\end{corrige}

\coursec{Débat guidé : « Ciudadanos del mundo »}

\courssub{Préparation : arguments et vocabulaire}

\begin{aretenir}[Expressions pour le débat (Expresión oral)]
\begin{center}
\begin{tabularx}{\linewidth}{|X|X|}
\hline
\rowcolor{popblueL} Función & Expresiones útiles \\
\hline
Dar la opinión & \esp{Para mí... / A mi modo de ver... / Creo que... / Opino que...} \\
\hline
Estar de acuerdo & \esp{Estoy de acuerdo / Tienes razón / Exactamente / Sin duda.} \\
\hline
No estar de acuerdo & \esp{No estoy de acuerdo / No lo veo así / Discrepo / En absoluto.} \\
\hline
Argumentar & \esp{Por un lado... / Por otro lado... / Además... / Por eso... / Porque...} \\
\hline
Pedir opinión & \esp{¿Qué opinas tú? / ¿Tú qué piensas? / ¿Estás de acuerdo?} \\
\hline
Matizar & \esp{Es verdad, pero... / Puede ser, sin embargo... / Depende.} \\
\hline
\end{tabularx}
\end{center}
\end{aretenir}

\begin{experience}[Débat guidé : « ¿Ciudadanos del mundo o ciudadanos de un país? »]
\textbf{Consigne :} En groupes de 4. Chaque groupe tire un rôle (A, B, C, D). Préparez 3 arguments en 5 min. Débattez 10 min. Un secrétaire note les idées clés pour une synthèse écrite finale.

\textbf{Rôles :}
\begin{itemize}
\item \textbf{A : Le cosmopolite convaincu.} « Les frontières sont artificielles. L'humanité est une seule famille. Je me sens citoyen du monde avant d'être camerounais/espagnol. »
\item \textbf{B : Le patriote attaché aux racines.} « On ne peut aimer l'humanité en général si on n'aime pas sa patrie concrète. La citoyenneté implique des devoirs envers un État (vote, impôts, défense). »
\item \textbf{C : Le pragmatique (juriste/politologue).} » « Le "citoyen du monde" n'a pas de statut juridique. Pas de passeport mondial, pas de tribunal mondial pour tous. C'est une belle métaphore, pas une réalité légale. »
\item \textbf{D : Le militant des droits humains.} « La citoyenneté mondiale, c'est la responsabilité solidaire : climat, migrants, pandémies. \esp{Por eso} agissons \textbf{para} un droit international contraignant. »
\end{itemize}

\textbf{Questions déclencheurs :}
\begin{enumerate}
\item \esp{¿Tiene sentido hablar de "ciudadanía mundial" sin un Estado mundial?}
\item \esp{¿Es compatible amar a tu país y ser "ciudadano del mundo"?}
\item \esp{¿Por qué luchamos por los derechos humanos: por identidad nacional o por dignidad universal?}
\end{enumerate}
\end{experience}

\begin{methode}[Synthèse écrite post-débat (Modèle de rédaction)]
Rédigez un paragraphe argumenté (80-100 mots) en espagnol :
\textbf{Structure :}
\begin{enumerate}
\item \textbf{These} : \esp{En mi opinión, la ciudadanía mundial es... / no es...}
\item \textbf{Arguments} : \esp{Por un lado... Por otro lado... Además...}
\item \textbf{Exemple personnel/concret} : \esp{Por ejemplo, en Camerún / en Guinea Ecuatorial...}
\item \textbf{Conclusion} : \esp{Por lo tanto / En definitiva, ser ciudadano del mundo implica...}
\end{enumerate}
\textbf{Contraintes :} Utiliser au moins 3 fois \esp{por} (cause, durée, agent...) et 3 fois \esp{para} (but, destinataire, limite...). Soulignez-les.
\end{methode}

\begin{exemplebox}[Modèle de production écrite (Corrigé type)]
\esp{En mi opinión, la ciudadanía mundial \textbf{no es} un documento legal, \textbf{sino} un compromiso moral. \textbf{Por} un lado, tenemos deberes \textbf{para con} nuestro país (votar, pagar impuestos). \textbf{Por} otro, los retos globales (clima, migración) nos obligan a actuar \textbf{para} el bien común de la humanidad. \textbf{Por ejemplo}, en Camerún, los jóvenes luchan \textbf{por} la justicia \textbf{para} que sus voces cuenten \textbf{para} siempre. \textbf{Por lo tanto}, ser ciudadano del mundo es actuar \textbf{por} la dignidad \textbf{para} todos.}
\end{exemplebox}

\courssub{Évaluation formative : Auto-évaluation}

\begin{exercice}[Testez vos acquis]
Cochez la case correspondante. Si vous hésitez, relisez la règle correspondante.
\begin{center}
\begin{tabularx}{\linewidth}{|X|c|c|c|}
\hline
\rowcolor{popblueL} Compétence & \textbf{Acquis} & \textbf{En cours} & \textbf{À revoir} \\
\hline
Je distingue \esp{por} (cause) et \esp{para} (but) dans une phrase simple. & & & \\
\hline
J'utilise correctement \esp{para + infinitif} et \esp{para que + subjonctif}. & & & \\
\hline
Je sais traduire « pendant 2 heures », « le matin », « par téléphone », « écrit par... ». & & & \\
\hline
Je participe au débat en utilisant \esp{por/para} et les connecteurs logiques. & & & \\
\hline
Je traduis un vers poétique en respectant le sens et la musicalité. & & & \\
\hline
\end{tabularx}
\end{center}
\end{exercice}

\begin{aretenir}[L'essentiel de la leçon E12 -- \cle{RÉVISION BAC}]
\begin{itemize}
\item \cle{Traduction poétique} : sens littéral $\rightarrow$ registre français $\rightarrow$ musicalité. Attention au \esp{para que + subjonctif} (finalité).
\item \cle{PARA} = \textbf{AVANT} (But, Destinataire, Destination, Date limite, Opinion). \esp{Para + inf. / Para que + subj. / Para + nom.}
\item \cle{POR} = \textbf{ARRIÈRE / AUTOUR} (Cause, Agent passif, Échange/Prix, Durée, Trajet/Lieu approx., Moyen, Moment journée).
\item \cle{Question magique} : \esp{¿Por qué?} (Cause $\rightarrow$ \esp{por}) vs \esp{¿Para qué?} (But $\rightarrow$ \esp{para}).
\item \cle{Pièges} : \esp{gracias por}, \esp{por dos horas} (durée), \esp{para mañana} (limite), \esp{escrito por} (agent), \esp{por que} (subjonctif).
\item \cle{Débat} : Structurer (These/Arguments/Exemple/Conclusion), vocabulaire de l'opinion (\esp{para mí, opino que}), connecteurs (\esp{por tanto, por un lado}).
\end{itemize}
\end{aretenir}

\coursec{Exercices de thème et de version}

Dans la section précédente, nous avons établi la règle générale qui gouverne l'emploi de \esp{por} et \esp{para} : \esp{por} regarde en arrière (la \cle{cause}, l'origine, l'agent, le moyen, le trajet), tandis que \esp{para} regarde en avant (la \cle{finalité}, le but, le destinataire, la destination). Il est temps de vérifier cette compétence par la pratique, car c'est exactement ce type d'exercice — thème, version et justification grammaticale — qui vous attend au baccalauréat. Nous progresserons en quatre étapes : d'abord la version (de l'espagnol vers le français), puis le thème (du français vers l'espagnol), ensuite un exercice à trous avec justification, et enfin la correction d'erreurs typiques et une mini-traduction poétique.

\courssub{Version : quatre phrases en contexte citoyen}

\begin{methode}[Méthode de version avec por/para]
Pour traduire correctement une phrase contenant \esp{por} ou \esp{para} :
\begin{enumerate}
\item Repérer la préposition et son complément.
\item Identifier sa \textbf{fonction} : cause, finalité, agent, lieu de passage, destination, destinataire, échange.
\item Choisir l'équivalent français naturel : « pour », « par », « afin de », « à cause de », « en échange de », etc.
\item Traduire la phrase entière en français correct, sans calquer mot à mot.
\end{enumerate}
Au baccalauréat, la justification de la fonction est ce qui rapporte les points de grammaire : ne vous contentez jamais d'une traduction globale.
\end{methode}

\begin{exoresolu}[Version guidée]
\textbf{Énoncé.} Traduisez en français et justifiez l'emploi de \esp{por} ou \esp{para} :
\begin{enumerate}
\item \esp{Los jóvenes marcharon por la libertad de expresión.}
\item \esp{Esta ley fue aprobada para proteger a los niños.}
\item \esp{El discurso fue pronunciado por la presidenta de la asociación.}
\item \esp{Mañana salgo para Malabo a participar en un encuentro de estudiantes.}
\end{enumerate}
\tcblower
\textbf{Solution détaillée.}
\begin{enumerate}
\item \esp{por la libertad de expresión} : \textbf{cause} (la liberté est le motif de la marche). Traduction : « Les jeunes ont manifesté pour la liberté d'expression. »
\item \esp{para proteger} : \textbf{finalité} (\esp{para} + infinitif = « afin de »). Traduction : « Cette loi a été adoptée pour protéger les enfants. »
\item \esp{por la presidenta} : \textbf{agent} du passif (\esp{fue pronunciado}). Traduction : « Le discours a été prononcé par la présidente de l'association. »
\item \esp{para Malabo} : \textbf{destination} (mouvement vers un lieu). Traduction : « Demain je pars pour Malabo afin de participer à une rencontre d'étudiants. »
\end{enumerate}
\end{exoresolu}

\begin{exemplebox}[Por et para dans la même phrase]
\esp{El pueblo lucha por sus derechos para construir un mundo mejor.}

« Le peuple lutte pour ses droits afin de construire un monde meilleur. »

Observez : \esp{por sus derechos} exprime la \textbf{cause} (on lutte à cause de, au nom de ses droits) ; \esp{para construir} exprime le \textbf{but} (l'objectif visé). Le français utilise « pour » dans les deux cas, mais l'espagnol distingue rigoureusement les deux fonctions. C'est toute la difficulté — et toute la beauté — de ce point de grammaire.
\end{exemplebox}

\courssub{Thème : cinq phrases françaises à traduire}

Le thème est plus exigeant que la version : c'est vous qui devez choisir entre \esp{por} et \esp{para}. Avant d'écrire, posez-vous la question rituelle : le « pour » français regarde-t-il en arrière (cause, motif, soutien) ou en avant (but, destinataire, destination) ?

\begin{exoresolu}[Thème guidé]
\textbf{Énoncé.} Traduisez en espagnol :
\begin{enumerate}
\item Je me bats pour la justice.
\item Ce livre est pour toi.
\item Il passe par le marché.
\item Il a été élu par le peuple.
\item Je pars pour Bafoussam demain.
\end{enumerate}
\tcblower
\textbf{Solution détaillée.}
\begin{enumerate}
\item « Pour la justice » = cause, motif du combat $\rightarrow$ \esp{por}. \esp{Yo lucho por la justicia.} (ou \esp{Yo me bato por la justicia} : \esp{batirse} existe au sens littéraire de « se battre », mais \esp{luchar por} est l'expression la plus naturelle.)
\item « Pour toi » = destinataire $\rightarrow$ \esp{para}. \esp{Este libro es para ti.} (Après une préposition, « toi » se dit \esp{ti}, jamais \esp{tú}.)
\item « Par le marché » = lieu de passage, trajet $\rightarrow$ \esp{por}. \esp{Pasa por el mercado.}
\item « Par le peuple » = agent du passif $\rightarrow$ \esp{por}. \esp{Ha sido elegido por el pueblo.} (ou \esp{Fue elegido por el pueblo.})
\item « Pour Bafoussam » = destination $\rightarrow$ \esp{para}. \esp{Mañana salgo para Bafoussam.} (ou \esp{Me marcho para Bafoussam mañana.})
\end{enumerate}
\end{exoresolu}

\begin{attention}[Deux « pour » qui ne se traduisent pas par « para »]
\begin{itemize}
\item « Être reconnaissant pour » se dit \esp{estar agradecido por} : la gratitude a une \textbf{cause}. \esp{Estoy agradecido por tu ayuda.} « Je te suis reconnaissant pour ton aide. » Ne dites jamais \esp{agradecido para}.
\item « Voter pour quelqu'un » se dit \esp{votar por alguien} : on vote \textbf{en faveur de}, par soutien (cause). \esp{Los ciudadanos votaron por la candidata joven.} « Les citoyens ont voté pour la jeune candidate. »
\item Mais attention : « voter pour élire » (but de l'action de voter) exige \esp{para} : \esp{Votamos para elegir a nuestros representantes.} « Nous votons pour élire nos représentants. » Même verbe, deux fonctions, deux prépositions.
\end{itemize}
\end{attention}

\begin{center}
\begin{tabularx}{\linewidth}{|X|X|X|}
\hline
\rowcolor{popblueL} Français & Español & Fonction \\
\hline
Je me bats pour la justice. & \esp{Yo lucho por la justicia.} & cause, motif \\
\hline
Ce livre est pour toi. & \esp{Este libro es para ti.} & destinataire \\
\hline
Il passe par le marché. & \esp{Pasa por el mercado.} & lieu de passage \\
\hline
Il a été élu par le peuple. & \esp{Ha sido elegido por el pueblo.} & agent du passif \\
\hline
Je pars pour Bafoussam demain. & \esp{Mañana salgo para Bafoussam.} & destination \\
\hline
Merci pour ton aide. & \esp{Gracias por tu ayuda.} & cause (remerciement) \\
\hline
Il étudie pour devenir avocat. & \esp{Estudia para ser abogado.} & finalité \\
\hline
\end{tabularx}
\end{center}

\begin{popfigure}
\begin{center}
\begin{tikzpicture}[
  box/.style={draw=popblue, fill=popblueL, rounded corners=3pt, align=center, font=\small, inner sep=4pt},
  just/.style={draw=poporange, fill=poporangeL, rounded corners=3pt, align=center, font=\footnotesize, inner sep=3pt},
  arr/.style={-{Stealth[length=2.5mm]}, thick, popdark}
]
\node[box, align=center] (p1) at (0,3.2){Je me bats\\ pour la justice};
\node[box, fill=poppinkL, draw=poppink, align=center] (p2) at (6.2,3.2){Ce livre\\ est pour toi};
\node[box, align=center] (p3) at (0,1.6){Il a été élu\\ par le peuple};
\node[box, fill=poppinkL, draw=poppink, align=center] (p4) at (6.2,1.6){Je pars\\ pour Bafoussam};
\node[box, align=center] (p5) at (0,0){Il passe\\ par le marché};
\node[box, fill=poppinkL, draw=poppink, align=center] (p6) at (6.2,0){J'étudie pour\\ devenir médecin};

\node[just, align=center] (j1) at (3.1,3.2){\esp{por}\\ cause};
\node[just, fill=popgreenL, draw=popgreen, align=center] (j2) at (9.3,3.2){\esp{para}\\ destinataire};
\node[just, align=center] (j3) at (3.1,1.6){\esp{por}\\ agent};
\node[just, fill=popgreenL, draw=popgreen, align=center] (j4) at (9.3,1.6){\esp{para}\\ destination};
\node[just, align=center] (j5) at (3.1,0){\esp{por}\\ passage};
\node[just, fill=popgreenL, draw=popgreen, align=center] (j6) at (9.3,0){\esp{para}\\ finalité};

\draw[arr] (p1) -- (j1);
\draw[arr] (p2) -- (j2);
\draw[arr] (p3) -- (j3);
\draw[arr] (p4) -- (j4);
\draw[arr] (p5) -- (j5);
\draw[arr] (p6) -- (j6);
\end{tikzpicture}
\end{center}
\legende{Tableau de correction : pour chaque phrase française, le choix por/para et sa justification par la fonction grammaticale.}
\end{popfigure}

\courssub{Exercice à trous : compléter puis justifier}

\begin{exercice}[Por ou para ?]
Complétez avec \esp{por} ou \esp{para}, puis justifiez votre choix par la fonction (cause, finalité, agent, passage, destination, destinataire, durée, échange).
\begin{enumerate}
\item \esp{Los ciudadanos luchan \ldots\ sus derechos.}
\item \esp{Compré un regalo \ldots\ mi hermana.}
\item \esp{El tren pasa \ldots\ muchas ciudades antes de llegar.}
\item \esp{El poema fue escrito \ldots\ un poeta de Guinea Ecuatorial.}
\item \esp{Trabajamos mucho \ldots\ construir un mundo más justo.}
\item \esp{Estuve en Douala \ldots\ tres semanas.}
\item \esp{Cambié mi bicicleta \ldots\ un balón de fútbol.}
\item \esp{Salimos temprano \ldots\ Yaoundé mañana por la mañana.}
\end{enumerate}
\end{exercice}

\begin{corrige}[Exercice : por ou para ?]
\begin{enumerate}
\item \esp{por} : cause, motif de la lutte (les droits sont la raison du combat).
\item \esp{para} : destinataire du cadeau.
\item \esp{por} : lieu de passage, trajet.
\item \esp{por} : agent du passif (\esp{fue escrito}).
\item \esp{para} : finalité (\esp{para} + infinitif = « afin de »).
\item \esp{por} : durée (« pendant trois semaines »).
\item \esp{por} : échange (« en échange de »).
\item \esp{para} : destination (mouvement vers Yaoundé).
\end{enumerate}
\end{corrige}

\courssub{Phrase à corriger : l'erreur typique du francophone}

Le francophone raisonne avec un seul mot, « pour », et le projette partout en espagnol. Résultat : il dit \esp{para} là où l'espagnol exige \esp{por}. Voici l'exercice de détection qui vous immunisera.

\begin{exoresolu}[Repérez et corrigez l'erreur]
\textbf{Énoncé.} Chaque phrase contient une erreur de préposition. Repérez-la, corrigez-la et justifiez.
\begin{enumerate}
\item \esp{Gracias para tu visita.}
\item \esp{El puente fue construido para los ingenieros.}
\item \esp{Votamos para este candidato porque defiende a los jóvenes.}
\item \esp{Pasamos para el mercado de camino a la escuela.}
\end{enumerate}
\tcblower
\textbf{Solution détaillée.}
\begin{enumerate}
\item Erreur : \esp{para} $\rightarrow$ correction : \esp{Gracias \textbf{por} tu visita.} Le remerciement exprime une \textbf{cause} ; en espagnol, « merci pour » = \esp{gracias por}. C'est l'erreur la plus fréquente des francophones.
\item Erreur : \esp{para} $\rightarrow$ correction : \esp{El puente fue construido \textbf{por} los ingenieros.} « Le pont a été construit par les ingénieurs » : agent du passif, toujours introduit par \esp{por}.
\item Erreur : \esp{para} $\rightarrow$ correction : \esp{Votamos \textbf{por} este candidato.} « Voter pour quelqu'un » = \esp{votar por alguien} (soutien, cause). Réservez \esp{para} au but : \esp{votamos para elegir}.
\item Erreur : \esp{para} $\rightarrow$ correction : \esp{Pasamos \textbf{por} el mercado.} « Passer par » un lieu (trajet, passage) se dit \esp{pasar por}. \esp{Para} marquerait la destination finale, ce qui changerait le sens.
\end{enumerate}
\end{exoresolu}

\begin{attention}[Le réflexe anti-calque]
Avant d'écrire \esp{para}, demandez-vous toujours : « Puis-je remplacer le "pour" français par "afin de", "à destination de" ou "destiné à" ? »
\begin{itemize}
\item Si oui $\rightarrow$ \esp{para} (finalité, destination, destinataire).
\item Si non — si « pour » signifie « à cause de », « en faveur de », « pendant », « en échange de » — $\rightarrow$ \esp{por}.
\end{itemize}
Exemple : « Il pleut, je reste pour la pluie » n'a pas de sens en français avec « afin de » ; on dirait « à cause de la pluie » $\rightarrow$ \esp{Me quedo en casa por la lluvia.}
\end{attention}

\courssub{Mini-traduction poétique : à vous de jouer}

Pour clore cette séance d'entraînement, voici deux vers originaux, dans l'esprit du poème engagé étudié en début de leçon. Traduisez-les en français en respectant le sens, le rythme et la fonction de chaque préposition.

\begin{textebox}[Dos versos para traducir]
\esp{Caminamos por la noche buscando la aurora,}\\\esp{sembramos palabras para que florezca la libertad.}
\end{textebox}

Glossaire : \esp{caminar} = marcher, cheminer ; \esp{la aurora} = l'aurore, l'aube ; \esp{sembrar} = semer ; \esp{florecer} = fleurir, s'épanouir.

\begin{exercice}[Traduction poétique en autonomie]
\begin{enumerate}
\item Traduisez les deux vers en français soigné.
\item Justifiez l'emploi de \esp{por} dans le premier vers et de \esp{para} dans le second.
\item Question d'ouverture : quel effet produit, selon vous, le contraste entre \esp{la noche} et \esp{la aurora} ?
\end{enumerate}
\end{exercice}

\begin{corrige}[Mini-traduction poétique]
\begin{enumerate}
\item Traduction possible : « Nous cheminons dans la nuit en cherchant l'aurore, / nous semons des paroles pour que fleurisse la liberté. » Toute traduction respectant le sens et un ton soutenu est acceptable : « Nous marchons à travers la nuit à la recherche de l'aube... » convient aussi.
\item \esp{por la noche} : \textbf{lieu/temps de passage} — on chemine \emph{à travers} la nuit (comme \esp{por el mercado}). \esp{para que florezca} : \textbf{finalité} — le but de semer des paroles ; notez que \esp{para que} + subjonctif introduit un but avec un sujet différent (« pour que la liberté fleurisse »).
\item Le contraste nuit/aurore symbolise le passage de l'oppression à l'espérance : la marche (\esp{por la noche}) est l'épreuve présente, l'aurore est l'objectif (\esp{para}) d'un peuple qui lutte. Les deux prépositions dessinent ainsi le mouvement même du poème : partir d'une cause douloureuse vers un but lumineux.
\end{enumerate}
\end{corrige}

\begin{aretenir}[L'essentiel de la section]
\begin{itemize}
\item \esp{por} regarde en arrière : cause, motif, agent du passif, moyen, passage, durée, échange, soutien (\esp{luchar por}, \esp{votar por}, \esp{gracias por}, \esp{estar agradecido por}).
\item \esp{para} regarde en avant : finalité (\esp{para} + infinitif, \esp{para que} + subjonctif), destination, destinataire, date limite.
\item Le « pour » français est ambigu : il faut toujours identifier la \textbf{fonction} avant de choisir.
\item Au baccalauréat, en thème comme en version, \textbf{justifier le choix par la fonction grammaticale} est la compétence évaluée : c'est elle qui rapporte les points.
\end{itemize}
\end{aretenir}

\coursec{Débat guidé : « Ciudadanos del mundo »}

Après les exercices de thème et de version, qui ont consolidé l'emploi de \esp{por} et \esp{para} dans des phrases isolées, il est temps de mettre ces outils en action dans une situation de communication authentique : le \cle{débat guidé}. Traduire, c'est choisir ; débattre, c'est choisir aussi — mais devant un auditoire, dans le temps limité, avec des arguments. Cette séance finale du module 3 mobilise tout ce qui a été appris : le lexique des droits et devoirs, les connecteurs d'argumentation, et la grammaire de \esp{por} y \esp{para}, désormais obligatoire dans chaque prise de parole.

\courssub{Présentation du sujet}

Le manuel \emph{Didáctica V} (unité 3) propose la question suivante, que nous écrivons au tableau et que chaque élève recopie :

\begin{center}
\textbf{\esp{¿Somos ciudadanos del mundo o sólo de nuestro país?}}
\end{center}

« Sommes-nous citoyens du monde ou seulement de notre pays ? »

Cette question est d'une grande actualité : un jeune de Yaoundé suit en direct sur son téléphone un match du Real Madrid, échange avec des amis de Malabo ou de Mexico, et se sent concerné par le climat, les migrations ou les conflits lointains. En même temps, il vote (ou votera) au Cameroun, paie des impôts au Cameroun, et c'est le Cameroun qui lui délivre sa carte d'identité. Où commence et où s'arrête sa citoyenneté ?

\begin{definition}[Ciudadanía]
\esp{La ciudadanía} (la citoyenneté) désigne l'appartenance d'une personne à une communauté politique, avec des \cle{droits} (\esp{derechos}) et des \cle{devoirs} (\esp{deberes}). \esp{Ciudadano del mundo} (citoyen du monde) traduit l'idée que certains droits et devoirs dépassent les frontières nationales. Attention au genre : \esp{el ciudadano / la ciudadana}.
\end{definition}

\courssub{Phrases d'appui : observer avant de débattre}

Avant de lancer le débat, lisons à voix haute les deux phrases d'appui ci-dessous. Elles montrent comment une opinion s'appuie sur un argument, et comment \esp{por} et \esp{para} s'y glissent naturellement.

\begin{textebox}[Frases de apoyo para el debate]
— Estoy a favor porque los derechos humanos son para todos, no sólo para los que nacieron en un país rico. Lucho por la justicia, aquí y en todas partes.\par
— No estoy de acuerdo: primero tenemos deberes por nuestro país. Para construir una nación fuerte, hay que trabajar por ella antes de pensar en el mundo entero.\par
— Sin embargo, en mi opinión, cuidar el planeta es un deber para con todos, porque el clima no conoce fronteras.\par
— Ya que vivimos en la aldea global, creo que debemos aprender idiomas para entender a los demás. Por eso estudio español.
\end{textebox}

\textbf{Glossaire :} \esp{luchar por} = lutter pour ; \esp{un deber para con todos} = un devoir envers tous ; \esp{la aldea global} = le village planétaire ; \esp{las fronteras} = les frontières.

\begin{attention}[Observer avant de généraliser]
Dans la première phrase, \esp{para todos} exprime le \cle{destinataire} (les droits sont destinés à tous) et \esp{lucho por la justicia} exprime la \cle{cause} (je lutte en faveur de, à cause de). Dans la deuxième, \esp{deberes por nuestro país} = devoirs envers (pour) notre pays, et \esp{Para construir...} = \esp{para} + infinitif de \cle{but} : « pour construire ». Retenez : le français « pour » seul ne suffit pas à choisir ; il faut se demander : cause/motif (\esp{por}) ou but/destinataire (\esp{para}) ?
\end{attention}

\courssub{Le lexique des droits et devoirs}

Le débat exige un vocabulaire précis. Voici le tableau récapitulatif à recopier et à réemployer à l'oral.

\begin{center}
\begin{tabularx}{\linewidth}{|X|X|X|}\hline
\rowcolor{popblueL} \textbf{Español} & \textbf{Français} & \textbf{Remarque} \\ \hline
\esp{los derechos humanos} & les droits de l'homme & toujours au pluriel dans ce sens \\ \hline
\esp{el derecho a la educación} & le droit à l'éducation & \esp{derecho a} + nom \\ \hline
\esp{la libertad de expresión} & la liberté d'expression & \esp{libertad} est féminin \\ \hline
\esp{los deberes del ciudadano} & les devoirs du citoyen & \esp{deber} $\neq$ \esp{deber de} (devoir probable) \\ \hline
\esp{votar en las elecciones} & voter aux élections & verbe régulier en \esp{-ar} \\ \hline
\esp{respetar las leyes} & respecter les lois & \esp{la ley}, pluriel \esp{las leyes} \\ \hline
\esp{la identidad nacional} & l'identité nationale & faux ami partiel : \esp{nacional} = national \\ \hline
\esp{la solidaridad} & la solidarité & \esp{ser solidario con} = être solidaire de \\ \hline
\esp{las migraciones} & les migrations & \esp{emigrar} (quitter) $\neq$ \esp{inmigrar} (arriver) \\ \hline
\esp{la globalización / mundialización} & la mondialisation & deux termes corrects \\ \hline
\esp{la frontera} & la frontière & \esp{cruzar la frontera} \\ \hline
\esp{la paz mundial} & la paix mondiale & \esp{la paz}, s'utilise généralement au singulier \\ \hline
\end{tabularx}
\end{center}

\begin{savaistu}[Un précurseur espagnol de la citoyenneté universelle]
Le philosophe et théologien espagnol \cle{Francisco de Vitoria} (XVI\textsuperscript{e} siècle), professeur à l'université de Salamanque, est considéré comme l'un des pères du \cle{droit international}. Il défendit l'idée que tous les peuples, y compris les Indiens d'Amérique, possèdent des droits par nature — une intuition qui annonce la notion moderne de citoyenneté universelle. Quatre siècles plus tard, la Déclaration universelle des droits de l'homme (1948) reprend cette idée : les droits ne s'arrêtent pas aux frontières.
\end{savaistu}

\courssub{Les outils linguistiques du débat}

Un débat n'est pas une dispute : c'est un échange réglé où chacun doit \cle{opiner}, \cle{argumenter} et \cle{réagir} avec des formules précises. Voici la boîte à outils.

\begin{propriete}[Opinar, argumentar, reaccionar]
\textbf{1. Opiner (donner son avis) :} \esp{Creo que...} (je crois que), \esp{Pienso que...}, \esp{En mi opinión...} (à mon avis), \esp{Me parece que...}, \esp{Para mí...}. Après ces expressions en forme affirmative, on utilise l'\textbf{indicatif} : \esp{Creo que somos ciudadanos del mundo.} En revanche, à la forme négative, le subjonctif apparaît : \esp{No creo que seamos sólo ciudadanos de nuestro país.}\par
\textbf{2. Argumentar (justifier) :} \esp{porque} (parce que), \esp{ya que} / \esp{puesto que} (puisque), \esp{por eso} (c'est pourquoi), \esp{además} (de plus), \esp{en primer lugar} (d'abord).\par
\textbf{3. Reaccionar (réagir) :} \esp{Estoy de acuerdo (con...)} (je suis d'accord avec), \esp{No estoy de acuerdo} (je ne suis pas d'accord), \esp{Sin embargo} (cependant), \esp{Tienes razón, pero...} (tu as raison, mais), \esp{Al contrario...} (au contraire).
\end{propriete}

\begin{attention}[« Être d'accord avec » = \esp{estar de acuerdo con}]
Les francophones oublient souvent la préposition \esp{de} et disent \esp{*estar acuerdo} — c'est une erreur grave. La formule complète est \esp{estar de acuerdo con alguien/algo} : \esp{Estoy de acuerdo con María.} « Je suis d'accord avec Marie. » Autre piège : ne dites pas \esp{*yo pienso sí} pour « je pense que oui » ; dites \esp{Creo que sí} / \esp{Pienso que sí}. Et « je ne suis pas d'accord » se dit \esp{No estoy de acuerdo}, jamais \esp{*no soy acuerdo} : ici l'espagnol utilise \esp{estar} (état), pas \esp{ser}.
\end{attention}

\begin{methode}[Structurer son intervention en trois temps]
Pour prendre la parole efficacement, chaque élève suit le schéma \textbf{opinion → argument → exemple} :
\begin{enumerate}
\item \textbf{Opinion} : annoncez clairement votre position avec une formule d'opinion. \esp{En mi opinión, somos ciudadanos del mundo.}
\item \textbf{Argument} : justifiez avec \esp{porque}, \esp{ya que} ou \esp{por eso}, en glissant au moins un \esp{por} ou \esp{para} correct. \esp{...porque los derechos humanos son para todos.}
\item \textbf{Exemple concret} : illustrez avec un fait du Cameroun ou du monde hispanophone. \esp{Por ejemplo, los estudiantes de Yaoundé aprenden español para comunicarse con Guinea Ecuatorial.}
\end{enumerate}
Une intervention complète tient en trois ou quatre phrases. Mieux vaut une idée claire et bien appuyée qu'un discours long et confus.
\end{methode}

\begin{popfigure}
\begin{tikzpicture}[
  node distance=0.55cm and 0.9cm,
  caja/.style={draw=popblue, thick, rounded corners=3pt, fill=popblueL, align=center, inner sep=5pt, font=\small},
  caja2/.style={draw=poporange, thick, rounded corners=3pt, fill=poporangeL, align=center, inner sep=5pt, font=\small},
  caja3/.style={draw=popgreen, thick, rounded corners=3pt, fill=popgreenL, align=center, inner sep=5pt, font=\small},
  flecha/.style={-{Stealth[length=2.5mm]}, thick, popdark}
]
\node[caja, fill=poppurpleL, draw=poppurple, align=center] (tema){\textbf{TEMA}\\ ¿Ciudadanos del mundo?};
\node[caja, below left=0.7cm and 1.6cm of tema, align=center] (fav){\textbf{A FAVOR}\\ Equipo 1};
\node[caja, below right=0.7cm and 1.6cm of tema, align=center] (con){\textbf{EN CONTRA}\\ Equipo 2};
\node[caja2, below=0.7cm of fav, align=center] (op1){OPINIÓN\\ Creo que...};
\node[caja2, below=0.7cm of con, align=center] (op2){OPINIÓN\\ No estoy de acuerdo...};
\node[caja2, below=0.6cm of op1, align=center] (arg1){ARGUMENTO\\ porque + por/para};
\node[caja2, below=0.6cm of op2, align=center] (arg2){ARGUMENTO\\ ya que + por/para};
\node[caja2, below=0.6cm of arg1, align=center] (rep1){RÉPLICA\\ Sin embargo...};
\node[caja2, below=0.6cm of arg2, align=center] (rep2){RÉPLICA\\ Al contrario...};
\node[caja3, below=1.5cm of $(rep1)!0.5!(rep2)$, align=center] (fin){\textbf{CONCLUSIÓN}\\ síntesis + voto de la clase};
\draw[flecha] (tema) -- (fav);
\draw[flecha] (tema) -- (con);
\draw[flecha] (fav) -- (op1);
\draw[flecha] (op1) -- (arg1);
\draw[flecha] (arg1) -- (rep1);
\draw[flecha] (con) -- (op2);
\draw[flecha] (op2) -- (arg2);
\draw[flecha] (arg2) -- (rep2);
\draw[flecha] (rep1) -- (fin);
\draw[flecha] (rep2) -- (fin);
\end{tikzpicture}
\legende{Organigramme du débat guidé : du sujet à la conclusion collective.}
\end{popfigure}

\courssub{Organisation matérielle du débat}

\begin{propriete}[Règles du jeu]
\begin{itemize}
\item La classe est divisée en \textbf{deux équipes} : \esp{equipo a favor} (pour la citoyenneté mondiale) et \esp{equipo en contra} (pour la primauté de la citoyenneté nationale).
\item Un \textbf{modérateur} (\esp{el moderador / la moderadora}) distribue la parole, veille au respect du temps et rappelle les règles. C'est un rôle exigeant, à confier à un élève à l'aise.
\item \textbf{Temps de parole réglementé} : chaque intervention dure au maximum \textbf{une minute}. Le modérateur annonce : \esp{Tiempo, por favor.} ou \esp{Te quedan diez segundos.}
\item \textbf{Contrainte grammaticale obligatoire} : chaque argument doit contenir au moins un \esp{por} ou un \esp{para} correctement employé. Si la préposition est fautive, le modérateur peut demander une reformulation : \esp{¿Puedes repetir con «por» o «para» correcto?}
\item On s'exprime \textbf{uniquement en espagnol} ; le français est réservé aux éclaircissements demandés au professeur.
\end{itemize}
\end{propriete}

\courssub{Arguments attendus des deux équipes}

Le professeur veille à ce que les idées suivantes émergent, sans les imposer.

\textbf{Équipe « a favor » (citoyens du monde) :}
\begin{itemize}
\item \esp{La globalización nos une: vemos las mismas noticias y los mismos partidos por internet.} (La mondialisation nous unit : nous voyons les mêmes nouvelles et les mêmes matchs sur internet.)
\item \esp{Los derechos humanos son universales: valen para todos los seres humanos.} (Les droits humains sont universels : ils valent pour tous les êtres humains.)
\item \esp{Las migraciones crean familias entre países: muchos cameruneses trabajan por sus familias desde Europa o América.} (Les migrations créent des familles entre pays.)
\item \esp{El clima no tiene fronteras: luchamos por el planeta entero.} (Le climat n'a pas de frontières.)
\end{itemize}

\textbf{Équipe « en contra » (d'abord citoyens de son pays) :}
\begin{itemize}
\item \esp{La identidad nacional nos da raíces: primero somos cameruneses.} (L'identité nationale nous donne des racines.)
\item \esp{Tenemos deberes por nuestro país: pagar impuestos, votar, respetar las leyes.} (Nous avons des devoirs envers notre pays.)
\item \esp{Para cambiar el mundo, primero hay que construir una nación fuerte.} (Pour changer le monde, il faut d'abord construire une nation forte.)
\item \esp{El pasaporte que llevamos es de un país, no del mundo.} (Le passeport que nous portons est celui d'un pays, pas du monde.)
\end{itemize}

\begin{exoresolu}[Transformer une idée en argument complet]
\textbf{Énoncé.} Un élève dit simplement : \esp{Los derechos humanos son importantes.} Transformez cette idée en argument complet (opinion + justification + exemple), avec au moins un \esp{por} ou un \esp{para}.
\tcblower
\textbf{Solution détaillée.} On applique le schéma en trois temps :
\begin{enumerate}
\item Opinion : \esp{En mi opinión, los derechos humanos son la base de la ciudadanía.}
\item Argument avec \esp{para} (destinataire) : \esp{Creo que son para todos, porque no dependen del país donde nacemos.}
\item Exemple concret avec \esp{por} (cause) : \esp{Por ejemplo, en Guinea Ecuatorial, país hispanohablante vecino del Camerún, los jóvenes luchan por la educación para todos.}
\end{enumerate}
Intervention finale : \esp{En mi opinión, los derechos humanos son la base de la ciudadanía. Creo que son para todos, porque no dependen del país donde nacemos. Por ejemplo, en Guinea Ecuatorial los jóvenes luchan por la educación para todos.} Vérification : \esp{para todos} = destinataire (correct) ; \esp{luchar por} = cause, motif (correct).
\end{exoresolu}

\begin{experience}[Le débat en classe (30 à 40 minutes)]
\begin{enumerate}
\item \textbf{Préparation (10 min)} : chaque équipe rédige trois arguments sur papier, en vérifiant chaque \esp{por}/\esp{para} avec le tableau de la leçon.
\item \textbf{Round 1 — Opiniones (5 min)} : un porte-parole par équipe expose la position générale.
\item \textbf{Round 2 — Argumentos y réplicas (15 min)} : échange libre mais réglementé ; le modérateur alterne la parole ; chaque élève doit intervenir au moins une fois.
\item \textbf{Round 3 — Conclusiones (5 min)} : chaque équipe résume en deux phrases.
\item \textbf{Vote de la classe} : à main levée, y compris les indécis convaincus ; on peut voter contre son équipe si l'adversaire a été plus convaincant.
\end{enumerate}
\end{experience}

\begin{exercice}[À vous de débattre]
Rédigez par écrit votre intervention personnelle (5 à 6 phrases) sur le thème \esp{¿Somos ciudadanos del mundo?} en respectant le schéma opinion → argument → exemple. Vous devez employer : une formule d'opinion, un connecteur d'argumentation (\esp{porque}, \esp{ya que} ou \esp{por eso}), une formule de réaction à un adversaire imaginaire (\esp{sin embargo} ou \esp{no estoy de acuerdo}), et au moins un \esp{por} et un \esp{para} correctement employés.
\end{exercice}

\begin{corrige}[Exercice : modèle de réponse]
\esp{No estoy de acuerdo con los que dicen que sólo somos ciudadanos de nuestro país. En mi opinión, somos ciudadanos del mundo porque los problemas de hoy, como el clima, son para todos. Por eso estudio español: para entender a los jóvenes de Guinea Ecuatorial y de América Latina. Sin embargo, también tengo deberes por mi país: respeto las leyes y trabajo por el desarrollo del Camerún. Las dos cosas son posibles a la vez.}\par
« Je ne suis pas d'accord avec ceux qui disent que nous sommes seulement citoyens de notre pays. À mon avis, nous sommes citoyens du monde parce que les problèmes d'aujourd'hui, comme le climat, concernent tout le monde. C'est pourquoi j'étudie l'espagnol : pour comprendre les jeunes de Guinée équatoriale et d'Amérique latine. Cependant, j'ai aussi des devoirs envers mon pays : je respecte les lois et je travaille pour le développement du Cameroun. Les deux choses sont possibles à la fois. »\par
Vérification grammaticale : \esp{para todos} (destinataire), \esp{para entender} (but + infinitif), \esp{deberes por mi país} (cause/envers), \esp{trabajo por el desarrollo} (motif) : quatre emplois corrects.
\end{corrige}

\courssub{Bilan du débat : synthèse collective et vote}

Après le vote, le professeur conduit une \cle{synthèse collective} au tableau, rédigée en espagnol avec la classe. Elle peut prendre cette forme :

\begin{center}
\esp{Conclusión de la clase: somos ciudadanos de nuestro país \textbf{y} del mundo. Tenemos deberes por nuestra nación, pero los derechos humanos son para todos, sin fronteras.}
\end{center}

\begin{aretenir}[Ce qu'il faut retenir de cette séance]
\begin{itemize}
\item Débattre, c'est suivre le schéma \textbf{opinion → argument → exemple}, avec des formules figées : \esp{creo que}, \esp{en mi opinión}, \esp{porque}, \esp{ya que}, \esp{por eso}, \esp{estoy de acuerdo con}, \esp{sin embargo}.
\item « Être d'accord avec » se dit \esp{estar de acuerdo con} — jamais \esp{*estar acuerdo}.
\item Chaque argument doit contenir un \esp{por} (cause, motif, moyen) ou un \esp{para} (but, destinataire) correctement choisi : c'est la grammaire mise au service de la pensée.
\item La citoyenneté a deux dimensions, nationale et mondiale : un bon débat ne choisit pas l'une contre l'autre, il les articule.
\end{itemize}
\end{aretenir}

\newpage

\coursec{Activité d'intégration}

\courssub{Objectifs de l'activité}

À l'issue de cette activité d'intégration, l'élève sera capable de :
\begin{itemize}
\item traduire en français soigné deux vers d'un poème engagé espagnol, en respectant le sens, le rythme et les images (\cle{traduction littéraire} vs \cle{traduction poétique}) ;
\item rédiger un court plaidoyer citoyen en espagnol contenant au moins trois emplois corrects et justifiés de \esp{por} ou \esp{para} (\cle{choix prépositionnel}) ;
\item présenter oralement son travail et réagir aux productions des autres groupes avec les outils du \cle{débat guidé} (\cle{expression orale}, \cle{interaction}) ;
\item mobiliser le lexique de la citoyenneté et de la solidarité (\cle{los derechos humanos}, \cle{la solidaridad}, \cle{la tolerancia}...) dans une tâche authentique de communication.
\end{itemize}

Cette activité met en synergie les trois savoir-faire de la leçon : la traduction poétique, le choix entre \esp{por} et \esp{para}, et l'argumentation orale sur le thème \esp{Ciudadanos del mundo}.

\courssub{Situation}

Dans le cadre de la \esp{Semana de la Lengua Española} organisée par le club d'espagnol de ton lycée à Yaoundé, le proviseur a lancé un concours intitulé \esp{Un verso para el mundo}. Chaque groupe de quatre élèves de Terminale reçoit deux vers d'un poème engagé espagnol. La tâche : les traduire en français soigné, puis rédiger et présenter un plaidoyer citoyen. Le meilleur plaidoyer, élu par la classe, sera affiché au mur de la salle d'espagnol et lu lors de la cérémonie de clôture devant les élèves de Seconde.

Voici le document remis à chaque groupe (les deux vers sont les mêmes pour tous ; c'est le plaidoyer qui fera la différence) :

\begin{textebox}[Poema original : « Caminantes » (extrait)]
No hay fronteras para quien sueña con justicia,\\
caminamos por un mundo sin miedo y sin banderas.
\end{textebox}

\textbf{Glossaire :} \esp{frontera} = frontière ; \esp{soñar con} = rêver de ; \esp{justicia} = justice ; \esp{caminar} = marcher, cheminer ; \esp{miedo} = peur ; \esp{bandera} = drapeau ; \esp{quien} = celui qui / qui (pronom relatif antéposé à préposition).

\courssub{Compréhension et lexique thématique}

\begin{methode}[Comment aborder un texte poétique court]
\begin{enumerate}
\item \textbf{Lecture à voix haute :} repérer le rythme, les rimes (assonances : \esp{justicia / banderas}), la césure.
\item \textbf{Analyse grammaticale :} identifier le verbe principal (\esp{hay}, \esp{caminamos}), les prépositions (\esp{para}, \esp{por}), les compléments.
\item \textbf{Sens littéral :} traduire mot à mot pour verrouiller la structure.
\item \textbf{Interprétation :} qui parle ? (\cle{nous} collectif). Quelle image ? (\cle{cheminement}, \cle{frontières}, \cle{drapeaux} comme symboles de division).
\item \textbf{Rédaction de la version poétique :} chercher l'élégance française (inversion, vocabulaire soutenu, concision) sans trahir l'original.
\end{enumerate}
\end{methode}

\textbf{Tâche 1 — Compréhension du poème (10 min).} En groupe, lisez les deux vers à voix haute. Répondez par écrit, en espagnol, aux questions suivantes :
\begin{enumerate}
\item ¿Quién es el « \esp{quien} » del primer verso : una persona concreta o cualquier ser humano?
\item ¿Qué significa « \esp{sin banderas} » : que el mundo ideal no tiene países, o que los caminantes no llevan símbolos de odio ni división?
\item ¿Qué valores defiende el poema? Cita dos palabras del texto que lo demuestren.
\end{enumerate}

\courssub{Traduction de vers : méthode et pratique}

\begin{exemplebox}[Traduction littérale vs poétique]
\textbf{Original :} \esp{No hay fronteras para quien sueña con justicia, / caminamos por un mundo sin miedo y sin banderas.}

\textbf{1. Traduction littérale (version) :} « Il n'y a pas de frontières pour celui qui rêve de justice, / nous cheminons pour un monde sans peur et sans drapeaux. »
\par \textbf{2. Traduction poétique (création) :} « Nulle frontière pour qui rêve de justice : / marchons vers un monde sans peur ni bannière. »

\emph{Justification des choix :}
\begin{itemize}
\item \esp{No hay fronteras} $\rightarrow$ « Nulle frontière » (négation littéraire, plus fort que « il n'y a pas de »).
\item \esp{para quien sueña} $\rightarrow$ « pour qui rêve » (conserve la préposition de destinataire \esp{para}, rythme binaire).
\item \esp{caminamos} $\rightarrow$ « marchons » (impératif exhortatif / \cle{hortatif} pour engager le lecteur, ou présent de narration « nous marchons »).
\item \esp{por un mundo} $\rightarrow$ « vers un monde » (\esp{por} exprime le but du mouvement, le français préfère « vers » pour la direction/but).
\item \esp{sin miedo y sin banderas} $\rightarrow$ « sans peur ni bannière » (\cle{ni} littéraire pour \esp{y} coordonnant deux négations, \esp{bannière} plus solennel que \esp{drapeau}).
\end{itemize}
\end{exemplebox}

\textbf{Tâche 2 — Traduction des vers (15 min).} Traduisez les deux vers en français soigné. Deux propositions sont attendues : une traduction \emph{littérale} (fidèle au sens, mot à mot) et une traduction \emph{poétique} (qui respecte le rythme et, si possible, la musicalité). Justifiez en deux phrases le choix que vous avez fait pour traduire \esp{para quien sueña} et \esp{por un mundo}.

\courssub{Grammaire : POR et PARA — Règle de choix}

\begin{propriete}[Tableau récapitulatif : POR vs PARA]
\begin{center}
\begin{tabularx}{\linewidth}{|>{\bfseries}X|X|X|X|}
\hline
\rowcolor{popblueL}
\textbf{Valeur} & \textbf{POR} & \textbf{PARA} & \textbf{Test de substitution (français)} \\
\hline
\cle{Cause / Motif} & \esp{Lucho \textbf{por} la justicia.} & -- & « À cause de », « en raison de », « pour défendre » \\
\hline
\cle{But / Objectif (nom)} & \esp{Trabajamos \textbf{por} un mundo mejor.} & -- & « Dans le but d'obtenir », « en vue de » (nom) \\
\hline
\cle{Finalité (verbe)} & -- & \esp{Estudio \textbf{para} comprender.} & « Pour » + \textbf{infinitif} (but) \\
\hline
\cle{Destinataire} & -- & \esp{Este regalo es \textbf{para} ti.} & « Pour » (bénéficiaire), « à l'intention de » \\
\hline
\cle{Durée} & \esp{Estudié \textbf{por} dos horas.} & -- & « Pendant », « durant » \\
\hline
\cle{Date limite} & -- & \esp{La tarea es \textbf{para} mañana.} & « Pour » (échéance) \\
\hline
\cle{Lieu (traversée)} & \esp{Caminamos \textbf{por} el parque.} & -- & « À travers », « par », « dans » \\
\hline
\cle{Direction / Destination} & -- & \esp{Salgo \textbf{para} Madrid.} & « Pour », « vers », « en direction de » \\
\hline
\cle{Échange / Prix} & \esp{Lo compré \textbf{por} diez euros.} & -- & « Contre », « en échange de » \\
\hline
\cle{Agent (passif)} & \esp{Fue escrito \textbf{por} Lorca.} & -- & « Par » (agent) \\
\hline
\cle{Substitution} & -- & \esp{Para mí, es fácil.} & « Selon », « pour ma part » \\
\hline
\end{tabularx}
\end{center}
\end{propriete}

\begin{methode}[Algorithme de choix : POR ou PARA ?]
\begin{enumerate}
\item \textbf{Devant un infinitif ?} $\rightarrow$ Presque toujours \esp{PARA} (finalité). \emph{Ex : \esp{Vengo para ayudar.}}
\item \textbf{Indique un bénéficiaire / un destinataire ?} $\rightarrow$ \esp{PARA}. \emph{Ex : \esp{Es para ti.}}
\item \textbf{Indique une échéance (date/heure) ?} $\rightarrow$ \esp{PARA}. \emph{Ex : \esp{Para el lunes.}}
\item \textbf{Indique une direction physique ?} $\rightarrow$ \esp{PARA}. \emph{Ex : \esp{Salgo para Douala.}}
\item \textbf{Sinon (cause, durée, lieu vague, échange, agent, motif) ?} $\rightarrow$ \esp{POR}.
\end{enumerate}
\begin{center}
\begin{tikzpicture}[node distance=1.2cm, >=Stealth, font=\small]
\node[draw, rounded corners, fill=popblueL] (start) {Devant un nom / pronom ?};
\node[draw, diamond, aspect=2, below of=start, fill=poporangeL] (inf) {Devant un infinitif ?};
\node[draw, rounded corners, right of=inf, xshift=2.5cm, fill=popgreenL] (para) {\textbf{PARA} (Finalité)};
\node[draw, diamond, aspect=2, below of=inf, fill=poporangeL] (dest) {Destinataire / Bénéficiaire ?};
\node[draw, rounded corners, right of=dest, xshift=2.5cm, fill=popgreenL] (para2) {\textbf{PARA} (Destinataire)};
\node[draw, diamond, aspect=2, below of=dest, fill=poporangeL] (date) {Date limite / Échéance ?};
\node[draw, rounded corners, right of=date, xshift=2.5cm, fill=popgreenL] (para3) {\textbf{PARA} (Échéance)};
\node[draw, diamond, aspect=2, below of=date, fill=poporangeL] (dir) {Direction physique (aller vers) ?};
\node[draw, rounded corners, right of=dir, xshift=2.5cm, fill=popgreenL] (para4) {\textbf{PARA} (Direction)};
\node[draw, rounded corners, below of=dir, fill=poppinkL] (por) {\textbf{POR} (Cause, Durée, Lieu, Échange, Agent, Motif...)};
\draw[->] (start) -- (inf);
\draw[->] (inf) -- node[above, midway] {OUI} (para);
\draw[->] (inf) -- node[left, midway] {NON} (dest);
\draw[->] (dest) -- node[above, midway] {OUI} (para2);
\draw[->] (dest) -- node[left, midway] {NON} (date);
\draw[->] (date) -- node[above, midway] {OUI} (para3);
\draw[->] (date) -- node[left, midway] {NON} (dir);
\draw[->] (dir) -- node[above, midway] {OUI} (para4);
\draw[->] (dir) -- node[left, midway] {NON} (por);
\end{tikzpicture}
\end{center}
\legende{Organigramme de décision POR / PARA (simplifié pour la production écrite).}
\end{methode}

\begin{aretenir}[Autres valeurs à reconnaître (niveau Terminale)]
\begin{itemize}
\item \esp{POR} : \textbf{Agent passif} (\esp{Escrito por García Márquez}), \textbf{Moyen/Transport} (\esp{Envié por correo, Voy por tren}), \textbf{Substitution} (\esp{Voto por él = en su lugar}), \textbf{Distribution} (\esp{Un por uno}).
\item \esp{PARA} : \textbf{Opinion} (\esp{Para mí, es injusto}), \textbf{Comparaison} (\esp{Para su edad, es alto}), \textbf{Sur le point de} (\esp{Estaba para salir}), \textbf{Proportion} (\esp{Trabaja mucho para lo que gana}).
\end{itemize}
\end{aretenir}

\begin{attention}[Pièges fréquents des francophones — \cle{Faux amis} et \cle{Calques}]
\begin{itemize}
\item \textbf{« Pour » + infinitif de but} = \esp{PARA} + infinitif, \textbf{jamais POR}. \esp{Vengo \textbf{para} ayudar.} (Pas \esp{*por ayudar}).
\item \textbf{« Remercier pour »} = \esp{agradecer / dar las gracias \textbf{POR}}. \esp{Gracias \textbf{por} vuestra atención.}
\item \textbf{« Lutter pour »} = \esp{luchar \textbf{POR}} ; \textbf{« Rêver de »} = \esp{soñar \textbf{CON}} (ni \esp{por} ni \esp{para} !).
\item \textbf{« S'inquiéter pour »} = \esp{preocuparse \textbf{POR}} ; \textbf{« S'occuper de »} = \esp{ocuparse \textbf{DE}}.
\item \textbf{Accents obligatoires :} \esp{está} (verbe), \esp{también} (aussi), \esp{corazón}, \esp{atención}, \esp{país}, \esp{inglés}, \esp{estación}, \esp{acción}, \esp{razón}, \esp{lección}.
\item \textbf{Ser / Estar :} \esp{Es ciudadano} (identité, nationalité) vs \esp{Está en Yaoundé} (localisation). Dans le plaidoyer : \esp{Soy ciudadano} (identité essentielle).
\end{itemize}
\end{attention}

\courssub{Ressources lexicales et outils du débat}

\begin{definition}[Lexique de la citoyenneté et de la solidarité (Módulo 3)]
\begin{center}
\begin{tabularx}{\linewidth}{|l|X|l|X|}
\hline
\rowcolor{popblueL}
\textbf{Español} & \textbf{Français} & \textbf{Español} & \textbf{Français} \\
\hline
\esp{los derechos humanos} & les droits de l'homme & \esp{la solidaridad} & la solidarité \\
\esp{la tolerancia} & la tolérance & \esp{el respeto} & le respect \\
\esp{la paz} & la paix & \esp{el diálogo} & le dialogue \\
\esp{la igualdad} & l'égalité & \esp{la justicia} & la justice \\
\esp{la libertad} & la liberté & \esp{la fraternidad} & la fraternité \\
\esp{compartir} & partager & \esp{defender} & défendre \\
\esp{comprometerse (con)} & s'engager (pour) & \esp{convivir} & cohabiter, vivre ensemble \\
\esp{la diversidad} & la diversité & \esp{la inclusión} & l'inclusion \\
\hline
\end{tabularx}
\end{center}
\end{definition}

\begin{definition}[Outils du débat guidé — \cle{Interaction orale}]
\begin{center}
\begin{tabularx}{\linewidth}{|>{\bfseries}l|X|}
\hline
\rowcolor{popblueL}
\textbf{Fonction} & \textbf{Formules types} \\
\hline
Accord & \esp{Estoy de acuerdo con vosotros porque...} / \esp{Comparto vuestra opinión.} \\
\hline
Désaccord poli & \esp{No estoy de acuerdo. / Sin embargo, creo que...} / \esp{Me parece que no es así porque...} \\
\hline
Nuance & \esp{Es verdad, pero...} / \esp{Por un lado... por otro lado...} \\
\hline
Demander avis & \esp{¿Qué pensáis de...?} / \esp{¿No creéis que...?} / \esp{¿Estás de acuerdo?} \\
\hline
Argumenter & \esp{Mi argumento es que...} / \esp{La razón es que...} / \esp{Por ejemplo...} \\
\hline
Conclure & \esp{En conclusión...} / \esp{Por todo ello...} / \esp{Por tanto...} \\
\hline
\end{tabularx}
\end{center}
\end{definition}

\courssub{Production écrite : Plaidoyer « Soy ciudadano del mundo porque... »}

\textbf{Tâche 3 — Rédaction du plaidoyer (20 min).} Rédigez un plaidoyer de 6 à 8 phrases intitulé \esp{Soy ciudadano del mundo porque...}. Contraintes obligatoires :
\begin{enumerate}
\item Au moins trois emplois de \esp{por} ou \esp{para}, de valeurs \textbf{différentes} (cause, finalité, destinataire, temps, lieu, échange, agent...).
\item Au moins deux mots du lexique de la citoyenneté (cf. tableau ci-dessus).
\item Au moins une phrase commençant par \esp{Creo que...} ou \esp{Estoy convencido de que...}.
\item Une conclusion forte qui reprend l'image du poème (le chemin, le rêve, les frontières).
\end{enumerate}

\courssub{Expression orale : Présentation et débat}

\textbf{Tâche 4 — Présentation orale et débat (20 min).} Chaque groupe désigne un porte-parole qui lit la traduction puis le plaidoyer. Après chaque passage, les autres groupes réagissent avec les outils du débat. Le porte-parole doit répondre à au moins une objection.

\textbf{Tâche 5 — Vote et affichage (5 min).} La classe vote à main levée pour le meilleur plaidoyer (on ne vote pas pour son propre groupe). Le texte gagnant est recopié au propre et affiché au mur de la salle d'espagnol.

\courssub{Proposition de production et corrigé commenté}

\begin{exoresolu}[Modèle de traduction et de plaidoyer]
\textbf{Énoncé.} Traduisez les deux vers du poème « Caminantes » et rédigez le plaidoyer « Soy ciudadano del mundo porque... » (6 à 8 phrases, trois emplois justifiés de \esp{por}/\esp{para}).

\tcblower

\textbf{1. Traduction littérale :} « Il n'y a pas de frontières pour celui qui rêve de justice, / nous cheminons pour un monde sans peur et sans drapeaux. »

\textbf{2. Traduction poétique :} « Nulle frontière pour qui rêve de justice : / marchons vers un monde sans peur ni bannière. »

\emph{Commentaire :} \esp{para quien sueña} est traduit par « pour qui rêve » (destinataire / bénéficiaire de l'absence de frontières) ; la version poétique resserre le vers (« nulle frontière ») pour retrouver le rythme hendécasyllabique implicite. \esp{por un mundo} exprime le but de la marche (cause finale) : « pour un monde » $\rightarrow$ « vers un monde » (préposition de direction en français). On a choisi « bannière » (registre soutenu) pour éviter la répétition « drapeaux » et garder la musicalité (assonance en \textipa{/iɛʁ/}).

\textbf{3. Plaidoyer modèle :}

\esp{Soy ciudadano del mundo porque creo que los derechos humanos no tienen fronteras. Lucho \textbf{por} la justicia y \textbf{por} la igualdad, en mi barrio de Yaoundé y en todas partes. Estudio español \textbf{para} comprender otras culturas y \textbf{para} dialogar con los jóvenes de Guinea Ecuatorial y de América Latina. \textbf{Estoy convencido de que} la solidaridad es un regalo \textbf{para} todos, no solo para los amigos. Doy mi tiempo \textbf{por} los demás, porque compartir es vivir. \textbf{Sin embargo}, sé que el camino es largo : caminamos \textbf{por} un mundo sin miedo, paso a paso. Este poema es \textbf{para} quien sueña : \textbf{¡soñemos juntos!}}

\emph{Traduction :} « Je suis citoyen du monde parce que je crois que les droits de l'homme n'ont pas de frontières. Je lutte pour la justice et pour l'égalité, dans mon quartier de Yaoundé et partout. J'étudie l'espagnol pour comprendre d'autres cultures et pour dialoguer avec les jeunes de Guinée équatoriale et d'Amérique latine. Je suis convaincu que la solidarité est un cadeau pour tous, pas seulement pour les amis. Je donne mon temps pour les autres, car partager, c'est vivre. Cependant, je sais que le chemin est long : nous cheminons pour un monde sans peur, pas à pas. Ce poème est pour celui qui rêve : rêvons ensemble ! »

\emph{Commentaire des choix de langue (justification grammaticale) :}
\begin{itemize}
\item \esp{Lucho \textbf{por} la justicia / igualdad} : \esp{POR} de \textbf{cause / motif} (on lutte \emph{en faveur de}, \emph{à cause de} la justice).
\item \esp{Estudio español \textbf{para} comprender / dialogar} : \esp{PARA} + \textbf{infinitif de finalité} (but poursuivi par le sujet).
\item \esp{Un regalo \textbf{para} todos} : \esp{PARA} de \textbf{destinataire / bénéficiaire}.
\item \esp{Doy mi tiempo \textbf{por} los demás} : \esp{POR} d'\textbf{échange / substitution} (je donne mon temps \emph{en échange de} leur bien-être / \emph{à la place de} eux).
\item \esp{Caminamos \textbf{por} un mundo sin miedo} : \esp{POR} de \textbf{but / objectif nominal} (reprise fidèle de l'image du poème : « nous marchons \emph{en vue d'un} monde... »).
\item \esp{Este poema es \textbf{para} quien sueña} : \esp{PARA} de \textbf{destinataire} (le poème s'adresse aux rêveurs).
\item Connecteurs du débat intégrés : \esp{Sin embargo} (concession), \esp{Estoy convencido de que} (opinion ferme + indicatif), \esp{porque} (cause), \esp{por tanto} (conséquence implicite).
\end{itemize}
\end{exoresolu}

\courssub{Entraînement individuel : Thème et Version}

\begin{exercice}[Application : POR / PARA en contexte]
Complétez les phrases suivantes par \esp{POR} ou \esp{PARA}. Justifiez brièvement en français (valeur).
\begin{enumerate}
\item Los estudiantes de Terminale estudian \_\_\_\_\_\_ aprobar el bac.
\item \_\_\_\_\_\_ mi parte, la paz es indispensable.
\item Caminamos \_\_\_\_\_\_ las calles de Douala \_\_\_\_\_\_ manifestar \_\_\_\_\_\_ la paz.
\item Este libro fue escrito \_\_\_\_\_\_ un autor camerunés.
\item Gracias \_\_\_\_\_\_ vuestra solidaridad.
\item Trabajo \_\_\_\_\_\_ ganar dinero \_\_\_\_\_\_ mi familia.
\item El avión sale \_\_\_\_\_\_ Madrid \_\_\_\_\_\_ diez minutos.
\item \_\_\_\_\_\_ quién es este regalo? — \_\_\_\_\_\_ mi hermano.
\end{enumerate}
\end{exercice}

\begin{corrige}[Exercice : POR / PARA]
\begin{enumerate}
\item \esp{PARA} (Finalité + infinitif : \esp{estudian \textbf{para} aprobar}).
\item \esp{PARA} (Opinion / Point de vue : \esp{\textbf{Para} mi parte}).
\item \esp{POR} (Lieu / Traversée : \esp{caminamos \textbf{por} las calles}) ; \esp{PARA} (Finalité + infinitif : \esp{manifestar \textbf{para} la paz} — but de la manifestation). \emph{Note :} \esp{manifestar por la paz} (cause) est aussi possible, mais \esp{para} insiste sur l'objectif visé.
\item \esp{POR} (Agent passif : \esp{fue escrito \textbf{por} un autor}).
\item \esp{POR} (Remerciement / Cause : \esp{Gracias \textbf{por} vuestra solidaridad}).
\item \esp{PARA} (Finalité : \esp{trabajo \textbf{para} ganar dinero}) ; \esp{POR} (Destinataire/Bénéficiaire au sens large / Cause : \esp{ganar dinero \textbf{por} mi familia} = pour subvenir à leurs besoins). \emph{Variante :} \esp{para mi familia} (destinataire direct de l'argent).
\item \esp{PARA} (Direction : \esp{sale \textbf{para} Madrid}) ; \esp{POR} (Durée / Échéance proche : \esp{sale \textbf{por} diez minutos} = « part dans dix minutes »). \emph{Attention :} \esp{para diez minutos} = « pour dans dix minutes » (échéance).
\item \esp{PARA} (Destinataire : \esp{\textbf{Para} quién}) ; \esp{PARA} (Destinataire : \esp{\textbf{Para} mi hermano}).
\end{enumerate}
\end{corrige}

\begin{exercice}[Version : Traduire en espagnol]
Traduisez les phrases suivantes en espagnol en soignant le choix \esp{por} / \esp{para}.
\begin{enumerate}
\item Je me bats pour la liberté. (luchar)
\item Ce cadeau est pour toi.
\item J'ai marché pendant trois heures. (caminar)
\item Nous étudions l'espagnol pour voyager. (viajar)
\item Le pont a été construit par des ingénieurs. (construir / ingenieros)
\item Je vais à Yaoundé demain. (salir / mañana)
\item Merci pour votre aide. (ayuda)
\item Il est sur le point de partir. (estar para)
\end{enumerate}
\end{exercice}

\begin{corrige}[Exercice : Version]
\begin{enumerate}
\item \esp{Lucho \textbf{por} la libertad.} (Cause)
\item \esp{Este regalo es \textbf{para} ti.} (Destinataire)
\item \esp{Caminé \textbf{por} tres horas.} (Durée)
\item \esp{Estudiamos español \textbf{para} viajar.} (Finalité + infinitif)
\item \esp{El puente fue construido \textbf{por} ingenieros.} (Agent passif)
\item \esp{Salgo \textbf{para} Yaoundé mañana.} (Direction + temps)
\item \esp{Gracias \textbf{por} vuestra ayuda.} (Cause / Remerciement)
\item \esp{Está \textbf{para} salir.} (Sur le point de + infinitif)
\end{enumerate}
\end{corrige}

\courssub{Bilan de l'activité}

Cette activité a permis de vérifier l'acquisition des trois savoir-faire de la leçon dans une tâche authentique et motivante :
\begin{itemize}
\item \textbf{Traduction poétique :} l'élève a distingué traduction littérale (version) et traduction d'élégance (création), et justifié ses choix (rythme, images, sens, registres de langue).
\item \textbf{Grammaire :} l'élève a employé \esp{por} (cause, but nominal, durée, lieu, échange, agent) et \esp{para} (finalité verbale, destinataire, direction, échéance, opinion) de manière correcte et consciente, en évitant les calques du français (\esp{*por ayudar} $\rightarrow$ \esp{para ayudar}).
\item \textbf{Expression orale et débat :} l'élève a défendu un point de vue sur la citoyenneté mondiale, réagi aux objections avec les connecteurs logiques (\esp{Sin embargo}, \esp{Por tanto}, \esp{Es verdad pero...}), et écouté les arguments des autres pour construire une pensée collective.
\end{itemize}

Le plaidoyer gagnant, affiché au mur de la salle d'espagnol, devient un texte de référence pour toute la classe : il pourra être relu et enrichi lors du prochain module. À la maison, chaque élève recopie le plaidoyer de son groupe en corrigeant les erreurs signalées pendant le débat, et ajoute une phrase supplémentaire avec un \textbf{nouvel emploi} de \esp{por} ou \esp{para} (ex: \esp{por} agent passif, \esp{para} opinion, \esp{por} moyen de transport).

\begin{savaistu}[Le saviez-vous ? — La Guinée équatoriale, seule nation hisphone d'Afrique]
La \esp{Guinea Ecuatorial} (capitale : \esp{Malabo}, sur l'île de Bioko ; \esp{Bata} sur le continent) est le seul pays africain où l'espagnol est langue officielle (avec le français et le portugais). Elle a obtenu son indépendance de l'Espagne en 1968. Sa littérature (Juan Tomás Ávila Laurel, Donato Ndongo-Bidyogo) explore l'identité, le post-colonialisme et la \cle{ciudadanía}. Yaoundé et Malabo sont reliées par une coopération universitaire et culturelle active. Être \esp{ciudadano del mundo} au Cameroun, c'est aussi connaître son voisin hispanophone !
\end{savaistu}

\newpage

\coursec{Méthodes et exercices résolus}

\coursec{Leçon E12 : Traduction de quelques vers d'un poème ; POR y PARA ; Debate guiado : « Ciudadanos del mundo »}

\begin{textebox}[Poème original — Support de la leçon]
\esp{No quiero ser un árbol sin raíces,}\\
\esp{quiero la libertad de mi mañana,}\\
\esp{porque el silencio nunca fue mi lengua}\\
\esp{y el pan del pueblo sabe a esperanza.}
\legende{Poème engagé original, écrit pour le cours. Thème : identité, liberté, justice sociale.}
\end{textebox}

\courssub{Savoir-faire 1 : Traduire quelques vers d'un poème}

\begin{methode}[Traduire un extrait de poème en vers]
\begin{enumerate}
\item \textbf{Lire le poème entier au moins deux fois} avant de toucher au stylo : repérer le thème (l'amour, la patrie, la liberté, la nature), le ton (tendre, engagé, mélancolique) et le locuteur (\esp{yo poético}).
\item \textbf{Identifier le sens global de chaque vers}, pas mot à mot. Un poème se traduit par \emph{idées} : chercher d'abord ce que le vers \emph{veut dire}, puis comment le dire en français.
\item \textbf{Repérer les procédés poétiques} : métaphore (\esp{mi pueblo es un río} = mon peuple est un fleuve), anaphore (répétition en tête de vers), hyperbole, personnification. Les conserver dans la traduction autant que possible.
\item \textbf{Traduire le lexique avec précision} : vérifier chaque mot dans le contexte (attention aux faux amis : \esp{largo} = long, pas « large » ; \esp{sentir} = ressentir/regretter selon le contexte).
\item \textbf{Reconstituer une phrase française correcte} : l'ordre des mots espagnol est souvent inversé en poésie (\esp{de mi tierra el dolor} = la douleur de ma terre). Remettre les compléments à leur place naturelle en français.
\item \textbf{Relire la traduction à voix haute} : elle doit être fluide, fidèle et garder le rythme ou l'émotion de l'original. Ne jamais ajouter d'idées absentes du texte.
\end{enumerate}
\end{methode}

\begin{exoresolu}[Traduction de vers engagés — Type Bac]
\textbf{Énoncé.} Traduis en français les vers suivants :

\begin{center}
\esp{No quiero ser un árbol sin raíces,}\\
\esp{quiero la libertad de mi mañana,}\\
\esp{porque el silencio nunca fue mi lengua}\\
\esp{y el pan del pueblo sabe a esperanza.}
\end{center}

\tcblower
\textbf{Solution détaillée.}

\textbf{Étape 1 — Sens global.} Le poète refuse d'être déraciné ; il revendique la liberté, la parole et l'espoir du peuple. Ton : engagé, volontaire.

\textbf{Étape 2 — Analyse vers par vers.}
\begin{itemize}
\item Vers 1 : \esp{árbol sin raíces} = métaphore de l'homme coupé de sa culture. \esp{raíces} = racines (faux ami : pas « raisons »).
\item Vers 2 : \esp{mi mañana} = mon avenir (métaphore temporelle : le matin = l'avenir).
\item Vers 3 : \esp{el silencio nunca fue mi lengua} = personnification : le silence n'a jamais été ma langue → le poète a toujours parlé, dénoncé. Attention au passé simple \esp{fue} (ser, indéfini) : « n'a jamais été » (passé composé français).
\item Vers 4 : \esp{saber a} = avoir le goût de. \esp{el pan del pueblo} = le pain du peuple (symbole de la vie quotidienne, de la dignité).
\end{itemize}

\textbf{Étape 3 — Traduction rédigée.}

« Je ne veux pas être un arbre sans racines,\\
je veux la liberté de mon avenir,\\
car le silence n'a jamais été ma langue\\
et le pain du peuple a le goût de l'espérance. »

\textbf{Conclusion.} La traduction respecte les métaphores (arbre/racines, matin/avenir, pain/espérance), l'ordre français naturel et le ton engagé. On n'a ni ajouté ni supprimé d'idée.
\end{exoresolu}

\courssub{Savoir-faire 2 : Choisir correctement entre POR et PARA}

\begin{propriete}[Valeurs fondamentales de POR et PARA]
\begin{center}
\begin{tabularx}{\linewidth}{|X|X|X|}
\hline
\rowcolor{popblueL} \textbf{Valeur} & \textbf{Préposition} & \textbf{Exemple traduit} \\
\hline
Cause, motif, raison (« à cause de ») & \esp{POR} & \esp{Lucho \textbf{por} la justicia.} (Je lutte pour la justice.) \\
\hline
Durée, trajet, lieu vague, moment approximatif & \esp{POR} & \esp{Caminamos \textbf{por} la ciudad \textbf{por} la noche.} (Nous marchons dans la ville le soir.) \\
\hline
Agent du passif, moyen, échange, prix & \esp{POR} & \esp{Escrito \textbf{por} Cervantes.} (Écrit par Cervantes.) \esp{Pagué 1000 francs \textbf{por} el libro.} \\
\hline
But, finalité (« afin de », « pour » + infinitif) & \esp{PARA} & \esp{Estudio \textbf{para} ser médico.} (J'étudie pour devenir médecin.) \\
\hline
Destination, échéance, date limite & \esp{PARA} & \esp{Salgo \textbf{para} Yaundé.} (Je pars pour Yaoundé.) \esp{La tarea es \textbf{para} mañana.} \\
\hline
Destinataire, opinion, comparaison & \esp{PARA} & \esp{Este regalo es \textbf{para} ti.} (Ce cadeau est pour toi.) \esp{\textbf{Para} mí, es injusto.} \\
\hline
\end{tabularx}
\end{center}
\end{propriete}

\begin{popfigure}
\begin{tikzpicture}[
  node distance=1.2cm and 1.5cm,
  decision/.style={diamond, draw=popdark, thick, fill=popblueL, text width=2.5cm, align=center, inner sep=4pt},
  action/.style={rectangle, rounded corners, draw=popdark, thick, fill=popgreenL, text width=2.8cm, align=center, inner sep=4pt},
  start/.style={ellipse, draw=popdark, thick, fill=popgoldL, text width=2cm, align=center, inner sep=4pt},
  arrow/.style={-Stealth, thick, popdark}
]
\node[start] (start) {Choisir POR / PARA};
\node[decision, below of=start, align=center] (q1){Question : \emph{¿Por qué?}\\(Cause / But ?)};
\node[action, below left=1.5cm and 1cm of q1, align=center] (por1){\textbf{POR}\\\textbullet Cause, motif\\\textbullet Durée, trajet\\\textbullet Agent passif\\\textbullet Moyen, prix};
\node[action, below right=1.5cm and 1cm of q1, align=center] (para1){\textbf{PARA}\\\textbullet But, finalité\\\textbullet Destination\\\textbullet Destinataire\\\textbullet Opinion};
\node[decision, below=2.5cm of q1, align=center] (q2){Question : \emph{¿Dónde? ¿Cuándo?}\\(Lieu / Temps ?)};
\node[action, below left=1.5cm and 1cm of q2, align=center] (por2){\textbf{POR}\\\textbullet Trajet, lieu vague\\\textbullet Durée, moment approx.};
\node[action, below right=1.5cm and 1cm of q2, align=center] (para2){\textbf{PARA}\\\textbullet Destination précise\\\textbullet Date limite (échéance)};
\draw[arrow] (start) -- (q1);
\draw[arrow] (q1) -- node[left, pos=0.5, align=center] {Cause / Durée /\\(Agent / Moyen} (por1);
\draw[arrow] (q1) -- node[right, pos=0.5, align=center] {But / Destination /\\(Destinataire} (para1);
\draw[arrow] (q1) -- (q2);
\draw[arrow] (q2) -- node[left, pos=0.5, align=center] {Trajet / Durée} (por2);
\draw[arrow] (q2) -- node[right, pos=0.5, align=center] {Destination / Échéance} (para2);
\end{tikzpicture}
\legende{Arbre de décision pour choisir entre POR et PARA.}
\end{popfigure}

\begin{methode}[Choisir entre POR et PARA : l'arbre de décision]
\begin{enumerate}
\item \textbf{Poser la question « pourquoi ? » ou « dans quel but ? »}
\begin{itemize}
\item \textbf{Cause, motif, raison} (« à cause de, en raison de ») $\rightarrow$ \cle{POR}. \esp{Lo hago \textbf{por} amor.} « Je le fais par amour. »
\item \textbf{But, finalité} (« afin de, pour + infinitif ») $\rightarrow$ \cle{PARA}. \esp{Estudio \textbf{para} ser médico.} « J'étudie pour devenir médecin. »
\end{itemize}
\item \textbf{Poser la question « où / quand ? »}
\begin{itemize}
\item \textbf{Trajet, durée, moment approximatif} $\rightarrow$ \cle{POR}. \esp{Paseamos \textbf{por} el mercado.} « Nous nous promenons dans le marché. » \esp{Trabajé \textbf{por} dos horas.} « J'ai travaillé pendant deux heures. »
\item \textbf{Destination, date limite} $\rightarrow$ \cle{PARA}. \esp{Salgo \textbf{para} Yaundé.} « Je pars pour Yaoundé. » \esp{La tarea es \textbf{para} mañana.} « Le devoir est pour demain. »
\end{itemize}
\item \textbf{Autres cas fixes à mémoriser :}
\begin{itemize}
\item \textbf{Agent du passif, moyen, échange, prix} $\rightarrow$ \cle{POR}. \esp{Fue escrito \textbf{por} Cervantes.} « Écrit par Cervantes. » \esp{Pagué mil francos \textbf{por} el libro.} « J'ai payé mille francs pour le livre. »
\item \textbf{Destinataire, opinion, comparaison} $\rightarrow$ \cle{PARA}. \esp{Este regalo es \textbf{para} ti.} « Ce cadeau est pour toi. » \esp{\textbf{Para} mí, la libertad es sagrada.} « Pour moi, la liberté est sacrée. »
\end{itemize}
\item \textbf{Expressions figées :} \esp{por ejemplo} (par exemple), \esp{por fin} (enfin), \esp{por eso} (c'est pourquoi), \esp{para siempre} (pour toujours), \esp{para qué} (à quoi bon).
\item \textbf{Vérification finale :} remplacer mentalement par « à cause de » (por) ou « dans le but de » (para) : si la phrase garde son sens, le choix est bon.
\end{enumerate}
\end{methode}

\begin{exoresolu}[Thème : POR ou PARA ? — Type Bac]
\textbf{Énoncé.} Traduis en espagnol :
\begin{enumerate}
\item Je lutte pour la justice et pour mes enfants.
\item Nous marchons dans la ville le soir.
\item Ce poème a été écrit par un jeune de Yaoundé pour sa mère.
\item Il a voyagé en bus pour voir sa grand-mère.
\end{enumerate}

\tcblower
\textbf{Solution détaillée.}

\begin{enumerate}
\item « Je lutte \textbf{pour} la justice (cause, motif) et \textbf{pour} mes enfants (destinataire/bénéficiaires). »
$\rightarrow$ \esp{Lucho \textbf{por} la justicia y \textbf{para} mis hijos.}
Justification : la justice est le \emph{motif} de la lutte (\esp{por}) ; les enfants sont les \emph{bénéficiaires} (\esp{para}).

\item « Nous marchons \textbf{dans} la ville \textbf{le} soir. » (trajet dans un lieu + moment approximatif)
$\rightarrow$ \esp{Caminamos \textbf{por} la ciudad \textbf{por} la noche.}
Justification : \esp{por + lieu} = à travers, dans ; \esp{por la noche} = expression de temps figée. \esp{Caminar} (marcher) préféré à \esp{pasear} (se promener) pour traduire « marcher ».

\item « Ce poème a été écrit \textbf{par} un jeune (agent du passif) de Yaoundé \textbf{pour} sa mère (destinataire). »
$\rightarrow$ \esp{Este poema fue escrito \textbf{por} un joven de \textbf{Yaundé} \textbf{para} su madre.}
Justification : l'agent d'une voix passive exige toujours \esp{por} ; la mère reçoit le poème, donc \esp{para}. Toponyme espagnol : \esp{Yaundé}.

\item « Il a voyagé en bus (moyen) \textbf{pour} voir sa grand-mère (but). »
$\rightarrow$ \esp{Viajó \textbf{por} autobús \textbf{para} ver a su abuela.}
Justification : le moyen de transport se dit avec \esp{por} (\esp{por autobús}, aussi \esp{en autobús}) ; \esp{para + infinitif} exprime la finalité. Ne pas oublier l'\esp{a} personnelle devant \esp{su abuela}.
\end{enumerate}

\textbf{Conclusion.} À chaque occurrence, on a identifié la valeur (cause, but, agent, destinataire, moyen) avant de choisir la préposition : c'est la seule méthode fiable.
\end{exoresolu}

\courssub{Savoir-faire 3 : Version — reconnaître POR et PARA dans un texte}

\begin{textebox}[Texte original pour version — Contexte citoyen]
\esp{Los jóvenes luchan por sus derechos y por un futuro mejor. Para ellos, la educación es la llave del mundo. Este programa fue creado por la comunidad para los barrios pobres, y por eso todos trabajan para construir una ciudad más justa.}
\legende{Texte original, niveau Terminale, thème : engagement jeunesse et citoyenneté.}
\end{textebox}

\begin{methode}[Traduire POR et PARA de l'espagnol vers le français]
\begin{enumerate}
\item \textbf{Repérer chaque occurrence} de \esp{por} ou \esp{para} dans le texte et la souligner.
\item \textbf{Déterminer sa valeur} d'après le contexte : cause, but, agent, durée, destination, échange, destinataire.
\item \textbf{Choisir l'équivalent français adapté}, qui n'est pas toujours « pour » :
\begin{itemize}
\item \esp{por} = par, à cause de, pendant, à travers, en échange de, de la part de ;
\item \esp{para} = pour, afin de, vers (destination), d'ici (date limite).
\end{itemize}
\item \textbf{Traduire les expressions figées en bloc} : \esp{por eso} = c'est pourquoi ; \esp{para siempre} = pour toujours ; \esp{por fin} = enfin.
\item \textbf{Relire la phrase française entière} pour vérifier qu'elle est idiomatique : une version mot à mot est sanctionnée.
\end{enumerate}
\end{methode}

\begin{exoresolu}[Version : texte citoyen — Type Bac]
\textbf{Énoncé.} Traduis en français le passage suivant :

\esp{Los jóvenes luchan por sus derechos y por un futuro mejor. Para ellos, la educación es la llave del mundo. Este programa fue creado por la comunidad para los barrios pobres, y por eso todos trabajan para construir una ciudad más justa.}

\tcblower
\textbf{Solution détaillée.}

Analyse des prépositions :
\begin{itemize}
\item \esp{luchan \textbf{por} sus derechos} : cause, motif $\rightarrow$ « luttent pour leurs droits ».
\item \esp{\textbf{por} un futuro mejor} : cause $\rightarrow$ « pour un avenir meilleur ».
\item \esp{\textbf{Para} ellos} : opinion, point de vue $\rightarrow$ « Pour eux ».
\item \esp{fue creado \textbf{por} la comunidad} : agent du passif $\rightarrow$ « a été créé par la communauté ».
\item \esp{\textbf{para} los barrios pobres} : destinataire $\rightarrow$ « pour les quartiers pauvres ».
\item \esp{\textbf{por eso}} : expression figée $\rightarrow$ « c'est pourquoi ».
\item \esp{trabajan \textbf{para} construir} : \esp{para + infinitif} = but $\rightarrow$ « travaillent pour construire / afin de construire ».
\end{itemize}

Vocabulaire clé : \esp{los derechos} = les droits ; \esp{la llave} = la clé ; \esp{el barrio} = le quartier ; \esp{justo/a} = juste, équitable ; \esp{la comunidad} = la communauté.

\textbf{Traduction rédigée :}

« Les jeunes luttent pour leurs droits et pour un avenir meilleur. Pour eux, l'éducation est la clé du monde. Ce programme a été créé par la communauté pour les quartiers pauvres, et c'est pourquoi tous travaillent pour construire une ville plus juste. »

\textbf{Conclusion.} Chaque préposition a été traduite selon sa valeur contextuelle ; la phrase française est fluide et fidèle.
\end{exoresolu}

\courssub{Savoir-faire 4 : Débattre des droits et de la citoyenneté}

\begin{aretenir}[Lexique essentiel : La citoyenneté et les droits]
\begin{center}
\begin{tabularx}{\linewidth}{|X|X|}
\hline
\rowcolor{popblueL} \textbf{Español} & \textbf{Français} \\
\hline
\esp{los derechos humanos} & les droits de l'homme \\
\esp{la libertad de expresión} & la liberté d'expression \\
\esp{la igualdad (de género, racial)} & l'égalité (des sexes, raciale) \\
\esp{la justicia social} & la justice sociale \\
\esp{el deber / los deberes} & le devoir / les devoirs \\
\esp{la solidaridad} & la solidarité \\
\esp{la participación ciudadana} & la participation citoyenne \\
\esp{votar / el voto / las elecciones} & voter / le vote / les élections \\
\esp{respetar / el respeto} & respecter / le respect \\
\esp{la diversidad cultural} & la diversité culturelle \\
\esp{la ciudadanía global / mundial} & la citoyenneté mondiale \\
\esp{las fronteras} & les frontières \\
\hline
\end{tabularx}
\end{center}
\end{aretenir}

\begin{savaistu}[Cameroun et monde hispanophone : citoyenneté en partage]
Le Cameroun, pays bilingue officiel (français/anglais), partage avec la \esp{Guinea Ecuatorial} (seul pays africain hispanophone) une frontière et des enjeux communs : jeunesse, migration, climat. Être \esp{ciudadano del mundo} ici, c'est défendre les droits dans le bassin du Congo comme à Madrid ou Mexico. Le \esp{bendskin} ou le \esp{makossa} portent aussi des messages citoyens (ex. \esp{Manu Dibango}, \esp{Lapiro de Mbanga}). La \esp{Francophonie} et l'\esp{Hispanophonie} ne s'opposent pas : elles s'enrichissent.
\end{savaistu}

\begin{methode}[Participer à un débat guidé en espagnol]
\begin{enumerate}
\item \textbf{Préparer sa position} avant de parler : \esp{estoy a favor / estoy en contra / estoy en parte de acuerdo}. Noter deux arguments et un exemple concret (Cameroun, Espagne, Amérique latine).
\item \textbf{Prendre la parole avec des formules d'opinion :}
\begin{itemize}
\item \esp{En mi opinión... / A mi juicio... / Para mí...} « À mon avis... »
\item \esp{Creo que / Pienso que / Me parece que + indicatif} « Je pense que... »
\item \esp{No creo que + subjonctif} : \esp{No creo que \textbf{sea} justo.} « Je ne crois pas que ce soit juste. »
\end{itemize}
\item \textbf{Réagir aux autres :}
\begin{itemize}
\item Accord : \esp{Estoy de acuerdo contigo. / Tienes razón.} « Je suis d'accord avec toi. / Tu as raison. »
\item Désaccord poli : \esp{No estoy de acuerdo porque... / Entiendo tu punto de vista, pero...} « Je ne suis pas d'accord parce que... / Je comprends ton point de vue, mais... »
\end{itemize}
\item \textbf{Argumenter avec des connecteurs :} \esp{primero, además, sin embargo, por eso, en conclusión}.
\item \textbf{Employer le lexique citoyen} (voir tableau ci-dessus).
\item \textbf{Conclure son intervention} : \esp{En conclusión, ser ciudadano del mundo significa...} « En conclusion, être citoyen du monde signifie... »
\end{enumerate}
\end{methode}

\begin{exoresolu}[Expression orale préparée : « Ciudadanos del mundo » — Type Bac]
\textbf{Énoncé.} Rédige en espagnol une intervention de 8 à 10 lignes pour le débat : \esp{¿Somos ciudadanos del mundo?} Donne ton opinion, deux arguments et une conclusion.

\tcblower
\textbf{Solution détaillée.}

\textbf{Plan de l'intervention :} (1) prise de position ; (2) argument 1 : universalité des droits ; (3) argument 2 : défis globaux ; (4) nuance (identité) ; (5) conclusion.

\textbf{Intervention rédigée (modèle) :}

\esp{En mi opinión, todos somos ciudadanos del mundo, aunque vivamos en países diferentes. Primero, los derechos humanos no tienen fronteras: la libertad y la dignidad son iguales para un niño de Yaundé y para un niño de Madrid. Además, los grandes problemas de hoy, como el cambio climático o la pobreza, no se resuelven en un solo país: por eso debemos trabajar juntos. Sin embargo, ser ciudadano del mundo no significa olvidar nuestra cultura; significa respetar la diversidad y luchar por la justicia en todas partes. En conclusión, pienso que la ciudadanía mundial es un deber y una esperanza para las nuevas generaciones.}

\textbf{Traduction :} « À mon avis, nous sommes tous citoyens du monde, même si nous vivons dans des pays différents. D'abord, les droits de l'homme n'ont pas de frontières : la liberté et la dignité sont les mêmes pour un enfant de Yaoundé et pour un enfant de Madrid. De plus, les grands problèmes d'aujourd'hui, comme le changement climatique ou la pauvreté, ne se résolvent pas dans un seul pays : c'est pourquoi nous devons travailler ensemble. Cependant, être citoyen du monde ne signifie pas oublier notre culture ; cela signifie respecter la diversité et lutter pour la justice partout. En conclusion, je pense que la citoyenneté mondiale est un devoir et un espoir pour les nouvelles générations. »

\textbf{Points grammaticaux réussis :}
\begin{itemize}
\item \esp{aunque vivamos} : subjonctif après \esp{aunque} (concession réelle).
\item \esp{para un niño} (destinataire/bénéficiaire), \esp{luchar \textbf{por} la justicia} (cause).
\item \esp{por eso} (connecteur figé = c'est pourquoi).
\item \esp{pienso que + indicatif} (opinion affirmative), \esp{no creo que + subjonctif} (évité ici car affirmatif).
\item \esp{la diversidad, la justicia, los derechos humanos} : lexique précis.
\end{itemize}

\textbf{Conclusion.} L'intervention suit le plan opinion-arguments-nuance-conclusion, utilise le lexique citoyen et les connecteurs exigés : c'est le format attendu à l'épreuve orale.
\end{exoresolu}

\courssub{Savoir-faire 5 : Questions de compréhension sur un texte poétique}

\begin{textebox}[Poème original — Compréhension écrite]
\esp{Mi voz no es solo mía,}\\
\esp{es la voz de los que callan,}\\
\esp{de los niños sin escuela,}\\
\esp{de las madres que trabajan.}
\legende{Poème original, thème : voix des sans-voix, engagement social.}
\end{textebox}

\begin{methode}[Répondre aux questions de compréhension]
\begin{enumerate}
\item \textbf{Lire la question avant le texte} pour savoir quoi chercher.
\item \textbf{Localiser la réponse dans le texte} : souligner le passage concerné et citer un mot-clé espagnol entre guillemets.
\item \textbf{Répondre en espagnol, par une phrase complète}, en reformulant : ne jamais recopier le texte tel quel.
\item \textbf{Justifier quand la question le demande} (\esp{¿Por qué?} $\rightarrow$ répondre avec \esp{porque...}).
\item \textbf{Pour les questions personnelles} (\esp{¿Y tú? ¿Qué opinas?}), donner son avis avec \esp{En mi opinión...} et un argument.
\end{enumerate}
\end{methode}

\begin{exoresolu}[Compréhension d'un poème engagé — Type Bac]
\textbf{Énoncé.} Lis le poème et réponds en espagnol aux questions.

\esp{Mi voz no es solo mía,}\\
\esp{es la voz de los que callan,}\\
\esp{de los niños sin escuela,}\\
\esp{de las madres que trabajan.}

\begin{enumerate}
\item ¿De quién es la voz del poeta?
\item ¿Qué dos grupos de personas menciona el poema?
\item ¿Por qué es un poema comprometido?
\end{enumerate}

\tcblower
\textbf{Solution détaillée.}

\begin{enumerate}
\item Question : « De qui est la voix du poète ? » Le vers 2 dit \esp{es la voz de los que callan}.
Réponse rédigée : \esp{La voz del poeta no es solo suya: es también la voz de las personas que no pueden hablar, es decir, « la voz de los que callan ».}
Justification : phrase complète, reformulation (\esp{las personas que no pueden hablar}) + citation courte.

\item Question : « Quels deux groupes de personnes le poème mentionne-t-il ? »
Réponse rédigée : \esp{El poema menciona a los niños que no van a la escuela y a las madres que trabajan.}
Justification : on reformule \esp{niños sin escuela} par \esp{los niños que no van a la escuela}, ce qui montre la compréhension.

\item Question : « Pourquoi est-ce un poème engagé ? »
Réponse rédigée : \esp{Es un poema comprometido porque el poeta defiende a los más débiles y denuncia la injusticia social: habla por los que no tienen voz.}
Justification : on répond à \esp{¿Por qué?} avec \esp{porque}, et on utilise le vocabulaire du commentaire (\esp{defender}, \esp{denunciar}, \esp{la injusticia social}).
\end{enumerate}

\textbf{Conclusion.} Chaque réponse est une phrase complète en espagnol, appuyée sur le texte mais reformulée : c'est exactement ce que le correcteur attend.
\end{exoresolu}

\begin{attention}[Les 5 erreurs qui coûtent des points au Bac]
\begin{enumerate}
\item \textbf{Confondre cause et but :} écrire \esp{Estudio por ser médico} au lieu de \esp{Estudio \textbf{para} ser médico}. \cle{Réflexe} : devant un infinitif de but, c'est toujours \esp{para}.
\item \textbf{Oublier les accents obligatoires :} écrire \esp{educacion} au lieu de \esp{educación} ; \esp{raices} au lieu de \esp{raíces} ; \esp{corazon} au lieu de \esp{corazón} ; \esp{expresion} au lieu de \esp{expresión}. \esp{Libertad} et \esp{poema} ne prennent \textbf{pas} d'accent. Chaque accent manquant ou mis à tort est une faute d'orthographe.
\item \textbf{Faux amis fréquents :} \esp{largo} = « long » (pas « large » = \esp{ancho}) ; \esp{actualmente} = « actuellement / en ce moment » (pas « en fait » = \esp{realmente}) ; \esp{una excusa} = « une excuse » (pas « un prétexte noble » = \esp{un pretexto}) ; \esp{sentir} (en poésie) = « regretter, ressentir profondément » (pas seulement « sentir »).
\item \textbf{Calquer la syntaxe française :} « Je suis d'accord avec toi » $\rightarrow$ \esp{Estoy de acuerdo \textbf{contigo}} (jamais \esp{con tú}) ; « pour moi » $\rightarrow$ \esp{\textbf{para} mí} (avec accent, sinon \esp{mi} = mon) ; « je m'appelle » $\rightarrow$ \esp{Me llamo} (pas \esp{Yo llamo}).
\item \textbf{Accord et temps :} après \esp{no creo que}, le subjonctif est obligatoire (\esp{No creo que \textbf{sea} justo}) ; dans la voix passive, le participe s'accorde en genre et nombre : \esp{la ley fue \textbf{aprobada} por el parlamento} ; dans une traduction de vers, respecter le temps du poème : \esp{fue} (indéfini) = « a été / fut » (passé composé ou passé simple littéraire), jamais un présent.
\end{enumerate}
\end{attention}

\newpage

\coursec{Exercices}

\courssub{Je vérifie mes connaissances}

\begin{exercice}[QCM : por ou para ?]
Pour chaque phrase, choisis la bonne réponse et entoure-la.
\begin{enumerate}
\item Este poema fue escrito \textbf{por / para} un poeta anónimo después de la guerra.
\item Luchamos \textbf{por / para} la paz desde hace muchos años.
\item Compré este diccionario \textbf{por / para} mi hermano que estudia español.
\item Salimos de Yaoundé \textbf{por / para} Douala mañana por la mañana.
\item \textbf{Por / Para} ser tan joven, habla muy bien el español.
\end{enumerate}
\end{exercice}

\begin{exercice}[Vrai ou faux ? Justifie ta réponse]
Dis si chaque affirmation est vraie ou fausse, puis justifie en une phrase en français.
\begin{enumerate}
\item \esp{Por} exprime toujours la finalité, le but.
\item Dans \esp{Gracias por tu ayuda}, \esp{por} exprime la cause ou le motif.
\item \esp{Para} s'emploie devant un infinitif pour exprimer le but : \esp{Estudio para aprender.}
\item Dans une phrase passive, l'agent est toujours introduit par \esp{para}.
\item \esp{Por} peut exprimer la durée : \esp{Viví en España por tres años.}
\end{enumerate}
\end{exercice}

\begin{exercice}[Vocabulaire : citoyenneté et poésie]
Complète chaque phrase espagnole avec le mot qui convient, choisi dans la liste : \esp{la ciudadanía} -- \esp{los derechos humanos} -- \esp{el verso} -- \esp{la fraternidad} -- \esp{la patria} -- \esp{el poeta}.
\begin{enumerate}
\item Un \ldots\ldots\ldots\ldots\ldots\ est une ligne d'un poème.
\item Celui qui écrit des poèmes est un \ldots\ldots\ldots\ldots\ldots\ .
\item La \ldots\ldots\ldots\ldots\ldots\ est le sentiment d'amitié et de solidarité entre tous les hommes.
\item Toute personne, où qu'elle vive, doit jouir de \ldots\ldots\ldots\ldots\ldots\ .
\item La \ldots\ldots\ldots\ldots\ldots\ est le statut juridique qui relie une personne à son pays.
\end{enumerate}
\end{exercice}

\begin{exercice}[Conjugaison à compléter]
Conjugue le verbe entre parenthèses au temps indiqué, dans des phrases sur la citoyenneté.
\begin{enumerate}
\item (ser, présent) Nosotros \ldots\ldots\ldots\ ciudadanos del mundo.
\item (luchar, présent) Los jóvenes \ldots\ldots\ldots\ por sus derechos.
\item (poder, futur) Un día, todos los hombres \ldots\ldots\ldots\ vivir en paz.
\item (viajar, pretérito perfecto) Este año yo \ldots\ldots\ldots\ a Guinea Ecuatorial.
\item (querer, imparfait) De niño, ella \ldots\ldots\ldots\ ser poeta.
\item (trabajar, présent du subjonctif) Es importante que tú \ldots\ldots\ldots\ para tu comunidad.
\end{enumerate}
\end{exercice}

\courssub{Je m'entraîne}

\begin{exercice}[Thème : por et para en action]
Traduis en espagnol les phrases suivantes. Attention au choix entre \esp{por} et \esp{para}.
\begin{enumerate}
\item Merci pour ton aide, mon ami.
\item Ce poème a été écrit par un jeune Camerounais.
\item J'apprends l'espagnol pour voyager en Amérique latine.
\item Il a vécu à Malabo pendant trois ans.
\item Ce cadeau est pour ma sœur.
\end{enumerate}
\end{exercice}

\begin{exercice}[Texte à trous : un jeune migrant]
Complète le texte suivant avec \esp{por} ou \esp{para} (10 occurrences), puis justifie chaque choix (cause, finalité, durée, agent, destinataire, lieu, échange).
\begin{textebox}[Un joven camerunés en España]
Samuel salió de Douala \ldots\ldots\ldots\ (1) Madrid hace dos años. Viajó \ldots\ldots\ldots\ (2) buscar un futuro mejor. \ldots\ldots\ldots\ (3) el camino, pasó \ldots\ldots\ldots\ (4) varios países. Ahora trabaja \ldots\ldots\ldots\ (5) una ONG que ayuda a los migrantes. \ldots\ldots\ldots\ (6) la mañana estudia español y \ldots\ldots\ldots\ (7) la tarde trabaja. Este poema fue escrito \ldots\ldots\ldots\ (8) él \ldots\ldots\ldots\ (9) sus amigos de Camerún. « Lucho \ldots\ldots\ldots\ (10) mis derechos y \ldots\ldots\ldots\ los derechos de todos », dice Samuel.
\end{textebox}
\end{exercice}

\begin{exercice}[Corrige les erreurs d'un francophone]
Un élève francophone a écrit les phrases suivantes. Chacune contient une erreur avec \esp{por} ou \esp{para}. Corrige la phrase et explique l'erreur en français.
\begin{enumerate}
\item \esp{Gracias para tu carta.}
\item \esp{Viví en España durante por dos años.}
\item \esp{Este libro fue escrito para Cervantes.}
\item \esp{Estudio español por aprender una lengua nueva.}
\item \esp{Salimos por la mañana para ir por Madrid.}
\end{enumerate}
\end{exercice}

\begin{exercice}[Appariements : fonction et exemple]
Relie chaque emploi de \esp{por} ou \esp{para} (colonne de gauche) à la phrase qui l'illustre (colonne de droite).
\begin{center}
\begin{tabularx}{\linewidth}{|X|X|}
\hline
\rowcolor{popblueL} Emploi & Exemple \\
\hline
1. PARA + destinataire & a. \esp{Luchamos por la libertad.} \\
\hline
2. POR + cause, motif & b. \esp{Estas flores son para mi madre.} \\
\hline
3. PARA + but (+ infinitif) & c. \esp{El tren pasa por Douala.} \\
\hline
4. POR + durée & d. \esp{Vino a clase para aprender.} \\
\hline
5. POR + lieu (passage) & e. \esp{Estuvo enfermo por una semana.} \\
\hline
\end{tabularx}
\end{center}
\end{exercice}

\courssub{Je me prépare au Bac}

\begin{exercice}[Version type Bac : un poème pour la paix]
Lis le poème suivant, puis traduis les vers soulignés (vers 1 à 5) en français. Glossaire partiel : \esp{sembrar} = semer ; \esp{la semilla} = la graine ; \esp{el surco} = le sillon ; \esp{el odio} = la haine.
\begin{textebox}[Sembrar la paz]
\underline{Sembremos la paz en cada surco,}\\
\underline{que la tierra no quiere más odio:}\\
\underline{una semilla de amor por el mundo,}\\
\underline{una palabra para el hermano,}\\
\underline{un pan partido para el que sufre.}\\
Los niños de Douala y de Madrid\\
sueñan con el mismo cielo azul,\\
y el viento no conoce fronteras\\
cuando canta la voz de la fraternidad.
\end{textebox}
Questions de langue :
\begin{enumerate}
\item Releve dans le poème deux emplois de \esp{por} et deux emplois de \esp{para} et justifie chacun (cause, but, destinataire, lieu).
\item Conjugue \esp{sembrar} au présent de l'indicatif à toutes les personnes.
\end{enumerate}
\end{exercice}

\begin{exercice}[Thème type Bac : cinq phrases]
Traduis en espagnol les cinq phrases suivantes. Chaque phrase mobilise un emploi différent de \esp{por} ou \esp{para} : cause, finalité, agent du passif, durée, destinataire.
\begin{enumerate}
\item Nous luttons pour la paix et pour la justice.
\item Ce poème célèbre a été écrit par un poète espagnol.
\item Elle étudie le droit pour défendre les citoyens.
\item Ils ont marché dans la rue pendant deux heures.
\item J'ai acheté ce livre de poésie pour mon professeur.
\end{enumerate}
\end{exercice}

\begin{exercice}[Expression écrite type Bac : « Ciudadanos del mundo »]
Sujet : \esp{¿Es posible ser ciudadano del mundo?}

Rédige en espagnol un paragraphe argumenté d'environ 10 lignes (entre 100 et 120 mots) dans lequel tu donnes ton opinion sur la citoyenneté mondiale.

Consignes obligatoires :
\begin{enumerate}
\item Utilise au moins trois fois \esp{por} ou \esp{para} avec des emplois différents (cause, but, destinataire, durée, opinion).
\item Emploie au moins deux connecteurs : \esp{en primer lugar}, \esp{además}, \esp{sin embargo}, \esp{por eso}, \esp{en conclusión}.
\item Utilise au moins une fois le présent du subjonctif après \esp{es importante que} ou \esp{no creo que}.
\end{enumerate}

\end{exercice}

\begin{experience}[Débat flash en binômes]
Par binômes, préparez un débat de 4 échanges minimum (2 interventions chacun) : l'élève A défend la citoyenneté nationale (\esp{Mi patria es lo primero}), l'élève B défend la citoyenneté mondiale (\esp{Todos somos ciudadanos del mundo}). Utilisez les expressions : \esp{Estoy de acuerdo contigo porque...}, \esp{No estoy de acuerdo, por eso...}, \esp{En mi opinión...}, \esp{Sin embargo...}. Chaque élève doit employer au moins deux fois \esp{por} ou \esp{para}.
\end{experience}


\newpage

\coursec{Corrigés des exercices}

\coursec{Leçon E12 : Traduction de vers, Por y Para, Debate « Ciudadanos del mundo »}

\courssub{Corrigés des exercices d'application et de synthèse}

\begin{corrige}[Exercice 1 — QCM : por ou para ?]
\begin{enumerate}
\item Este poema fue escrito \textbf{por} un poeta anónimo después de la guerra. \emph{(agent du passif : « par »)}
\item Luchamos \textbf{por} la paz desde hace muchos años. \emph{(cause, motif : « pour »)}
\item Compré este diccionario \textbf{para} mi hermano que estudia español. \emph{(destinataire)}
\item Salimos de Yaoundé \textbf{para} Douala mañana temprano. \emph{(destination, direction ; \esp{mañana por la mañana} est redondant)}
\item \textbf{Para} ser tan joven, habla muy bien el español. \emph{(locution \esp{para ser} = « pour être », concession)}
\end{enumerate}
\textbf{Astuce de vérification :} si l'on peut remplacer par « à destination de » ou « destiné à », c'est \esp{para} ; si l'on peut remplacer par « à cause de », « par » (agent) ou « pendant », c'est \esp{por}.
\end{corrige}

\begin{corrige}[Exercice 2 — Vrai ou faux ?]
\begin{enumerate}
\item \textbf{Faux.} C'est \esp{para} qui exprime la finalité (\esp{Estudio para aprender}) ; \esp{por} exprime surtout la cause, le motif, la durée, l'agent ou le passage.
\item \textbf{Vrai.} Dans \esp{Gracias por tu ayuda} (« merci pour ton aide »), \esp{por} indique le motif du remerciement.
\item \textbf{Vrai.} \esp{Para} + infinitif exprime le but : \esp{Estudio para aprender} (« j'étudie pour apprendre »).
\item \textbf{Faux.} Dans une phrase passive, le complément d'agent est introduit par \esp{por} : \esp{El poema fue escrito por el poeta}.
\item \textbf{Vrai.} \esp{Por} + durée : \esp{Viví en España por tres años} (« j'ai vécu en Espagne pendant trois ans »). On peut aussi omettre la préposition ou utiliser \esp{durante}.
\end{enumerate}
\end{corrige}

\begin{corrige}[Exercice 3 — Vocabulaire : citoyenneté et poésie]
\begin{enumerate}
\item Un \textbf{verso} es una línea de un poema.
\item Celui qui écrit des poèmes est un \textbf{poeta}.
\item La \textbf{fraternidad} es el sentimiento de amistad y de solidaridad entre todos los hombres.
\item Toda persona, dondequiera que viva, debe disfrutar de \textbf{los derechos humanos}.
\item La \textbf{ciudadanía} es el estatuto jurídico que une una persona a su país.
\end{enumerate}
\textbf{Attention :} \esp{la patria} (« la patrie ») n'était pas utilisée ici ; c'était un intrus de la liste.
\end{corrige}

\begin{corrige}[Exercice 4 — Conjugaison à compléter]
\begin{enumerate}
\item Nosotros \textbf{somos} ciudadanos del mundo. \emph{(ser, présent, 1\iere{} pers. du pluriel — irrégulier)}
\item Los jóvenes \textbf{luchan} por sus derechos. \emph{(luchar, présent, 3\ieme{} pers. du pluriel — régulier en -ar)}
\item Un día, todos los hombres \textbf{podrán} vivir en paz. \emph{(poder, futur : radical irrégulier podr-)}
\item Este año yo \textbf{he viajado} a Guinea Ecuatorial. \emph{(pretérito perfecto : haber au présent + participe)}
\item De niña, ella \textbf{quería} ser poeta. \emph{(querer, imparfait : terminaison régulière -ía)}
\item Es importante que tú \textbf{trabajes} para tu comunidad. \emph{(subjonctif présent obligatoire après \esp{es importante que})}
\end{enumerate}
\textbf{Point de vigilance :} ne pas écrire \esp{*poderán} au futur ; le radical de \esp{poder} au futur est \esp{podr-}.
\end{corrige}

\begin{corrige}[Exercice 5 — Thème : por et para en action]
\begin{enumerate}
\item \esp{Gracias por tu ayuda, amigo mío.} \emph{(por + cause : motif du remerciement)}
\item \esp{Este poema fue escrito por un joven camerunés.} \emph{(por + agent du passif)}
\item \esp{Aprendo español para viajar por América Latina.} \emph{(para + infinitif = but ; por + lieu = déplacement à travers un espace)}
\item \esp{Vivió en Malabo por tres años.} \emph{(por + durée ; on accepte aussi \esp{durante tres años} ou sans préposition)}
\item \esp{Este regalo es para mi hermana.} \emph{(para + destinataire)}
\end{enumerate}
\textbf{Erreurs à éviter :} \esp{*gracias para tu ayuda} (calque du français « merci pour ») ; \esp{*escrito para} (agent du passif).
\end{corrige}

\begin{corrige}[Exercice 6 — Texte à trous : un jeune migrant]
\begin{enumerate}
\item \textbf{para} Madrid — destination, direction (« à destination de »).
\item \textbf{para} buscar — finalité, but (+ infinitif).
\item \textbf{Por} el camino — lieu, trajet (« en chemin »).
\item \textbf{por} varios países — lieu, passage à travers.
\item \textbf{para} una ONG — employeur / au service de (destinataire de l'activité).
\item \textbf{Por} la mañana — moment de la journée (temps).
\item \textbf{por} la tarde — moment de la journée (temps).
\item \textbf{por} él — agent du passif (\esp{fue escrito por él}).
\item \textbf{para} sus amigos — destinataire.
\item \textbf{por} mis derechos y \textbf{por} los derechos de todos — cause, motif (on lutte « pour » une cause).
\end{enumerate}
\textbf{Traduction du texte :} « Samuel a quitté Douala pour Madrid il y a deux ans. Il a voyagé pour chercher un avenir meilleur. En chemin, il est passé par plusieurs pays. Maintenant il travaille pour une ONG qui aide les migrants. Le matin il étudie l'espagnol et l'après-midi il travaille. Ce poème a été écrit par lui pour ses amis du Cameroun. "Je lutte pour mes droits et pour les droits de tous", dit Samuel. »
\end{corrige}

\begin{corrige}[Exercice 7 — Corrige les erreurs d'un francophone]
\begin{enumerate}
\item \esp{Gracias para tu carta.} $\rightarrow$ \esp{Gracias \textbf{por} tu carta.} Le remerciement exprime un motif : toujours \esp{por} (calque du français « merci pour »).
\item \esp{Viví en España durante por dos años.} $\rightarrow$ \esp{Viví en España \textbf{por} dos años} (ou \esp{\textbf{durante} dos años}). On ne cumule jamais \esp{durante} et \esp{por}.
\item \esp{Este libro fue escrito para Cervantes.} $\rightarrow$ \esp{Este libro fue escrito \textbf{por} Cervantes.} Le complément d'agent du passif s'introduit par \esp{por} ; \esp{para} marquerait le destinataire.
\item \esp{Estudio español por aprender una lengua nueva.} $\rightarrow$ \esp{Estudio español \textbf{para} aprender una lengua nueva.} Devant un infinitif de but, c'est \esp{para} ; \esp{por} + infinitif exprime la cause (\esp{por no estudiar} = « parce qu'il n'a pas étudié »).
\item \esp{Salimos por la mañana para ir por Madrid.} $\rightarrow$ \esp{Salimos por la mañana para ir \textbf{a} Madrid.} La destination précise exige \esp{a} (ou \esp{para} Madrid) ; \esp{por Madrid} signifierait « à travers Madrid ».
\end{enumerate}
\end{corrige}

\begin{corrige}[Exercice 8 — Appariements : fonction et exemple]
\begin{center}
\begin{tabularx}{\linewidth}{|X|X|}
\hline
\rowcolor{popblueL} Emploi & Exemple \\
\hline
1. PARA + destinataire & \textbf{b.} \esp{Estas flores son para mi madre.} \\
\hline
2. POR + cause, motif & \textbf{a.} \esp{Luchamos por la libertad.} \\
\hline
3. PARA + but (+ infinitif) & \textbf{d.} \esp{Vino a clase para aprender.} \\
\hline
4. POR + durée & \textbf{e.} \esp{Estuvo enfermo por una semana.} \\
\hline
5. POR + lieu (passage) & \textbf{c.} \esp{El tren pasa por Douala.} \\
\hline
\end{tabularx}
\end{center}
\textbf{Vérification :} 1-b (les fleurs sont destinées à la mère) ; 2-a (la liberté est la cause du combat) ; 3-d (apprendre est le but) ; 4-e (une semaine = durée) ; 5-c (le train traverse Douala).
\end{corrige}

\begin{corrige}[Exercice 9 — Version type Bac : un poème pour la paix]
\begin{textebox}[Poème original (vers 1--5)]
Sembremos la paz en cada surco, \\
que la tierra no quiere más odio: \\
una semilla de amor por el mundo, \\
una palabra para el hermano, \\
un pan partido para el que sufre.
\end{textebox}

\textbf{Traduction des vers 1 à 5 (proposition) :}

« Semons la paix dans chaque sillon, / car la terre ne veut plus de haine : / une graine d'amour pour le monde, / une parole pour le frère, / un pain partagé pour celui qui souffre. »

\textbf{Remarques de traduction :}
\begin{itemize}
\item \esp{Sembremos} est un subjonctif présent à valeur d'impératif collectif (« semons », invitation).
\item \esp{que} en tête du vers 2 a valeur explicative : « car ».
\item \esp{partido} = « partagé » (littéralement « divisé ») ; on préfère « partagé » en français.
\item Respecter le rythme : cinq vers courts, sans ajouter de mots.
\end{itemize}

\textbf{Questions de langue :}
\begin{enumerate}
\item Emplois de \esp{por} : \esp{una semilla de amor \textbf{por} el mundo} (lieu, diffusion à travers le monde). Emplois de \esp{para} : \esp{una palabra \textbf{para} el hermano} (destinataire) ; \esp{un pan partido \textbf{para} el que sufre} (destinataire).
\item Présent de l'indicatif de \esp{sembrar} (verbe à changement de radical e $\rightarrow$ ie) : yo \textbf{siembro}, tú \textbf{siembras}, él/ella/usted \textbf{siembra}, nosotros \textbf{sembramos}, vosotros \textbf{sembráis}, ellos/ustedes \textbf{siembran}.
\end{enumerate}

\textbf{Barème indicatif (sur 20) :} traduction des 5 vers : 10 pts (2 pts par vers) ; questions de langue : 6 pts (4 pts pour por/para justifiés, 2 pts pour la conjugaison) ; qualité de la langue française : 4 pts.

\textbf{Points de vigilance :} ne pas traduire \esp{surco} par « sillon » oublié ou par « champ » ; ne pas confondre \esp{el odio} (la haine) avec « l'odeur » ; accorder « partagé » avec « pain ».
\end{corrige}

\begin{corrige}[Exercice 10 — Thème type Bac : cinq phrases]
\begin{enumerate}
\item \esp{Luchamos \textbf{por} la paz y \textbf{por} la justicia.} \emph{(cause, motif — on répète la préposition devant chaque nom)}
\item \esp{Este famoso poema fue escrito \textbf{por} un poeta español.} \emph{(agent du passif)}
\item \esp{Ella estudia derecho \textbf{para} defender a los ciudadanos.} \emph{(but + infinitif ; a personnelle devant \esp{los ciudadanos})}
\item \esp{Caminaron por la calle \textbf{por} dos horas.} (ou \esp{\textbf{durante} dos horas}) \emph{(durée ; \esp{por la calle} = lieu, passage)}
\item \esp{He comprado este libro de poesía \textbf{para} mi profesor.} \emph{(destinataire)}
\end{enumerate}
\textbf{Barème indicatif (sur 10) :} 2 pts par phrase (1 pt pour le choix por/para, 1 pt pour l'exactitude globale : accents, a personnelle, temps).

\textbf{Points de vigilance :} \esp{derecho} = « le droit » (discipline juridique), sans article après \esp{estudiar} ; ne pas oublier l'a personnelle dans la phrase 3 ; \esp{poesía} prend un accent (hiato).
\end{corrige}

\begin{corrige}[Exercice 11 — Expression écrite type Bac : « Ciudadanos del mundo »]
\textbf{Proposition de rédaction modèle (environ 110 mots) :}

\begin{textebox}[¿Es posible ser ciudadano del mundo?]
En mi opinión, sí es posible ser ciudadano del mundo. En primer lugar, todos los seres humanos luchan \textbf{por} los mismos derechos: la libertad, la educación y la paz. Además, gracias a Internet, un joven de Yaoundé puede hablar con un joven de Madrid \textbf{para} compartir sus sueños. Sin embargo, no debemos olvidar nuestras raíces: amo Camerún y trabajo \textbf{para} mi comunidad. Es importante que cada joven respete las culturas de los demás. En conclusión, ser ciudadano del mundo no significa olvidar la patria, sino abrir el corazón \textbf{a} todos los hombres.
\end{textebox}

\textbf{Traduction :} « À mon avis, oui, il est possible d'être citoyen du monde. En premier lieu, tous les êtres humains luttent pour les mêmes droits : la liberté, l'éducation et la paix. De plus, grâce à Internet, un jeune de Yaoundé peut parler avec un jeune de Madrid pour partager ses rêves. Cependant, nous ne devons pas oublier nos racines : j'aime le Cameroun et je travaille pour ma communauté. Il est important que chaque jeune respecte les cultures des autres. En conclusion, être citoyen du monde ne signifie pas oublier la patrie, mais ouvrir son cœur à tous les hommes. »

\textbf{Vérification des consignes :} trois emplois différents de por/para : \esp{por los mismos derechos} (cause), \esp{para compartir} (but), \esp{para mi comunidad} (destinataire) ; connecteurs : \esp{en primer lugar}, \esp{además}, \esp{sin embargo}, \esp{en conclusión} ; subjonctif : \esp{es importante que cada joven \textbf{respete}}.

\textbf{Barème indicatif (sur 20) :} respect des consignes et de la longueur : 4 pts ; correction grammaticale (por/para, subjonctif, accords) : 8 pts ; richesse du lexique et des connecteurs : 4 pts ; cohérence de l'argumentation : 4 pts.

\textbf{Points de vigilance :} accorder \esp{ciudadano del mundo} ; ne pas écrire \esp{*es posible de ser} ; subjonctif obligatoire après \esp{es importante que} ; \esp{sin embargo} s'écrit en deux mots ; \esp{amar} pays/concept = COD direct sans \esp{a} (standard), \esp{a} personnel seulement pour personnes/animaux ou personification poétique.
\end{corrige}

\begin{corrige}[Exercice 12 — Débat flash en binômes (grille d'évaluation)]
\textbf{Exemple de débat réussi :}

\begin{textebox}[Diálogo modelo]
\textbf{A :} En mi opinión, mi patria es lo primero. Trabajo \textbf{para} mi país porque aquí están mi familia y mis raíces.\\
\textbf{B :} No estoy de acuerdo, por eso te digo que todos luchamos \textbf{por} los mismos derechos humanos, en Camerún o en España.\\
\textbf{A :} Sin embargo, sin una patria fuerte, ¿quién protege a los ciudadanos? Pago mis impuestos \textbf{para} construir escuelas aquí.\\
\textbf{B :} Estoy de acuerdo contigo en eso, pero el mundo es una gran familia: debemos abrir nuestro corazón \textbf{a} todos los pueblos.
\end{textebox}

\textbf{Grille d'évaluation orale (sur 10) :}
\begin{itemize}
\item 4 échanges minimum réalisés : 2 pts.
\item Au moins deux emplois corrects de \esp{por}/\esp{para} par élève : 3 pts.
\item Utilisation des expressions du débat (\esp{estoy de acuerdo}, \esp{sin embargo}, \esp{en mi opinión}) : 2 pts.
\item Prononciation, fluidité et écoute de l'autre : 3 pts.
\end{itemize}

\textbf{Conseils au professeur :} rappeler que \esp{estar de acuerdo} se construit avec \esp{con} (\esp{estoy de acuerdo contigo}) et sanctionner le calque \esp{*estoy de acuerdo a ti} ; valoriser les arguments culturels (Guinée équatoriale, francophonie, migrations) plutôt que la seule forme.
\end{corrige}

\newpage

\coursec{Fiche bilan}

\begin{popfigure}
\begin{tikzpicture}[mindmap, grow cyclic, every node/.style={concept, font=\small},
  root concept/.append style={concept color=popblue, font=\bfseries},
  level 1/.append style={level distance=3.6cm, sibling angle=60, font=\footnotesize},
  level 2/.append style={level distance=2.2cm, sibling angle=40, font=\scriptsize}]
\node[root concept]{E12 : Por / Para, poème, débat}
  child[concept color=poporange]{ node {Poème engagé}
    child { node {vers, estrofa, rima} }
    child { node {tono, tema, yo poético} }
  }
  child[concept color=popgreen]{ node {Traduire des vers}
    child { node {sens avant littéral} }
    child { node {rythme et images} }
  }
  child[concept color=poppurple]{ node {POR}
    child { node {cause, agent, échange} }
    child { node {durée, lieu (à travers)} }
  }
  child[concept color=poppink]{ node {PARA}
    child { node {finalité, destinataire} }
    child { node {direction, date limite} }
  }
  child[concept color=popgold]{ node {Débat}
    child { node {derechos, deberes} }
    child { node {opinar : creo que...} }
  }
  child[concept color=popblue!60]{ node {Citoyenneté}
    child { node {ciudadano del mundo} }
    child { node {solidaridad, tolerancia} }
  };
\end{tikzpicture}
\legende{Carte mentale de la leçon E12 : traduction de vers, POR/PARA et débat « Ciudadanos del mundo ».}
\end{popfigure}

\begin{bilan}
\textbf{Lexique clé (poésie et citoyenneté)}
\begin{itemize}
  \item \esp{el verso} : le vers ; \esp{la estrofa} : la strophe ; \esp{la rima} : la rime ; \esp{el poeta / la poetisa} : le poète / la poétesse.
  \item \esp{el yo poético} : le je poétique ; \esp{el tono} : le ton ; \esp{el tema} : le thème ; \esp{el mensaje} : le message.
  \item \esp{el/la ciudadano(a)} : le/la citoyen(ne) ; \esp{la ciudadanía} : la citoyenneté ; \esp{los derechos} : les droits ; \esp{los deberes} : les devoirs.
  \item \esp{la solidaridad} : la solidarité ; \esp{la tolerancia} : la tolérance ; \esp{la libertad} : la liberté ; \esp{la justicia} : la justice ; \esp{la patria} : la patrie.
  \item \esp{comprometerse} : s'engager ; \esp{luchar por} : lutter pour ; \esp{respetar} : respecter ; \esp{participar en} : participer à.
\end{itemize}

\textbf{Règles de grammaire à connaître par cœur : POR ou PARA}
\begin{center}
\begin{tabularx}{\linewidth}{|l|X|X|}\hline
\rowcolor{popblueL} Emploi & POR & PARA \\ \hline
Cause / motif & \esp{Lo hago \textbf{por} ti.} (à cause de, pour) & --- \\ \hline
Finalité / but & --- & \esp{Estudio \textbf{para} aprender.} (pour + infinitif) \\ \hline
Temps & durée : \esp{\textbf{por} dos horas} ; moment : \esp{\textbf{por} la mañana} & date limite : \esp{\textbf{para} mañana} \\ \hline
Lieu & à travers, le long de : \esp{pasear \textbf{por} la ciudad} & direction, destination : \esp{salir \textbf{para} Duala} \\ \hline
Agent (passif) & \esp{escrito \textbf{por} el poeta} & --- \\ \hline
Échange / prix & \esp{gracias \textbf{por} todo} ; \esp{pagar 500 francos \textbf{por} el libro} & --- \\ \hline
Destinataire & --- & \esp{un regalo \textbf{para} mi madre} \\ \hline
Opinion & \esp{\textbf{por} mí, no hay problema} & \esp{\textbf{para} mí, es importante} \\ \hline
\end{tabularx}
\end{center}

\textbf{Truc express} : POR regarde en arrière (la cause, l'origine, le moyen) ; PARA regarde en avant (le but, la destination, le destinataire).

\textbf{Expressions figées à mémoriser}
\begin{itemize}
  \item Avec POR : \esp{por favor}, \esp{por ejemplo}, \esp{por eso}, \esp{por fin}, \esp{por lo tanto}, \esp{por supuesto}, \esp{por aquí}, \esp{por teléfono}.
  \item Avec PARA : \esp{para siempre}, \esp{para qué}, \esp{para nada}, \esp{ser para + inf.} (\esp{Es para llorar.}).
\end{itemize}

\textbf{Méthode express : traduire des vers}
\begin{enumerate}
  \item Lire le poème entier et identifier le thème et le ton.
  \item Traduire le \textbf{sens} avant les mots : repérer les images et les trouver en français.
  \item Respecter les temps de verbe et l'ordre logique ; ne jamais calquer mot à mot.
  \item Relire à voix haute : le français doit être naturel et garder l'émotion du texte.
\end{enumerate}

\textbf{Méthode express : participer au débat}
\begin{itemize}
  \item Donner son avis : \esp{creo que...}, \esp{en mi opinión...}, \esp{me parece que...}, \esp{pienso que...}.
  \item Être d'accord : \esp{estoy de acuerdo (con...)}, \esp{es verdad que...}.
  \item Ne pas être d'accord : \esp{no estoy de acuerdo}, \esp{sin embargo...}, \esp{al contrario...}.
  \item Argumenter : \esp{porque...}, \esp{ya que...}, \esp{por eso...}, \esp{por ejemplo...}.
\end{itemize}
\end{bilan}

\begin{aretenir}[Checklist avant le Bac]
\begin{itemize}
  \item $\square$ Je connais le vocabulaire de la poésie : \esp{verso, estrofa, rima, tono, tema, yo poético}.
  \item $\square$ Je sais traduire des vers en privilégiant le sens et l'émotion, jamais le mot à mot.
  \item $\square$ Je sais que POR exprime la cause, l'agent, l'échange, la durée et le passage dans un lieu.
  \item $\square$ Je sais que PARA exprime la finalité (\esp{para + infinitif}), le destinataire, la direction et la date limite.
  \item $\square$ Je connais par cœur les expressions figées : \esp{por favor, por ejemplo, por fin, para siempre, para nada}.
  \item $\square$ Je ne confonds pas \esp{por la mañana} (le matin, durée habituelle) et \esp{para mañana} (pour demain, échéance).
  \item $\square$ Je sais employer l'agent du passif : \esp{un poema escrito \textbf{por} Bécquer}.
  \item $\square$ Je sais donner mon avis et argumenter en espagnol : \esp{creo que, en mi opinión, porque, por eso}.
  \item $\square$ Je connais le vocabulaire de la citoyenneté : \esp{derechos, deberes, solidaridad, tolerancia, justicia}.
  \item $\square$ Je peux définir ce qu'est un \esp{ciudadano del mundo} et défendre mon point de vue en quelques phrases.
  \item $\square$ Je vérifie mes accents et la ponctuation espagnole (¿ ? ¡ !) dans mes productions écrites.
\end{itemize}
\end{aretenir}

\findecours

\end{document}
